% Dieu et la religion — Philosophie Tle Tech — Cours signé OnBuch+
% Programme officiel MINESEC. Compiler avec : tectonic N12-dieu-religion.tex
\newcommand{\DOCMATIERE}{Philosophie}
\newcommand{\DOCNIVEAU}{Tle Tech}
\newcommand{\DOCMODULE}{Philosophie – classe de Terminale technique}
\newcommand{\DOCLECON}{N12}
\newcommand{\DOCTITRE}{Dieu et la religion}
% OnBuch+ — préambule « cours signés » ENRICHI (compilé avec tectonic / XeLaTeX)
% Système visuel « Pop » d'OnBuch+ : fond crème (jamais blanc), bordures encre
% épaisses, ombres « dures » décalées, titres Archivo Black en capitales.
% Version MATHÉMATIQUES : graphes (pgfplots/tikz), boîtes pédagogiques
% (définition, propriété, méthode, attention, le savais-tu, exercice résolu,
% exercice, TP, bilan), page de garde et signature OnBuch+ ancrée en bas.
%
% Macros attendues dans main.tex AVANT \input{preamble} :
%   \DOCMATIERE  \DOCNIVEAU  \DOCMODULE  \DOCLECON (numéro)  \DOCTITRE
\documentclass[11pt]{article}

\usepackage{fontspec}
\usepackage[french]{babel}
\usepackage[a4paper,margin=2cm,top=3.4cm,bottom=2.8cm,headheight=1.4cm,footskip=1.3cm]{geometry}
\usepackage{amsmath,amssymb,amsfonts}
\usepackage{mathtools} % \xmapsto, \coloneqq... souvent utilisés par les modèles
\usepackage{cancel} % simplifications barrées (bilans de forces)
% \abs{z} (module, valeur absolue) et \norm{u} (norme), utilisés par les modèles
\providecommand{\abs}{}\let\abs\relax\DeclarePairedDelimiter{\abs}{\lvert}{\rvert}
\providecommand{\norm}{}\let\norm\relax\DeclarePairedDelimiter{\norm}{\lVert}{\rVert}
\usepackage[table]{xcolor} % [table] : fournit aussi \rowcolors (alternance de lignes)
\usepackage{pagecolor}
\usepackage{tikz}
\usetikzlibrary{arrows.meta,positioning,calc,shapes.geometric,shapes.misc,shapes.arrows,decorations.pathmorphing,decorations.pathreplacing,decorations.markings,patterns,shadows,mindmap,trees,fit,backgrounds,matrix,intersections,angles,quotes}
\usepackage{pgfplots}
\usepackage[version=4]{mhchem} % équations nucléaires \ce{^{235}_{92}U}
\usepackage[european,siunitx]{circuitikz} % schémas électriques
\pgfplotsset{compat=1.18}
\usepgfplotslibrary{fillbetween,groupplots} % aires entre courbes (calcul intégral)
\usepackage{enumitem}
\usepackage{fancyhdr}
% [most] charge toutes les bibliothèques tcolorbox (listings, minted, raster,
% documentation...) et gonfle inutilement la mémoire TeX sur un document long
% (~300 boîtes) : on ne charge que ce qui est réellement utilisé.
\usepackage[skins,breakable]{tcolorbox}
\usepackage{graphicx}
\usepackage{colortbl}
\usepackage{booktabs}
\usepackage{tabularx}
\usepackage{multirow}
\usepackage{diagbox} % case d'en-tête diagonale des tableaux à double entrée
% Colonnes extensibles centrée / gauche / droite (C, L, R), souvent utilisées par les modèles
\newcolumntype{C}{>{\centering\arraybackslash}X}
\newcolumntype{L}{>{\raggedright\arraybackslash}X}
\newcolumntype{R}{>{\raggedleft\arraybackslash}X}
\usepackage{array}
\usepackage{multicol}
\usepackage{siunitx}
% parse-numbers=false : accepte aussi \SI{2,5 \times 10^{-3}}{...}
\sisetup{locale=FR,output-decimal-marker={,},group-minimum-digits=4,parse-numbers=false}
\DeclareSIUnit\ohm{\ensuremath{\symup{\Omega}}} % U+2126 absent de la police
\DeclareSIUnit\division{div} % réglages d'oscilloscope (ms/div, V/div)
\DeclareSIUnit\tr{tr} % tours (vitesse de rotation, ex. \SI{1500}{\tr\per\minute})
\DeclareSIUnit\molar{M} % concentration molaire (ex. \SI{3}{\molar}, \SI{1}{\micro\molar})
\DeclareSIUnit\an{an} % année (le modèle confond parfois avec \year, un compteur TeX) — ex. \SI{5}{\kilo\meter\per\an}
\DeclareSIUnit\FCFA{FCFA} % franc CFA (ex. \SI{500}{\FCFA\per\kilo\gram})
\DeclareSIUnit\degreeBrix{\textdegree Brix} % teneur en sucre d'un sirop/jus (réfractométrie)
\DeclareSIUnit\month{mois} % le modèle utilise parfois \month, un compteur TeX, comme unité de durée
\DeclareSIUnit\microliter{\micro\liter} % le modèle écrit parfois \microliter en un seul token
\DeclareSIUnit\million{millions} % dénombrement (ex. \SI{15}{\million\per\mL} spermatozoïdes)
\usepackage{qrcode}
% \degree : commande fréquemment utilisée par les modèles pour un angle ou une
% température en dehors de \SI{}{} (non fournie par les paquets chargés).
\providecommand{\degree}{^{\circ}}
% \qed (fin de démonstration), utilisé par les modèles sans amsthm
\providecommand{\qed}{\hfill$\square$}
% \barre : double barre des valeurs interdites dans un tableau de variations (tkz-tab)
\providecommand{\barre}{\ensuremath{\|}}
% environnement proof (amsthm non chargé) utilisé par les modèles
\ifdefined\proof\else\newenvironment{proof}[1][Démonstration]{\par\noindent\textit{#1.}\ }{\qed\par}\fi
\usepackage{lastpage}
\usepackage{hyperref}
\hypersetup{hidelinks}

% ============================================================
% Palette « Pop » OnBuch+ — mêmes hex que la classe P (plus_ui.dart)
% ============================================================
\definecolor{poporange}{HTML}{F59321}
\definecolor{popdark}{HTML}{E07A0C}
\definecolor{popcream}{HTML}{FFF6E8}
\definecolor{popink}{HTML}{1C1714}
\definecolor{poppurple}{HTML}{7A5AE0}
\definecolor{popgreen}{HTML}{2E9E6A}
\definecolor{poppink}{HTML}{E0457B}
\definecolor{popblue}{HTML}{2F7DE1}
\definecolor{popgold}{HTML}{C98A12}
\definecolor{popmuted}{HTML}{8A7F76}
\definecolor{popline}{HTML}{EADFD2}
\definecolor{popgray}{HTML}{8A7F76}
\colorlet{popgrey}{popgray}
% Alias tolérants (le modèle nomme parfois les couleurs différemment)
\colorlet{popwhite}{white}
\colorlet{popblack}{popink}
\colorlet{popred}{poppink}
\colorlet{popyellow}{popgold}
% Teintes claires (fonds de tableaux/encadrés) — le modèle nomme parfois
% les couleurs avec un suffixe L (ex. popblueL) sans les avoir définies.
\colorlet{poporangeL}{poporange!20}
\colorlet{popdarkL}{popdark!20}
\colorlet{poppurpleL}{poppurple!20}
\colorlet{popgreenL}{popgreen!20}
\colorlet{poppinkL}{poppink!20}
\colorlet{popblueL}{popblue!20}
\colorlet{popgoldL}{popgold!20}
\colorlet{popredL}{popred!20}
\colorlet{popgrayL}{popgray!20}
% Teintes claires pour fonds de tableaux / zones
\colorlet{popblueL}{popblue!10!white}
\colorlet{poporangeL}{poporange!14!white}
\colorlet{popgreenL}{popgreen!12!white}
\colorlet{poppurpleL}{poppurple!12!white}
\colorlet{poppinkL}{poppink!12!white}
\colorlet{popgoldL}{popgold!14!white}

\pagecolor{popcream}

% ============================================================
% Typographie
% ============================================================
\defaultfontfeatures{Ligatures=TeX}
\setmainfont{PlusJakartaSans-Medium.ttf}[Path=../../../fonts/,BoldFont=PlusJakartaSans-Bold.ttf]
% Maths en sans-serif (Fira Math) avec chiffres et lettres droites de Plus
% Jakarta, pour que les formules se fondent dans le texte.
\usepackage[math-style=ISO,bold-style=ISO]{unicode-math}
\setmathfont{FiraMath-Regular.otf}[Path=../../../fonts/]
\setmathfont{PlusJakartaSans-Medium.ttf}[Path=../../../fonts/,range={up/num,up/{Latin,latin},\mathup}]
\usepackage{icomma} % virgule décimale sans espace en mode math
% Caractères mathématiques Unicode absents des polices de texte (Plus Jakarta,
% Archivo Black) : rendus par leur équivalent mathématique, en texte comme en
% formule — sinon ils s'affichent en carré vide (ex. « Arithmétique dans ℤ »).
\usepackage{newunicodechar}
\newunicodechar{ℝ}{\ensuremath{\mathbb{R}}}
\newunicodechar{ℤ}{\ensuremath{\mathbb{Z}}}
\newunicodechar{ℂ}{\ensuremath{\mathbb{C}}}
\newunicodechar{ℕ}{\ensuremath{\mathbb{N}}}
\newunicodechar{ℚ}{\ensuremath{\mathbb{Q}}}
\newunicodechar{𝔻}{\ensuremath{\mathbb{D}}}
\newunicodechar{α}{\ensuremath{\alpha}}
\newunicodechar{β}{\ensuremath{\beta}}
\newunicodechar{γ}{\ensuremath{\gamma}}
\newunicodechar{δ}{\ensuremath{\delta}}
\newunicodechar{ε}{\ensuremath{\varepsilon}}
\newunicodechar{θ}{\ensuremath{\theta}}
\newunicodechar{λ}{\ensuremath{\lambda}}
\newunicodechar{μ}{\ensuremath{\mu}}
\newunicodechar{σ}{\ensuremath{\sigma}}
\newunicodechar{τ}{\ensuremath{\tau}}
\newunicodechar{φ}{\ensuremath{\varphi}}
\newunicodechar{ψ}{\ensuremath{\psi}}
\newunicodechar{ω}{\ensuremath{\omega}}
\newunicodechar{Σ}{\ensuremath{\Sigma}}
% majuscules produites par \MakeUppercase dans les titres (α-aminés -> Α)
\newunicodechar{Α}{\ensuremath{\alpha}}
\newunicodechar{Β}{\ensuremath{\beta}}
\newunicodechar{Γ}{\ensuremath{\Gamma}}
\newunicodechar{Δ}{\ensuremath{\Delta}}
\newunicodechar{Ε}{\ensuremath{\varepsilon}}
\newunicodechar{Θ}{\ensuremath{\Theta}}
\newunicodechar{Λ}{\ensuremath{\Lambda}}
\newunicodechar{Μ}{\ensuremath{\mu}}
\newunicodechar{Τ}{\ensuremath{\tau}}
\newunicodechar{Φ}{\ensuremath{\Phi}}
\newunicodechar{Ψ}{\ensuremath{\Psi}}
\newunicodechar{Ω}{\ensuremath{\Omega}}
\newunicodechar{↦}{\ensuremath{\mapsto}}
\newunicodechar{∈}{\ensuremath{\in}}
\newunicodechar{∉}{\ensuremath{\notin}}
\newunicodechar{∧}{\ensuremath{\wedge}}
\newunicodechar{⇔}{\ensuremath{\Leftrightarrow}}
\newunicodechar{⟺}{\ensuremath{\Longleftrightarrow}}
\newunicodechar{⇒}{\ensuremath{\Rightarrow}}
\newunicodechar{⟹}{\ensuremath{\Longrightarrow}}
\newunicodechar{∘}{\ensuremath{\circ}}
\newunicodechar{∪}{\ensuremath{\cup}}
\newunicodechar{∩}{\ensuremath{\cap}}
\newunicodechar{≡}{\ensuremath{\equiv}}
\newunicodechar{⊂}{\ensuremath{\subset}}
\newunicodechar{⊥}{\ensuremath{\perp}}
\newunicodechar{∃}{\ensuremath{\exists}}
\newunicodechar{∀}{\ensuremath{\forall}}
\newunicodechar{∛}{\ensuremath{\sqrt[3]{\,}}}
\newunicodechar{∜}{\ensuremath{\sqrt[4]{\,}}}
\newunicodechar{⃗}{\ensuremath{{}^{\to}}}
\newunicodechar{ⁿ}{\ifmmode^{n}\else\textsuperscript{\ensuremath{n}}\fi}
\newunicodechar{ˣ}{\ifmmode^{x}\else\textsuperscript{\ensuremath{x}}\fi}
\newunicodechar{ᵇ}{\ifmmode^{b}\else\textsuperscript{\ensuremath{b}}\fi}
\newunicodechar{ʳ}{\ifmmode^{r}\else\textsuperscript{\ensuremath{r}}\fi}
\newunicodechar{ᵘ}{\ifmmode^{u}\else\textsuperscript{\ensuremath{u}}\fi}
\newunicodechar{ʸ}{\ifmmode^{y}\else\textsuperscript{\ensuremath{y}}\fi}
\newunicodechar{ᵃ}{\ifmmode^{a}\else\textsuperscript{\ensuremath{a}}\fi}
\newunicodechar{ᵉ}{\ifmmode^{e}\else\textsuperscript{\ensuremath{e}}\fi}
\newunicodechar{ᵏ}{\ifmmode^{k}\else\textsuperscript{\ensuremath{k}}\fi}
\newunicodechar{ᵖ}{\ifmmode^{p}\else\textsuperscript{\ensuremath{p}}\fi}
\newunicodechar{ᴺ}{\ifmmode^{N}\else\textsuperscript{\ensuremath{N}}\fi}
\newunicodechar{ᶜ}{\ifmmode^{c}\else\textsuperscript{\ensuremath{c}}\fi}
\newunicodechar{ʰ}{\ifmmode^{h}\else\textsuperscript{\ensuremath{h}}\fi}
\newunicodechar{ᵠ}{\ifmmode^{q}\else\textsuperscript{\ensuremath{q}}\fi}
\newunicodechar{ₙ}{\ifmmode_{n}\else\textsubscript{\ensuremath{n}}\fi}
\newunicodechar{ₐ}{\ifmmode_{a}\else\textsubscript{\ensuremath{a}}\fi}
\newunicodechar{ᵢ}{\ifmmode_{i}\else\textsubscript{\ensuremath{i}}\fi}
\newunicodechar{ᵣ}{\ifmmode_{r}\else\textsubscript{\ensuremath{r}}\fi}
\newunicodechar{ₖ}{\ifmmode_{k}\else\textsubscript{\ensuremath{k}}\fi}
\newunicodechar{ᵦ}{\ifmmode_{\beta}\else\textsubscript{\ensuremath{\beta}}\fi}
\newunicodechar{ₓ}{\ifmmode_{x}\else\textsubscript{\ensuremath{x}}\fi}
\newfontfamily\popdisplay{ArchivoBlack-Regular.ttf}[Path=../../../fonts/]
\newfontfamily\poplogo{SpaceGrotesk-Bold.ttf}[Path=../../../fonts/]
\newfontfamily\popsemi{PlusJakartaSans-SemiBold.ttf}[Path=../../../fonts/]
\newfontfamily\popxbold{PlusJakartaSans-ExtraBold.ttf}[Path=../../../fonts/]
\newfontfamily\popxboldls{PlusJakartaSans-ExtraBold.ttf}[Path=../../../fonts/,LetterSpace=3.0]

\setlist[enumerate]{leftmargin=1.7em,itemsep=5pt,topsep=4pt}
\setlist[itemize]{leftmargin=1.7em,itemsep=4pt,topsep=4pt,label={\color{poporange}\textbullet}}
\setlength{\parindent}{0pt}
\setlength{\parskip}{5pt}
\renewcommand{\arraystretch}{1.25}

% pgfplots : style Pop par défaut
\pgfplotsset{
  every axis/.append style={axis line style={popink,line width=1pt},tick style={popink},
    grid style={popline,line width=0.5pt},label style={font=\small},tick label style={font=\footnotesize}},
  popcurve/.style={poporange,line width=1.8pt,no markers,smooth},
}
% Même style disponible dans un \draw TikZ simple (hors pgfplots)
\tikzset{popcurve/.style={poporange,line width=1.8pt}}
\tikzset{popgrid/.style={popline,very thin}}
\colorlet{popcurve}{poporange} % le nom du style est parfois utilisé comme couleur

% ============================================================
% Pill / squiggle
% ============================================================
\newcommand{\poppill}[3]{%
  \tikz[baseline=-0.55ex]\node[draw=popink,line width=1.2pt,rounded corners=2.6pt,%
    fill=#2,inner xsep=6.5pt,inner ysep=3pt]{%
    \popxboldls\fontsize{7}{7}\selectfont\color{#3}\MakeUppercase{#1}};%
}
\newcommand{\popsquiggle}[1]{%
  \tikz[baseline=-0.4ex]{
    \draw[#1,line width=2.6pt,line cap=round]
      plot [smooth] coordinates {(0,0.09) (0.34,0.01) (0.68,0.21) (1.02,0.01) (1.36,0.10)};
  }%
}
% Mot-clé surligné (marqueur orange sous le texte)
\newcommand{\cle}[1]{{\popxbold\color{popdark}#1}}

% ============================================================
% En-tête / pied
% ============================================================
\pagestyle{fancy}
\fancyhf{}
\renewcommand{\headrulewidth}{0pt}
\renewcommand{\footrulewidth}{0pt}
\lhead{\raisebox{-0.15ex}{\poplogo\fontsize{13}{13}\selectfont\color{popink}OnBuch\color{poporange}+}}
\rhead{\poppill{\DOCMATIERE\ \textbullet\ \DOCNIVEAU\ \textbullet\ Leçon \DOCLECON}{white}{popink}}
\lfoot{\popxbold\fontsize{7.6}{7.6}\selectfont\color{popmuted}\MakeUppercase{Cours signé OnBuch+}\ \textbullet\ {\popsemi\DOCTITRE}}
\rfoot{\tikz[baseline=-0.6ex]\node[rounded corners=8.5pt,fill=popink,minimum height=17pt,minimum width=17pt,inner xsep=5pt,inner sep=0]{\popxbold\fontsize{8}{8}\selectfont\color{popcream}\thepage\,/\,\pageref*{LastPage}};}

% ============================================================
% Titres
% ============================================================
\newcommand{\chaptitre}[2]{%
  \begin{tcolorbox}[enhanced,colback=poporange,colframe=popink,boxrule=2.6pt,arc=11pt,
      drop shadow=popink,left=15pt,right=15pt,top=12pt,bottom=11pt,before skip=3pt,after skip=18pt]
    \poppill{#2}{popink}{popcream}\par\vspace{9pt}
    {\raggedright\hyphenpenalty=10000\exhyphenpenalty=10000\popdisplay\fontsize{22}{24}\selectfont\color{popcream}\MakeUppercase{#1}\par}
  \end{tcolorbox}
}
% Section numérotée : pastille orange avec le numéro + titre Archivo
\newcounter{popsec}
\newcommand{\coursec}[1]{%
  \stepcounter{popsec}\setcounter{popsub}{0}%
  \par\vspace{16pt}\noindent
  \begin{minipage}{\linewidth}
  \tikz[baseline=-0.7ex]\node[draw=popink,line width=1.6pt,fill=poporange,rounded corners=4pt,minimum size=22pt,inner sep=0]{\popdisplay\fontsize{12}{12}\selectfont\color{popcream}\thepopsec};%
  \hspace{8pt}{\popdisplay\fontsize{15}{17}\selectfont\color{popink}\MakeUppercase{#1}}\par
  \vspace{4pt}\noindent{\color{poporange}\rule{36pt}{3.6pt}}
  \end{minipage}\par\nopagebreak\vspace{8pt}
}
\newcounter{popsub}
\newcommand{\courssub}[1]{%
  \stepcounter{popsub}%
  \par\vspace{9pt}\noindent{\popxbold\fontsize{11.5}{13}\selectfont\color{popink}{\color{poporange}\thepopsec.\thepopsub}\hspace{6pt}#1}\par\nopagebreak\vspace{4pt}
}

% ============================================================
% Boîtes pédagogiques — même recette : fond blanc, bordure épaisse colorée,
% ombre dure encre, badge plein. Titre optionnel [#1].
% ============================================================
\tcbset{popbase/.style={breakable,enhanced,colback=white,boxrule=2.3pt,arc=9pt,
  shadow={3pt}{-3pt}{0pt}{popink},left=14pt,right=14pt,top=11pt,bottom=11pt,before skip=11pt,after skip=15pt,
  fontupper=\popsemi\fontsize{10.6}{14.5}\selectfont\color{popink},
  fontlower=\popsemi\fontsize{10.6}{14.5}\selectfont\color{popink}}}
% Coche dessinée (le glyphe ✓ n'existe pas dans les polices du document)
\newcommand{\popcheck}[1]{\tikz[baseline=-0.5ex]\draw[#1,line width=1.6pt,line cap=round,line join=round](-0.13,0.0)--(-0.04,-0.09)--(0.13,0.1);}
\providecommand{\coche}{\popcheck{popgreen}}
\providecommand{\resultat}[1]{\par\smallskip\noindent\fcolorbox{popgreen}{popgreenL}{\parbox{\dimexpr\linewidth-2\fboxsep-2\fboxrule\relax}{\textbf{Résultat.} #1}}\par\smallskip}
\newcommand{\popboxhead}[3]{\poppill{#1}{#2}{white}\ifx\relax#3\relax\else\hspace{6pt}{\popxbold\fontsize{10.6}{12}\selectfont\color{popink}#3}\fi\par\vspace{7pt}}

\newtcolorbox{aretenir}[1][]{popbase,colframe=popgreen,before upper={\popboxhead{À retenir}{popgreen}{#1}}}
\newtcolorbox{exemplebox}[1][]{popbase,colframe=poporange,before upper={\popboxhead{Exemple}{poporange}{#1}}}
\newtcolorbox{definition}[1][]{popbase,colframe=popblue,before upper={\popboxhead{Définition}{popblue}{#1}}}
\newtcolorbox{propriete}[1][]{popbase,colframe=poppurple,before upper={\popboxhead{Propriété}{poppurple}{#1}}}
\newtcolorbox{methode}[1][]{popbase,colframe=popgold,before upper={\popboxhead{Méthode}{popgold}{#1}}}
\newtcolorbox{attention}[1][]{popbase,colframe=poppink,before upper={\popboxhead{Attention}{poppink}{#1}}}
\newtcolorbox{experience}[1][]{popbase,colframe=popblue,colback=popblueL,before upper={\popboxhead{Expérience}{popblue}{#1}}}
% « Le savais-tu ? » : carte violette pleine (contexte, culture, Cameroun)
\newtcolorbox{savaistu}[1][]{popbase,colback=poppurple,colframe=popink,
  fontupper=\popsemi\fontsize{10.4}{14.2}\selectfont\color{white},
  before upper={\poppill{Le savais-tu\,?}{white}{poppurple}\ifx\relax#1\relax\else\hspace{6pt}{\popxbold\color{white}#1}\fi\par\vspace{7pt}}}
% Exercice résolu : énoncé en haut, \tcblower puis la solution sur fond crème
\newcounter{popexo}\newcounter{popexr}
\newtcolorbox{exoresolu}[1][]{popbase,colframe=poporange,colbacklower=popcream,
  segmentation style={popink,line width=1.2pt,dashed},
  before upper={\stepcounter{popexr}\popboxhead{Exercice résolu \thepopexr}{poporange}{#1}},
  before lower={\poppill{Solution}{popink}{popcream}\par\vspace{6pt}}}
% Exercice (non résolu en place) — la correction est dans la partie Corrigés
\newtcolorbox{exercice}[1][]{popbase,colframe=popink,
  before upper={\stepcounter{popexo}\popboxhead{Exercice \thepopexo}{popink}{#1}}}
% Corrigé
\newtcolorbox{corrige}[1][]{popbase,colframe=popgreen,colback=popgreenL,
  before upper={\popboxhead{Corrigé}{popgreen}{#1}}}
% Objectifs / prérequis / bilan
\newtcolorbox{objectifs}{popbase,colframe=popink,colback=poporangeL,
  before upper={\popboxhead{Objectifs de la leçon}{popink}{}}}
\newtcolorbox{prerequis}{popbase,colframe=popink,colback=popblueL,
  before upper={\popboxhead{Prérequis}{popblue}{}}}
\newtcolorbox{bilan}{popbase,colframe=popink,colback=white,boxrule=2.8pt,
  before upper={\popboxhead{Fiche bilan}{poporange}{L'essentiel à connaître pour l'examen}}}
% Figure encadrée avec légende
\newtcolorbox{popfigure}[1][]{enhanced,breakable=false,colback=white,colframe=popline,boxrule=1.4pt,arc=8pt,
  left=10pt,right=10pt,top=10pt,bottom=8pt,before skip=10pt,after skip=12pt,halign=center,
  lower separated=false,halign lower=center,
  fontlower=\popsemi\fontsize{9}{11.5}\selectfont\color{popmuted},
  #1}
\newcommand{\legende}[1]{\par\vspace{6pt}{\popsemi\fontsize{9}{11.5}\selectfont\color{popmuted}#1}}

% ============================================================
% Signature OnBuch+
% ============================================================
\newcommand{\poplogoinline}[1]{{\poplogo\fontsize{#1}{#1}\selectfont\color{popink}OnBuch\color{poporange}+}}
\newcommand{\signatureonbuch}{%
  \begin{tcolorbox}[enhanced,colback=popink,colframe=popink,boxrule=2.2pt,arc=10pt,
      drop shadow=poporange,left=14pt,right=14pt,top=10pt,bottom=10pt,before skip=0pt,after skip=0pt]
    \tikz[baseline=-0.7ex]\node[circle,fill=popgreen,draw=popcream,line width=1.3pt,minimum size=22pt,inner sep=0]{\popcheck{white}};%
    \hspace{9pt}\begin{minipage}[c]{0.62\linewidth}
      {\popxbold\fontsize{12}{13}\selectfont\color{popcream}Signé\ \textbullet\ {\poplogo OnBuch\color{poporange}+}}\par\vspace{2pt}
      {\raggedright\popsemi\fontsize{8}{10.5}\selectfont\color{popline}\MakeUppercase{Équipe pédagogique OnBuch+}\\\MakeUppercase{Conforme au programme officiel MINESEC}\par}
    \end{minipage}\hfill
    {\poplogo\fontsize{22}{22}\selectfont\color{popcream}OnBuch\color{poporange}+}
  \end{tcolorbox}%
}

% ============================================================
% Page de garde
% #1 titre · #2 matière · #3 niveau · #4 module · #5 n° leçon · #6 durée ·
% #7 description / objectifs (texte ou itemize)
% ============================================================
\newcommand{\pagegarde}[7]{%
  \thispagestyle{empty}
  \newgeometry{a4paper,margin=1.7cm,top=1.6cm,bottom=1.4cm}%
  \begin{tikzpicture}[remember picture,overlay]
    \fill[poporange!18] ([xshift=-3.2cm,yshift=-3.6cm]current page.north east) circle (4.6cm);
    \fill[poppurple!14] ([xshift=2.4cm,yshift=7.6cm]current page.south west) circle (3.8cm);
  \end{tikzpicture}%
  {\poplogo\fontsize{30}{30}\selectfont\color{popink}OnBuch\color{poporange}+}\hfill
  \raisebox{0.6ex}{\poppill{Cours signé \textbullet\ Premium}{popink}{popcream}}\par
  \vspace{1pt}\popsquiggle{poppurple}\par
  \vspace{8pt}
  \begin{tcolorbox}[enhanced,colback=poporange,colframe=popink,boxrule=2.8pt,arc=13pt,
      drop shadow=popink,left=18pt,right=18pt,top=12pt,bottom=13pt,after skip=10pt]
    \poppill{#2 \textbullet\ #3}{popink}{popcream}\hspace{5pt}\poppill{#4}{white}{popink}\par\vspace{8pt}
    {\popdisplay\fontsize{14}{14}\selectfont\color{popink}LEÇON #5}\par\vspace{4pt}
    {\raggedright\hyphenpenalty=10000\exhyphenpenalty=10000\popdisplay\fontsize{25}{27}\selectfont\color{popcream}\MakeUppercase{#1}\par}
    \vspace{8pt}
    {\popxbold\fontsize{9.5}{9.5}\selectfont\color{popink}\MakeUppercase{Durée conseillée : #6}}
  \end{tcolorbox}
  \begin{tcolorbox}[enhanced,colback=white,colframe=popink,boxrule=2.2pt,arc=10pt,
      drop shadow=popink,left=16pt,right=16pt,top=10pt,bottom=10pt,after skip=8pt]
    \poppill{Dans cette leçon}{poppurple}{white}\par\vspace{5pt}
    \popsemi\fontsize{9.3}{12.6}\selectfont\color{popink}#7
  \end{tcolorbox}
  \noindent\begin{minipage}[t]{0.487\linewidth}
    \begin{tcolorbox}[enhanced,colback=popgreen,colframe=popink,boxrule=2.2pt,arc=10pt,
        drop shadow=popink,left=11pt,right=11pt,top=10pt,bottom=10pt,height=3.1cm,valign=center]
      \tcbox[colback=white,colframe=popink,boxrule=1.3pt,arc=5pt,boxsep=0pt,nobeforeafter,
        left=3pt,right=3pt,top=3pt,bottom=3pt,valign=center]{\qrcode[height=1.9cm]{https://play.google.com/store/apps/details?id=com.luvvix.onbuch&hl=fr}}%
      \hspace{8pt}\begin{minipage}[c]{0.5\linewidth}
        {\popdisplay\fontsize{11}{12.5}\selectfont\color{white}\MakeUppercase{Télécharge OnBuch}}\par\vspace{4pt}
        {\raggedright\popsemi\fontsize{8.4}{11}\selectfont\color{white}Gratuit \textbullet\ Google Play\\Tuteur IA, annales, concours, résultats\par}
      \end{minipage}
    \end{tcolorbox}
  \end{minipage}\hfill
  \begin{minipage}[t]{0.487\linewidth}
    \begin{tcolorbox}[enhanced,colback=poppurple,colframe=popink,boxrule=2.2pt,arc=10pt,
        drop shadow=popink,left=11pt,right=11pt,top=10pt,bottom=10pt,height=3.1cm,valign=center]
      \tikz[baseline=-0.7ex]\node[circle,fill=white,draw=popink,line width=1.3pt,minimum size=1.6cm,inner sep=0]{\popdisplay\fontsize{24}{24}\selectfont\color{poppurple}+};%
      \hspace{8pt}\begin{minipage}[c]{0.6\linewidth}
        {\popdisplay\fontsize{11}{12.5}\selectfont\color{white}\MakeUppercase{Passe à OnBuch+}}\par\vspace{4pt}
        {\raggedright\popsemi\fontsize{8.4}{11}\selectfont\color{white}Tous les cours signés, Tuteur IA illimité, fiches PDF — onglet OnBuch+ de l'app\par}
      \end{minipage}
    \end{tcolorbox}
  \end{minipage}
  \vspace{8pt}
  \popcontenu
  \vfill
  \signatureonbuch
  \restoregeometry
  \newpage
}

% Rangée « Ce cours contient » (page de garde)
\newcommand{\poptile}[3]{% #1 couleur #2 gros texte #3 légende
  \begin{tcolorbox}[enhanced,colback=#1,colframe=popink,boxrule=2pt,arc=9pt,shadow={2.5pt}{-2.5pt}{0pt}{popink},
      left=6pt,right=6pt,top=9pt,bottom=9pt,halign=center,valign=center,height=1.75cm,width=\linewidth,nobeforeafter]
    {\popdisplay\fontsize{12.5}{13}\selectfont\color{white}\MakeUppercase{#2}}\par\vspace{3pt}
    {\centering\popsemi\fontsize{7.8}{9.5}\selectfont\color{white}#3\par}
  \end{tcolorbox}}
\newcommand{\popcontenu}{%
  \par\noindent{\popxboldls\fontsize{7.5}{7.5}\selectfont\color{popmuted}CE COURS SIGNÉ CONTIENT}\par\vspace{7pt}
  \noindent\begin{minipage}[t]{0.235\linewidth}\poptile{popblue}{Cours}{complet, illustré, pas à pas}\end{minipage}\hfill
  \begin{minipage}[t]{0.235\linewidth}\poptile{poporange}{Méthodes}{et exercices résolus}\end{minipage}\hfill
  \begin{minipage}[t]{0.235\linewidth}\poptile{poppink}{Exercices}{type examen + corrigés}\end{minipage}\hfill
  \begin{minipage}[t]{0.235\linewidth}\poptile{popgreen}{Fiche bilan}{l'essentiel à retenir}\end{minipage}\par}

% Sommaire
\newtcolorbox{sommaire}{enhanced,colback=white,colframe=popink,boxrule=2.2pt,arc=10pt,
  drop shadow=popink,left=18pt,right=18pt,top=13pt,bottom=13pt,before skip=6pt,after skip=20pt,
  before upper={\poppill{Sommaire}{poppurple}{white}\par\vspace{9pt}\popsemi\fontsize{10.6}{14.5}\selectfont\color{popink}}}

% Signature de fin de cours — poussée tout en bas de la dernière page
\newcommand{\findecours}{%
  \par\vfill\nopagebreak
  \begin{center}\popsquiggle{poporange}\end{center}\vspace{4pt}
  \signatureonbuch
}
\AtBeginDocument{\def\llbracket{\mathopen{[\mkern-3.5mu[}}\def\rrbracket{\mathclose{]\mkern-3.5mu]}}} % intervalles d'entiers (glyphe ⟦ absent de la police)
% Glyphes absents de Fira Math : redéfinis à partir de glyphes disponibles
\AtBeginDocument{%
  \def\setminus{\mathbin{\backslash}}\let\smallsetminus\setminus
  \def\vdots{\mathord{\vbox{\baselineskip3.2pt\lineskiplimit0pt\kern4pt\hbox{.}\hbox{.}\hbox{.}}}}%
  \def\ddots{\mathinner{\mkern1mu\raise6.5pt\hbox{.}\mkern2mu\raise3.6pt\hbox{.}\mkern2mu\raise0.7pt\hbox{.}\mkern1mu}}%
  \def\triangle{\mathord{\scalebox{0.9}{$\symup{\Delta}$}}}%
  \def\nabla{\mathord{\raisebox{\depth}{\scalebox{1}[-1]{$\symup{\Delta}$}}}}%
  \def\checkmark{\popcheck{popdark}}%
  \def\star{\mathbin{\ast}}\def\bigstar{\mathord{\ast}}%
  \def\diamond{\mathbin{\scalebox{0.6}{$\Diamond$}}}%
  \def\bigcup{\mathop{\mathchoice{\raisebox{-0.3em}{\scalebox{1.7}{$\cup$}}}{\scalebox{1.25}{$\cup$}}{\cup}{\cup}}\displaylimits}%
  \def\bigcap{\mathop{\mathchoice{\raisebox{-0.3em}{\scalebox{1.7}{$\cap$}}}{\scalebox{1.25}{$\cap$}}{\cap}{\cap}}\displaylimits}%
  \def\lesseqqgtr{\mathrel{\lessgtr}}\def\bigoplus{\mathop{\oplus}\displaylimits}%
}
\providecommand{\ssi}{si et seulement si} % abréviation fréquente
\AtBeginDocument{\providecommand{\wideparen}{\overparen}} % arc de cercle

%% FIGFIX-BEGIN v5 — correctifs globaux des figures (docs/onbuch-generation/tools/figfix)
\usepackage{adjustbox}\usepackage{etoolbox}\tcbuselibrary{raster}
% toute figure TikZ/pgfplots plus large que la ligne est réduite à la largeur disponible
\BeforeBeginEnvironment{tikzpicture}{\begin{adjustbox}{max width=\linewidth}}
\AfterEndEnvironment{tikzpicture}{\end{adjustbox}}
% légendes pgfplots lisibles par-dessus les courbes
\pgfplotsset{every axis/.append style={legend style={fill=white,fill opacity=0.92,text opacity=1,draw=black!15,inner sep=3pt}}}
% étiquettes d'arêtes (syntaxe edge["..."]) avec fond blanc
\tikzset{every edge quotes/.append style={fill=white,inner sep=1.2pt}}
%% FIGFIX-END
\begin{document}

\pagegarde{Dieu et la religion}{Philosophie}{Tle Tech}{Philosophie – classe de Terminale technique}{N12}{4 séances (fiche de progression harmonisée)}{Cette leçon explore philosophiquement le problème de l'existence de Dieu et la valeur de la religion. Elle analyse l'origine du sentiment religieux, les grandes conceptions du divin, les preuves classiques et leurs critiques, pour aboutir à une compréhension critique du fait religieux dans l'espace public camerounais, en conformité avec les exigences du Baccalauréat technique.
\vspace{4pt}
\begin{multicols}{2}\small
\begin{itemize}
\item Analyser les théories explicatives de l'origine du sentiment religieux (psychologiques, sociologiques, anthropologiques).
\item Différencier les principales conceptions de Dieu (théisme, déisme, panthéisme, athéisme, agnosticisme) et situer leur portée logique.
\item Reconstituer et critiquer les grandes preuves rationnelles de l'existence de Dieu (ontologique, cosmologique, physico-théologique, morale) à la lumière de Kant et des modernes.
\item Distinguer rigoureusement religion, superstition, magie et spiritualité à l'aide de critères conceptuels.
\item Évaluer la valeur sociale de la religion (cohésion, conflit, éthique, aliénation) en s'appuyant sur Marx, Durkheim, Freud et Weber.
\item Appliquer le principe de laïcité à des situations concrètes camerounaises (gestion du pluralisme, éducation, droit de la famille).
\end{itemize}\end{multicols}}

\begin{sommaire}
\begin{multicols}{2}
\begin{enumerate}
\item 1. L'origine du sentiment religieux : peur, merveilleux ou projection sociale ?
\item 2. Les conceptions de Dieu : une cartographie conceptuelle
\item 3. L'existence de Dieu : les preuves classiques et leurs critiques
\item 4. La religion : définition, formes et valeur sociale
\item 5. Foi, Raison et Laïcité : le défi camerounais
\item 6. Atelier méthodologique : Dissertation et Explication de texte (Bac Tech)
\item Activité d'intégration
\item Méthodes et exercices résolus
\item Exercices
\item Corrigés des exercices
\item Fiche bilan
\end{enumerate}
\end{multicols}
\end{sommaire}

\begin{objectifs}
\begin{itemize}
\item Analyser les théories explicatives de l'origine du sentiment religieux (psychologiques, sociologiques, anthropologiques).
\item Différencier les principales conceptions de Dieu (théisme, déisme, panthéisme, athéisme, agnosticisme) et situer leur portée logique.
\item Reconstituer et critiquer les grandes preuves rationnelles de l'existence de Dieu (ontologique, cosmologique, physico-théologique, morale) à la lumière de Kant et des modernes.
\item Distinguer rigoureusement religion, superstition, magie et spiritualité à l'aide de critères conceptuels.
\item Évaluer la valeur sociale de la religion (cohésion, conflit, éthique, aliénation) en s'appuyant sur Marx, Durkheim, Freud et Weber.
\item Appliquer le principe de laïcité à des situations concrètes camerounaises (gestion du pluralisme, éducation, droit de la famille).
\item Construire une argumentation philosophique structurée (dissertation ou explication de texte) mobilisant concepts, auteurs et exemples locaux.
\end{itemize}
\end{objectifs}

\begin{prerequis}
\begin{itemize}
\item Distinction savoir / croire / opinion (Programme de Première).
\item Notions de subjectivité, conscience et désir (Unité 1 de Terminale).
\item Maîtrise des structures logiques de base (déduction, induction, analogie, sophismes).
\item Connaissances générales sur les trois grandes religions monothéistes et les religions traditionnelles au Cameroun.
\item Lecture de la Constitution camerounaise (Préambule) sur la laïcité et la liberté de conscience.
\end{itemize}
\end{prerequis}

\newpage

\coursec{L'origine du sentiment religieux : peur, merveilleux ou projection sociale ?}

Au carrefour de la mosquée centrale, de la cathédrale Notre-Dame-des-Victoires et des lieux de culte traditionnels de Yaoundé, des milliers de Camerounais prient chaque vendredi, chaque dimanche ou lors des cérémonies coutumières, selon des rites différents. Les appels du muezzin, les cloches des églises et les tam-tams des rituels traditionnels rythment la même ville. Pourtant, la Constitution du 18 janvier 1996 proclame dans son préambule la laïcité de l'État et garantit la liberté de conscience et de culte. Comment comprendre cette coexistence entre une ferveur religieuse omniprésente et la neutralité de l'État ? Le sentiment religieux est-il une nécessité anthropologique, inscrite dans la nature humaine, ou une construction sociale, un produit de l'histoire et des rapports de force ?

\begin{savaistu}[D'où vient le mot « religion » ?]
Le mot « religion » vient du latin, mais son étymologie est disputée. Cicéron le rattache à \emph{religere} (recueillir, relire, respecter scrupuleusement les rites), tandis que Lactance, auteur chrétien du IV\textsuperscript{e} siècle, le rattache à \emph{religare} (relier l'homme au divin). Ce débat étymologique reflète le débat philosophique lui-même : la religion est-elle d'abord un ensemble de \cle{devoirs} et de rites à accomplir, ou un \cle{lien} vécu avec le divin ?
\end{savaistu}

\courssub{L'expérience du sacré : la définition phénoménologique}

Avant d'expliquer la religion, il faut d'abord la décrire telle qu'elle se présente à celui qui la vit. C'est la démarche \cle{phénoménologique} : suspendre provisoirement le jugement sur la vérité des croyances pour décrire l'expérience religieuse de l'intérieur.

L'intuition de départ est simple : entrer dans une cathédrale silencieuse, assister à une veillée traditionnelle ou contempler les chutes de la Lobé provoque un sentiment particulier, mélange de crainte respectueuse et d'attraction irrésistible. Ce sentiment n'est ni la peur ordinaire ni l'admiration esthétique : c'est l'expérience du \cle{sacré}.

\begin{definition}[Le \emph{numinosum} selon Rudolf Otto]
Dans \emph{Le Sacré} (1917), le théologien et philosophe allemand Rudolf Otto décrit le \emph{numinosum} (du latin \emph{numen}, puissance divine) comme une catégorie propre de l'esprit humain : le sentiment du « tout autre » (\emph{ganz Andere}), d'une réalité radicalement différente de nous. Cette expérience a deux pôles inséparables : le \emph{mysterium tremendum} (le mystère qui fait trembler, qui écrase par sa majesté) et le \emph{mysterium fascinans} (le mystère qui fascine, attire et comble). Le sacré est donc à la fois redoutable et désirable.
\end{definition}

L'historien des religions Mircea Eliade complète cette analyse par une distinction fondatrice : pour l'homme religieux, le monde n'est pas homogène. Il se partage entre le \cle{sacré} (ce qui est mis à part, interdit ou vénéré : le temple, le bois sacré, le jour saint) et le \cle{profane} (l'espace ordinaire, quotidien, sans valeur religieuse). Chez de nombreux peuples camerounais, la forêt sacrée où résident les ancêtres n'est pas un lieu comme les autres : on n'y entre qu'avec des rites précis. Le sacré structure l'espace, le temps (fêtes, jours de culte) et la vie sociale.

\begin{attention}[Piège fréquent]
Ne pas confondre \cle{critique de la religion} et \cle{athéisme}. Critiquer la religion, c'est soumettre rationnellement ses croyances, ses institutions ou sa valeur sociale à l'examen ; l'athéisme est la négation de l'existence de Dieu. On peut critiquer l'institution religieuse tout en étant croyant : c'est le cas de Luther réformant l'Église, ou des prophètes bibliques dénonçant les pratiques de leur peuple.
\end{attention}

\courssub{L'approche psychologique : Freud et la religion comme illusion}

Pourquoi l'homme est-il religieux ? La première grande réponse est psychologique : la religion répondrait à un \cle{besoin} de l'âme humaine.

Dans \emph{L'Avenir d'une illusion} (1927), Sigmund Freud soutient que la religion est une \cle{illusion} (non pas nécessairement une erreur, mais une croyance motivée par un désir). Son raisonnement part d'un constat : l'homme est fondamentalement impuissant devant les forces de la nature — maladies, tempêtes, mort — et devant la souffrance inévitable de l'existence. Comme l'enfant impuissant se réfugie auprès de son père protecteur, l'humanité enfantine projette dans le ciel un \cle{Père tout-puissant} qui la protège, la console et promet la justice dans l'au-delà.

Freud va plus loin : il compare les rites religieux (gestes répétitifs, interdits, purifications) aux symptômes de la névrose obsessionnelle. La religion serait ainsi une \cle{névrose obsessionnelle collective} de l'humanité, un stade infantile que l'espèce humaine, devenue adulte grâce à la science, devrait dépasser. D'où le titre provocateur de l'ouvrage : l'illusion religieuse a-t-elle un \emph{avenir} ?

\begin{exemplebox}[Illustration camerounaise]
Lors des épidémies ou des deuils, on observe souvent un recours accru à la prière, aux veillées et aux rites de protection. Pour Freud, ce comportement confirme sa thèse : plus l'angoisse et le sentiment d'impuissance grandissent, plus le besoin d'un protecteur divin s'intensifie. La religion serait alors un mécanisme de consolation face à la détresse.
\end{exemplebox}

La critique de cette approche est classique : expliquer l'\emph{origine psychologique} d'une croyance ne tranche pas la question de sa \emph{vérité}. Qu'une croyance console ne prouve ni qu'elle est vraie ni qu'elle est fausse : c'est le problème de la « faute génétique » (juger une idée par son origine plutôt que par ses arguments).

\courssub{L'approche sociologique : Durkheim et la société qui se divinise}

La deuxième grande réponse est sociologique. Dans \emph{Les Formes élémentaires de la vie religieuse} (1912), Émile Durkheim étudie le totémisme des aborigènes d'Australie, qu'il considère comme la religion la plus simple, la forme « élémentaire » permettant de comprendre toutes les autres.

Son analyse est rigoureuse : dans le clan, le \cle{totem} (un animal ou une plante emblème du groupe) est à la fois l'objet du culte et le symbole du clan lui-même. Or, se demande Durkheim, d'où vient la force que les fidèles éprouvent dans les cérémonies collectives, cette exaltation qu'il nomme « effervescence collective » ? Elle vient du groupe rassemblé. Conclusion célèbre : en adorant le totem, le clan adore en réalité \emph{lui-même}, transfiguré. \textbf{La religion est la société qui se divinise.}

La fonction de la religion est alors claire : elle produit la \cle{conscience collective}, c'est-à-dire l'ensemble des croyances et des sentiments communs qui font la cohésion du groupe. Les rites resserrent les liens, les interdits sacrés (tabous) organisent la vie commune, les fêtes renouvellent l'appartenance.

\begin{exemplebox}[Le totem et le clan au Cameroun]
Chez plusieurs peuples camerounais, certains animaux (panthère, python, éléphant) servent d'emblèmes à des lignages et font l'objet d'interdits alimentaires ou de vénération. De même, les cérémonies traditionnelles (funérailles, rites d'initiation, cultes des ancêtres) rassemblent périodiquement la communauté et réaffirment son unité : illustration presque parfaite de la thèse durkheimienne de la cohésion sociale par le sacré.
\end{exemplebox}

Mais la critique s'impose : si Dieu n'est que la société déguisée, comment expliquer les prophètes et les mystiques qui s'opposent à leur société au nom de Dieu ? La religion ne se réduit peut-être pas à une fonction sociale utile.

\courssub{L'approche anthropologique : de l'animisme au monothéisme}

Une troisième réponse, d'inspiration évolutionniste, est proposée par l'anthropologue anglais Edward Burnett Tylor (\emph{Primitive Culture}, 1871). Pour lui, la religion naît d'une \cle{erreur intellectuelle} naïve mais explicable : l'homme primitif, cherchant à expliquer les rêves, la mort et les phénomènes naturels, suppose l'existence d'\emph{esprits} ou d'\emph{âmes} animant toutes choses. C'est l'\cle{animisme}, « religion minimum ». De là, l'humanité serait passée progressivement du polythéisme (plusieurs dieux) au monothéisme (un Dieu unique), puis à la science, selon une évolution unilinéaire vers la rationalité.

\begin{attention}[Critique de l'évolutionnisme unilinéaire]
Cette théorie présente un grave défaut : elle classe les religions du « plus primitif » au « plus évolué », ce qui dévalorise les religions traditionnelles. Or les religions traditionnelles camerounaises (cultes des ancêtres, divination, rites agraires) ne sont pas des « brouillons » de religion : ce sont des \cle{systèmes complets} et cohérents, avec leur théologie (Dieu créateur, esprits intermédiaires, ancêtres), leur morale et leur sagesse. Parler d'évolution unilinéaire relève de l'ethnocentrisme.
\end{attention}

\courssub{L'approche marxiste : la religion, « opium du peuple »}

Karl Marx, dans l'introduction à la \emph{Critique de la philosophie du droit de Hegel} (1844), formule la thèse la plus radicale : « La religion est le soupir de la créature accablée, le cœur d'un monde sans cœur, l'esprit d'un état de choses sans esprit. Elle est \textbf{l'opium du peuple}. »

Comprendre la métaphore : l'opium soulage la douleur sans guérir la maladie. De même, la religion console l'ouvrier exploité en lui promettant le bonheur au ciel, ce qui l'empêche de se révolter contre sa misère terrestre. La religion est une \cle{idéologie} : un monde renversé, reflet d'un monde réel renversé. Elle masque l'\cle{aliénation économique} (l'homme dépossédé du fruit de son travail) au lieu de la combattre. Pour Marx, ce n'est donc pas la religion qu'il faut d'abord critiquer, mais les conditions sociales qui la rendent nécessaire : supprimez la misère, et l'illusion s'évanouira d'elle-même.

La critique inverse est tout aussi forte : la religion a souvent été un ferment de révolte et de justice sociale (théologie de la libération en Amérique latine, engagement des Églises pour la paix et la démocratie au Cameroun), et non seulement un instrument d'endormissement.

\courssub{Synthèse comparative : quatre réponses à une même question}

\begin{center}
\begin{tabularx}{\linewidth}{|l|X|X|X|}
\hline
\rowcolor{popblueL} \textbf{Auteur} & \textbf{Origine du sentiment religieux} & \textbf{Concept clé} & \textbf{Critique possible} \\
\hline
Otto / Eliade & Expérience irréductible du sacré & \emph{Numinosum}, sacré/profane & Décrit sans expliquer l'origine \\
\hline
Freud & Besoin psychique de protection & Illusion, névrose collective & Faute génétique : origine $\neq$ vérité \\
\hline
Durkheim & Fonction sociale de cohésion & Conscience collective, totem & N'explique pas le mystique dissident \\
\hline
Marx & Aliénation économique & Idéologie, « opium du peuple » & La religion peut aussi libérer \\
\hline
\end{tabularx}
\end{center}

Le sentiment religieux est-il une \emph{erreur intellectuelle} (Tylor), un \emph{besoin psychique} (Freud), une \emph{fonction sociale} (Durkheim, Marx) ou une \emph{ouverture} à une réalité transcendante (Otto, les croyants eux-mêmes) ? La synthèse s'impose : ces explications ne sont pas forcément exclusives, car elles répondent à des questions différentes — le \emph{pourquoi} psychologique, le \emph{à quoi ça sert} sociologique, le \emph{comment} historique. Mais aucune ne tranche la question philosophique décisive : Dieu existe-t-il ? Ce sera l'objet de la suite de la leçon.

\begin{popfigure}
\begin{center}
\begin{tikzpicture}[
  racine/.style={rectangle, rounded corners, draw=popdark, fill=popblueL, text width=2.6cm, align=center, font=\small\bfseries},
  branche/.style={rectangle, rounded corners, draw=popdark, fill=poporangeL, text width=2.8cm, align=center, font=\small},
  fruit/.style={rectangle, rounded corners, draw=popdark, fill=popgreenL, text width=2.8cm, align=center, font=\small\itshape},
  lien/.style={-{Stealth[length=2mm]}, thick, popdark}
]
\node[racine] (r1) at (0,0) {Peur / Besoin psychique};
\node[racine] (r2) at (6,0) {Société};
\node[racine] (r3) at (12,0) {Projection / Aliénation};
\node[branche] (b1) at (0,2.6) {Freud : illusion, névrose collective};
\node[branche] (b2) at (6,2.6) {Durkheim : conscience collective, totem};
\node[branche] (b3) at (12,2.6) {Marx : idéologie, « opium du peuple »};
\node[fruit] (f1) at (0,5.2) {Critique : origine $\neq$ vérité};
\node[fruit] (f2) at (6,5.2) {Critique : n'explique pas le dissident};
\node[fruit] (f3) at (12,5.2) {Critique : la religion peut libérer};
\draw[lien] (r1) -- (b1);
\draw[lien] (r2) -- (b2);
\draw[lien] (r3) -- (b3);
\draw[lien] (b1) -- (f1);
\draw[lien] (b2) -- (f2);
\draw[lien] (b3) -- (f3);
\end{tikzpicture}
\end{center}
\legende{Arbre des théories de l'origine du sentiment religieux : les racines (causes supposées), les branches (auteurs et concepts clés) et les fruits (critiques respectives).}
\end{popfigure}

\begin{popfigure}
\begin{center}
\begin{tikzpicture}
\begin{axis}[
  ybar, bar width=14pt,
  width=0.85\linewidth, height=6.5cm,
  ymin=0, ymax=45,
  ylabel={Pourcentage de la population (\%)},
  symbolic x coords={Catholiques, Protestants, Musulmans, Animistes, Autres},
  xtick=data,
  x tick label style={font=\small, rotate=15, anchor=east},
  nodes near coords, nodes near coords style={font=\small},
  axis lines=left,
  enlarge x limits=0.15
]
\addplot[fill=popblue, draw=popdark] coordinates {(Catholiques,38.4) (Protestants,26.3) (Musulmans,20.9) (Animistes,5.6) (Autres,8.8)};
\end{axis}
\end{tikzpicture}
\end{center}
\legende{Répartition religieuse estimée de la population camerounaise (en \%, ordres de grandeur d'après le recensement BUCREP 2005 et les estimations démographiques postérieures) : environ 65\,\% de chrétiens (catholiques et protestants), 21\,\% de musulmans, 5 à 6\,\% d'animistes, le reste se répartissant entre autres religions et sans religion.}
\end{popfigure}

\begin{exemplebox}[Un pluralisme remarquable]
Les données disponibles (recensement BUCREP 2005 et estimations ultérieures) font état d'environ 65 à 70\,\% de chrétiens, 20 à 25\,\% de musulmans, environ 5\,\% d'animistes, le reste se répartissant entre autres religions et personnes sans religion. Le ministère de l'Administration territoriale (MINAT) recense par ailleurs un très grand nombre de groupements religieux légalement reconnus. Ce pluralisme exceptionnel fait du Cameroun un laboratoire vivant de la question religieuse.
\end{exemplebox}

\begin{definition}[Laïcité (contexte camerounais)]
La \cle{laïcité} désigne la neutralité de l'État en matière religieuse : séparation des institutions religieuses et politiques, liberté de conscience et de culte garantie à tous (préambule de la Constitution du 18 janvier 1996). Attention : la laïcité ne signifie \emph{pas} l'athéisme d'État ; elle protège au contraire la liberté religieuse de chaque citoyen en empêchant toute religion d'État.
\end{definition}

\begin{methode}[Analyser un texte sur l'origine de la religion]
\begin{enumerate}
\item \textbf{Identifier la thèse} : l'auteur soutient-il que la religion est une illusion (Freud), une fonction sociale (Durkheim), une idéologie (Marx) ou une expérience sui generis (Otto) ?
\item \textbf{Repérer les arguments} : sont-ils psychologiques, sociologiques, historiques ou logiques ? Sur quels faits s'appuient-ils ?
\item \textbf{Évaluer la portée} : l'explication est-elle réductionniste (ramène-t-elle le religieux à autre chose) ? Prétend-elle tout expliquer ? Tranche-t-elle, oui ou non, la question de l'existence de Dieu ?
\end{enumerate}
\end{methode}

\begin{exoresolu}[Analyse d'une thèse : Freud]
Énoncé. « Les doctrines religieuses sont des illusions, qui ne répondent qu'à nos désirs les plus anciens, les plus urgents. » (Freud). Dégagez la thèse, l'argument et formulez une critique.
\tcblower
\textbf{Solution détaillée.}
1. \emph{Thèse} : la religion est une illusion, c'est-à-dire une croyance dont le contenu est déterminé par la réalisation d'un désir (être protégé, consolé, sauvé de la mort).
2. \emph{Argument} : l'homme est impuissant devant la nature et la souffrance ; comme l'enfant appelle son père, il projette un Dieu paternel tout-puissant. La force affective de la croyance s'explique donc par le besoin, non par la preuve.
3. \emph{Critique} : Freud commet peut-être une faute génétique : montrer qu'une croyance naît d'un désir ne prouve pas qu'elle est fausse. Le désir d'eau ne prouve pas que l'eau n'existe pas. L'argument psychanalytique explique l'origine du sentiment religieux, mais reste silencieux sur la question de l'existence de Dieu.
\end{exoresolu}

\begin{exercice}[Pour s'entraîner]
1. Définissez le \emph{numinosum} et distinguez sacré et profane à partir d'un exemple camerounais.
2. « En adorant Dieu, les hommes adorent la société sans le savoir. » Expliquez et discutez cette thèse durkheimienne.
3. Sujet de dissertation : « La religion n'est-elle qu'une consolation ? »
\end{exercice}

\begin{corrige}[Exercice 1]
1. Le \emph{numinosum} (Otto) est le sentiment du « tout autre », mêlant \emph{tremendum} (crainte) et \emph{fascinans} (attraction). Le sacré est ce qui est mis à part et vénéré (la forêt des ancêtres, l'église), le profane est l'espace ordinaire quotidien.
2. Pour Durkheim, le totem est l'emblème du clan ; la force éprouvée dans les rites vient du groupe rassemblé. Discussion : cette thèse explique bien la cohésion sociale, mais peine à rendre compte de l'expérience mystique individuelle et des prophètes critiques de leur propre société.
3. Plan possible : I. La religion comme consolation (Freud : illusion protectrice ; Marx : opium du peuple face à la misère). II. Mais la religion est plus qu'une consolation : expérience du sacré (Otto), quête de sens, exigence morale et engagement (martyrs, prophètes). III. Synthèse : la consolation est une fonction réelle mais non exhaustive du religieux ; la question de la vérité de la croyance demeure ouverte.
\end{corrige}

\coursec{Les conceptions de Dieu : une cartographie conceptuelle}

Dans la section précédente, nous avons interrogé l'\emph{origine} du sentiment religieux : naît-il de la peur, de l'admiration devant le merveilleux, ou d'une projection sociale ? Mais cette question généalogique laisse entière une autre question, plus conceptuelle : une fois que l'homme croit, \emph{à quoi} croit-il exactement ? Dire « je crois en Dieu » ne suffit pas : le Dieu de Voltaire n'est pas celui de Spinoza, et le Dieu de Spinoza n'est pas celui d'un chrétien de Yaoundé ni celui d'un sage bamiléké. Il faut donc dresser une \cle{cartographie} des grandes conceptions de Dieu, c'est-à-dire classer les réponses possibles à la question : \emph{qu'est-ce que Dieu et quel rapport entretient-il avec le monde ?} Nous verrons successivement le théisme, le déisme, le panthéisme, le panenthéisme, puis les positions négatives ou suspendues (athéisme, agnosticisme), avant d'interroger les conceptions africaines et camerounaises.

\courssub{Le théisme classique : le Dieu personnel des monothéismes}

Commençons par l'intuition la plus familière. Quand un croyant de Douala prie avant un examen, il s'adresse à \emph{quelqu'un} : un être qui l'entend, qui peut l'aider, qui le juge et qui l'aime. C'est l'intuition du \cle{théisme}.

\begin{definition}[Théisme]
Le \cle{théisme} est la conception de Dieu comme un \textbf{être personnel}, à la fois \textbf{transcendant} (radicalement distinct du monde qu'il a créé) et \textbf{immanent} (présent et agissant dans le monde par sa providence). C'est la conception des grands monothéismes : judaïsme, christianisme, islam.
\end{definition}

Le théisme classique attribue à Dieu des \cle{attributs} précis, qu'il faut connaître pour le Bac :

\begin{itemize}
\item \textbf{Omnipotence} : Dieu est tout-puissant. Il crée le monde \emph{ex nihilo}, c'est-à-dire « à partir de rien », sans matière préexistante. Le monde n'est donc pas une émanation de Dieu, mais une \emph{créature} librement voulue.
\item \textbf{Omniscience} : Dieu sait tout, passé, présent et avenir.
\item \textbf{Omnibonté} : Dieu est infiniment bon, source de toute valeur morale.
\item \textbf{Transcendance} : Dieu n'est pas une partie du monde ; il le dépasse infiniment. D'où la distinction fondamentale entre \cle{Créateur} et \cle{créature} : confondre les deux serait, pour le théiste, une idolâtrie.
\item \textbf{Immanence (providence)} : Dieu n'abandonne pas sa création ; il la gouverne par sa \emph{providence}, il répond aux prières, il peut accomplir des miracles.
\end{itemize}

\begin{exemplebox}[La distinction Créateur/créature]
Dans la Genèse, Dieu dit : « Que la lumière soit », et la lumière fut. Le monde dépend totalement de Dieu, mais Dieu ne dépend pas du monde. Image scolaire : le potier et le pot. Le potier existe sans le pot ; le pot n'existe que par le potier. Toute la différence avec le panthéisme (que nous verrons) tient dans cette asymétrie.
\end{exemplebox}

Le problème majeur du théisme, que nous retrouverons dans la section 3, est le \emph{problème du mal} : comment concilier l'omnipotence, l'omniscience et l'omnibonté avec l'existence de la souffrance innocente (un enfant victime d'un accident de la route sur l'axe Yaoundé--Bafoussam) ?

\courssub{Le déisme : l'horloger qui s'est retiré}

Imaginez maintenant un horloger qui fabrique une montre parfaite, la remonte, puis part en voyage pour toujours. La montre fonctionne seule, selon ses lois. C'est exactement l'image que \textbf{Voltaire} et les philosophes des \cle{Lumières} (XVIII\textsuperscript{e} siècle) donnent de Dieu.

\begin{definition}[Déisme]
Le \cle{déisme} affirme l'existence d'un Dieu \textbf{cause première intelligente}, architecte de l'univers, mais \textbf{non interventionniste} : après avoir créé le monde et ses lois, Dieu ne s'en occupe plus. Il n'y a ni révélation, ni miracles, ni providence, ni prière efficace.
\end{definition}

Le déisme repose sur un raisonnement par analogie : de même qu'une montre suppose un horloger, l'ordre admirable de l'univers (la régularité des saisons, la précision des lois physiques) suppose une intelligence ordonnatrice. Mais la raison seule suffit à connaître ce Dieu : c'est la \cle{religion naturelle}, fondée sur l'observation du monde, par opposition à la \cle{religion révélée}, fondée sur des textes sacrés (Bible, Coran) que le déiste rejette comme superstitions.

\begin{exemplebox}[Voltaire et l'horloger]
Dans le \emph{Dictionnaire philosophique}, Voltaire écrit que « l'univers m'embarrasse, et je ne puis songer que cette horloge existe et n'ait point d'horloger ». Le déiste admire donc Dieu comme on admire un ingénieur, mais il ne le \emph{prie} pas : à quoi bon prier un être qui, par définition, ne modifie jamais ses lois ?
\end{exemplebox}

\begin{attention}[Théisme $\neq$ déisme]
Les deux positions affirment un Dieu créateur et transcendant, mais le théiste affirme en plus la \textbf{providence} (Dieu agit dans l'histoire, répond aux prières), tandis que le déiste la nie. Erreur fréquente en dissertation : attribuer la providence à Voltaire. Retenez : théisme = Dieu \emph{présent} ; déisme = Dieu \emph{absent} après la création.
\end{attention}

\courssub{Le panthéisme : \emph{Deus sive Natura}}

Renversons complètement la perspective. Et si Dieu n'était pas \emph{au-delà} du monde, mais \emph{identique} au monde ? C'est la thèse audacieuse de \textbf{Spinoza} (XVII\textsuperscript{e} siècle), héritière sur ce point des \textbf{Stoïciens} de l'Antiquité.

\begin{definition}[Panthéisme]
Le \cle{panthéisme} (du grec \emph{pan}, tout, et \emph{theos}, Dieu) identifie Dieu et la Nature : \emph{Deus sive Natura}, « Dieu ou la Nature » (Spinoza, \emph{Éthique}). Dieu n'est ni une personne, ni une volonté, ni un être transcendant : il est la \textbf{substance unique et infinie} dont tout ce qui existe est un mode. L'immanence est \textbf{totale}.
\end{definition}

Chez Spinoza, tout ce qui arrive arrive par \cle{nécessité géométrique} : comme il suit de la nature du triangle que la somme de ses angles vaut deux angles droits, il suit de la nature de Dieu (c'est-à-dire de la Nature) tout ce qui existe. Il n'y a donc ni miracle, ni finalité, ni libre arbitre divin : prier un tel « Dieu » n'a aucun sens, car il n'entend pas et ne décide rien. Le panthéisme dissout la distinction Créateur/créature : nous sommes des \emph{modes} de Dieu, comme les vagues sont des modes de l'océan.

\begin{propriete}[Le Dieu des philosophes $\neq$ le Dieu de la révélation]
Il faut distinguer le \cle{Dieu des philosophes et des savants} — principe rationnel, ordre de la nature, substance spinoziste, auquel Einstein se référait quand il disait croire « au Dieu de Spinoza qui se révèle dans l'harmonie de tout ce qui existe » — et le \cle{Dieu d'Abraham, d'Isaac et de Jacob}, Dieu personnel de la révélation, que Pascal oppose au Dieu des philosophes dans son \emph{Mémorial} (« Dieu d'Abraham, Dieu d'Isaac, Dieu de Jacob, non des philosophes et des savants »). Cette distinction est cruciale pour le dialogue science/foi : un savant peut être « religieux » au sens spinoziste sans adhérer à aucune religion révélée.
\end{propriete}

\courssub{Le panenthéisme : le monde en Dieu}

Entre le théisme (Dieu au-delà du monde) et le panthéisme (Dieu = le monde), une position intermédiaire tente une synthèse : le \cle{panenthéisme} (littéralement : « tout-en-Dieu »).

\begin{definition}[Panenthéisme]
Le \cle{panenthéisme} affirme que le monde est \textbf{en} Dieu, mais que Dieu \textbf{dépasse} le monde : Dieu contient l'univers sans s'y réduire. On le trouve chez \textbf{Hegel} (l'Esprit se réalise dans l'histoire tout en la dépassant), chez \textbf{Whitehead} et les théologies processuelles (Dieu évolue \emph{avec} le monde, il est affecté par lui).
\end{definition}

L'image classique : l'éponge dans l'océan. L'océan (Dieu) pénètre entièrement l'éponge (le monde), mais l'océan est infiniment plus vaste que l'éponge. Le panenthéisme tente ainsi de concilier \textbf{immanence} (Dieu est présent en toute chose, il « vit » l'histoire du monde) et \textbf{transcendance} (Dieu n'est pas épuisé par le monde).

\begin{attention}[Panthéisme $\neq$ panenthéisme]
Piège classique du Bac. Le \textbf{panthéisme} dit : \emph{Tout est Dieu} (identité pure et simple ; il nie la distinction Créateur/créature). Le \textbf{panenthéisme} dit : \emph{Tout est en Dieu, mais Dieu est plus que le Tout} (il préserve la distinction et une forme de transcendance). Formule mnémotechnique : panthéisme = « Dieu = Monde » ; panenthéisme = « Monde $\subset$ Dieu ».
\end{attention}

\courssub{Athéisme et agnosticisme : la négation et la suspension}

Face à toutes ces conceptions positives, deux positions se définissent par un refus ou une réserve.

\begin{definition}[Athéisme]
L'\cle{athéisme} est la négation de l'existence de Dieu. On distingue l'\textbf{athéisme théorique (métaphysique)}, qui affirme « Dieu n'existe pas » (Marx : la religion est « l'opium du peuple », reflet aliéné de la misère sociale ; Nietzsche : « Dieu est mort » ; Sartre : l'existence de Dieu est incompatible avec la liberté humaine), et l'\textbf{athéisme pratique}, qui consiste simplement à \emph{vivre sans référence à Dieu}, sans même poser la question. On parle aussi d'\textbf{athéisme méthodologique} pour les sciences : le physicien n'a pas besoin de l'hypothèse Dieu pour expliquer la chute des corps, sans pour autant nier Dieu.
\end{definition}

\begin{definition}[Agnosticisme]
L'\cle{agnosticisme} (mot forgé par \textbf{Huxley} en 1869, du grec \emph{a-gnôsis}, « sans connaissance ») est la \textbf{suspension du jugement} : l'existence de Dieu est \emph{inconnaissable}. Pour \textbf{Kant}, la raison théorique ne peut ni prouver ni réfuter Dieu, car Dieu n'est pas un objet d'expérience possible ; pour le positivisme, seul ce qui est vérifiable par l'expérience a du sens.
\end{definition}

\begin{attention}[Agnosticisme $\neq$ scepticisme]
L'agnostique ne dit pas « on ne peut rien savoir en général » (scepticisme) ; il dit précisément : \emph{sur cette question}, la raison humaine est incompétente. Kant sauve d'ailleurs Dieu par la raison \emph{pratique} : il faut le postuler pour fonder la morale. L'agnosticisme n'est donc pas un athéisme timide, mais une thèse sur les \textbf{limites de la connaissance}.
\end{attention}

\courssub{Les conceptions africaines et camerounaises de l'Être Suprême}

Notre cartographie serait ethnocentrique si elle ignorait les traditions religieuses africaines. Chez les peuples du Cameroun et d'Afrique, on trouve presque partout la croyance en un \cle{Être Suprême} créateur : \textbf{Zambe} (ou Nzambe) chez les Fang, \textbf{Nsi} ou \textbf{Nyuy} (Nwob) chez les Bamiléké, \textbf{Roog} chez les Sérères du Sénégal, \textbf{Mawu} chez les Fons du Bénin. Mais cet Être Suprême présente souvent un trait étonnant pour un esprit formé au théisme : il s'est \emph{retiré}.

\begin{definition}[\emph{Deus otiosus}]
Le \cle{deus otiosus} (« Dieu oisif » ou « Dieu retiré ») est un concept anthropologique (W. Schmidt, Bolaji Idowu) désignant un \textbf{Haut Dieu créateur qui, après avoir créé le monde, s'en est éloigné}, laissant la gestion quotidienne aux esprits, aux divinités secondaires et aux ancêtres. Cette figure est fréquente dans les religions traditionnelles africaines : on n'offre que rarement des sacrifices directement à l'Être Suprême ; on s'adresse aux médiateurs.
\end{definition}

\begin{exemplebox}[La médiation des ancêtres]
Dans bien des traditions camerounaises, un chef de famille qui bénit un enfant ou scelle une réconciliation s'adresse aux \textbf{ancêtres} du lignage, non directement à l'Être Suprême. L'ancêtre, mort mais « vivant » dans l'invisible, est le médiateur naturel : proche des vivants, proche du divin. La religion vécue est donc moins un « monothéisme » pur qu'un système \emph{hiérarchisé} : un Dieu lointain au sommet, des esprits et des ancêtres actifs au quotidien.
\end{exemplebox}

Cette réalité pose un problème d'interprétation. Les premiers missionnaires, voulant établir une continuité avec le christianisme, ont présenté ces traditions comme des \textbf{monothéismes} imparfaits, traduisant les noms locaux par « Dieu ». Or certains anthropologues objectent que ces systèmes sont plutôt \textbf{polythéistes} (plusieurs dieux) ou \textbf{hénothéistes} (plusieurs dieux, mais un seul est adoré comme suprême par un groupe donné). La leçon philosophique est importante : les catégories occidentales (théisme, polythéisme) ne s'appliquent pas toujours sans déformation aux expériences religieuses africaines.

\courssub{Cartographie synthétique et méthode de comparaison}

\begin{popfigure}
\begin{center}
\begin{tikzpicture}[
  grow=down,
  level distance=1.5cm,
  level 1/.style={sibling distance=4.6cm},
  level 2/.style={sibling distance=2.6cm, level distance=1.6cm},
  every node/.style={draw, rounded corners=2pt, align=center, font=\small, inner sep=3pt},
  edge from parent/.style={draw=popmuted, -latex}
]
\node[fill=popblue, text=white, font=\small\bfseries, align=center]{Question : quel rapport\\ Dieu / Monde ?}
  child { node[fill=popblueL, align=center]{\textbf{Transcendance}\\ Dieu au-delà du monde}
    child { node[fill=popcream, align=center]{\textbf{Théisme}\\ Personnel,\\ interventionniste\\ (providence)} }
    child { node[fill=popcream, align=center]{\textbf{Déisme}\\ Intelligent,\\ non interventionniste\\ (horloger)} }
  }
  child { node[fill=popgreenL, align=center]{\textbf{Immanence}\\ Dieu dans le monde}
    child { node[fill=popcream, align=center]{\textbf{Panthéisme}\\ Dieu = Monde\\ impersonnel,\\ nécessité} }
    child { node[fill=popcream, align=center]{\textbf{Panenthéisme}\\ Monde en Dieu,\\ Dieu $>$ Monde} }
  }
  child { node[fill=poporangeL, align=center]{\textbf{Négation}\\ Pas de Dieu}
    child { node[fill=popcream, align=center]{\textbf{Athéisme}\\ théorique\\ (Marx, Nietzsche)} }
    child { node[fill=popcream, align=center]{\textbf{Athéisme}\\ pratique ou\\ méthodologique} }
  }
  child { node[fill=poppurpleL, align=center]{\textbf{Suspension}\\ Inconnaissable}
    child { node[fill=popcream, align=center]{\textbf{Agnosticisme}\\ (Kant, Huxley)\\ limites de la raison} }
  };
\end{tikzpicture}
\end{center}
\legende{Taxinomie des grandes conceptions de Dieu. Premier critère de classement : le rapport de Dieu au monde (transcendance, immanence, négation, suspension). Second critère : les attributs clés (personnel ou impersonnel ; interventionniste ou non).}
\end{popfigure}

\begin{methode}[Comparer deux conceptions de Dieu en 5 questions]
Pour comparer rigoureusement deux conceptions (exercice fréquent en dissertation), interrogez successivement :
\begin{enumerate}
\item \textbf{Statut de la transcendance} : Dieu est-il au-delà du monde, identique au monde, ou le monde est-il en Dieu ?
\item \textbf{Personnalité et intelligence} : Dieu est-il une personne qui veut et qui aime, ou un principe impersonnel ?
\item \textbf{Rapport au mal et à la souffrance} : le mal est-il un scandale (théisme), une nécessité (panthéisme), un fait sans sens (athéisme) ?
\item \textbf{Fondement de la morale} : le bien vient-il de la volonté divine, de la nature des choses, ou de l'homme seul ?
\item \textbf{Possibilité de la prière et du culte} : peut-on s'adresser à Dieu et espérer une réponse ?
\end{enumerate}
Appliquez cette grille à deux conceptions et vous obtenez un plan de comparaison solide.
\end{methode}

\begin{exemplebox}[Un fait religieux camerounais]
Selon les enquêtes Afrobarometer menées au Cameroun, l'immense majorité des Camerounais (plus de neuf sur dix) déclare que la religion est importante dans leur vie. Mais cette adhésion massive recouvre une grande diversité de représentations : pour les uns, Dieu est une \emph{personne} qui écoute et intervient (conception théiste) ; pour d'autres, une \emph{force} ou une \emph{puissance} impersonnelle qui ordonne le cosmos (conception proche du panthéisme ou du \emph{deus otiosus}). La statistique sociologique ne tranche donc pas la question philosophique : \emph{quelle} conception de Dieu est cohérente ?
\end{exemplebox}

\begin{savaistu}[Le Pari de Pascal]
Le célèbre « Pari » de Pascal (\emph{Pensées}) ne prétend pas \emph{prouver} l'existence de Dieu : il montre la \textbf{rationalité du pari} sur son existence. Si je parie que Dieu existe et qu'il existe, je gagne l'infini (le salut) ; s'il n'existe pas, je perds peu. Si je parie contre et que je me trompe, je perds l'infini. C'est donc un argument \emph{pragmatique} (une décision dans l'incertitude), non une démonstration \emph{théorique}. Ne le citez jamais comme une preuve de l'existence de Dieu : ce serait une erreur d'analyse sanctionnée au Bac.
\end{savaistu}

\begin{exoresolu}[Identifier une conception de Dieu]
\textbf{Énoncé.} Un élève déclare : « Je crois qu'une intelligence a organisé l'univers, car tout y est réglé comme une mécanique admirable ; mais il est absurde de prier : les lois de la nature ne changent pas pour nos demandes. » À quelle conception de Dieu cette position correspond-elle ? Justifiez en mobilisant deux attributs.
\tcblower
\textbf{Solution détaillée.} Il s'agit du \textbf{déisme}. Deux indices le montrent. (1) L'élève affirme une \emph{cause première intelligente} : l'ordre de l'univers suppose un ordonnateur — c'est l'argument de l'horloger cher à Voltaire. (2) Il nie la \emph{providence} : Dieu n'intervient pas dans le monde, la prière est inefficace, il n'y a ni miracle ni révélation. Ce n'est donc pas un théisme (qui affirme la providence et la possibilité de la prière), ni un panthéisme (le Dieu de l'élève est distinct du monde et intelligent, alors que le Dieu de Spinoza est impersonnel et identique à la Nature), ni un athéisme (il affirme l'existence de Dieu). Conclusion : déisme = Dieu architecte transcendant mais absent.
\end{exoresolu}

\begin{aretenir}[L'essentiel de la section]
\begin{itemize}
\item \textbf{Théisme} : Dieu personnel, transcendant (création \emph{ex nihilo}) et provident ; distinction Créateur/créature.
\item \textbf{Déisme} : Dieu horloger, cause première, mais non interventionniste ; religion naturelle contre religion révélée.
\item \textbf{Panthéisme} : \emph{Deus sive Natura} (Spinoza) ; immanence totale, nécessité, pas de personne.
\item \textbf{Panenthéisme} : le monde est en Dieu, mais Dieu dépasse le monde.
\item \textbf{Athéisme} : négation (théorique ou pratique) ; \textbf{agnosticisme} : suspension du jugement (Kant, Huxley).
\item \textbf{Traditions africaines} : Être Suprême souvent \emph{deus otiosus}, médiation des ancêtres ; prudence dans l'usage des catégories occidentales.
\end{itemize}
\end{aretenir}

\begin{exercice}[Entraînement au Bac]
1. Définissez le panenthéisme et distinguez-le soigneusement du panthéisme (4 points).\\
2. « Le Dieu des philosophes n'est pas le Dieu des croyants. » Expliquez cette formule en vous appuyant sur Spinoza et Pascal (6 points).\\
3. Dissertation : « Peut-on prier le Dieu de Spinoza ? » (Vous construirez un plan dialectique : thèse / antithèse / synthèse.)
\end{exercice}

\begin{corrige}[Exercice]
\textbf{1.} Le panenthéisme affirme que le monde est \emph{en} Dieu mais que Dieu \emph{dépasse} le monde (Hegel, Whitehead) : il combine immanence et transcendance. Le panthéisme, lui, \emph{identifie} purement Dieu et le monde (Spinoza) : il supprime toute transcendance et la distinction Créateur/créature, que le panenthéisme conserve.\\
\textbf{2.} Le Dieu des philosophes est un principe rationnel déduit de l'ordre du monde : substance impersonnelle chez Spinoza (\emph{Deus sive Natura}), harmonie intelligible chez Einstein. Le Dieu des croyants est le Dieu personnel de la révélation, « Dieu d'Abraham, d'Isaac et de Jacob » (Pascal, \emph{Mémorial}), qui entre en relation avec l'homme. Le premier se connaît par la raison, le second par la foi.\\
\textbf{3.} Pistes : \emph{Thèse} — non, car le Dieu de Spinoza est impersonnel, sans volonté ni audition ; tout est nécessité, prier est une superstition née de l'ignorance des causes. \emph{Antithèse} — mais Spinoza propose un équivalent rationnel : l'\emph{amor Dei intellectualis}, l'amour intellectuel de Dieu, joie de comprendre la nécessité ; c'est une forme de « piété » sans demande. \emph{Synthèse} — on ne peut pas \emph{supplier} le Dieu de Spinoza, mais on peut l'\emph{aimer} en le connaissant : la prière de demande disparaît au profit de la contemplation rationnelle.
\end{corrige}

\coursec{L'existence de Dieu : les preuves classiques et leurs critiques}

Après avoir cartographié les conceptions de Dieu — personnel ou impersonnel, transcendant ou immanent, créateur ou principe d'ordre — nous abordons maintenant la question décisive : \textbf{peut-on démontrer rationnellement que Dieu existe ?} Depuis l'Antiquité, les philosophes ont tenté de construire des preuves \emph{a priori} (à partir du concept de Dieu) ou \emph{a posteriori} (à partir de l'expérience du monde). Cette section examine les quatre grandes preuves classiques, leurs mécanismes logiques, et les critiques qui les ont ébranlées, sans trancher la question de la foi, qui relève d'un autre registre.

\courssub{La preuve ontologique : de l'idée à l'existence}

\paragraph{Intuition et énoncé.} L'idée centrale, formulée par \textbf{Anselme de Cantorbéry} (XI\textsuperscript{e} s.), est saisissante : Dieu est défini comme \emph{« celui dont on ne peut rien concevoir de plus grand »}. Or, un être qui n'existerait que dans l'entendement (dans l'idée) serait moins grand qu'un être qui existerait aussi dans la réalité. Donc, Dieu \emph{doit} exister dans la réalité, sinon il ne serait pas « celui dont on ne peut rien concevoir de plus grand ».

\begin{definition}[Preuve ontologique]
Démonstration \emph{a priori} qui déduit l'existence réelle de Dieu de son seul concept ou définition. Elle prétend montrer que l'existence fait partie de l'essence de Dieu (existence nécessaire), contrairement aux êtres contingents dont l'existence est accidentelle.
\end{definition}

\paragraph{Les grandes versions.}
\begin{itemize}
    \item \textbf{Descartes} (\emph{Méditations métaphysiques}, V) : L'idée de Dieu (être souverainement parfait) est innée en moi. La perfection implique l'existence (un être parfait qui n'existerait pas manquerait de perfection). Donc Dieu existe. L'existence est inséparable de l'essence divine comme la somme des angles de 180° est inséparable de l'essence du triangle.
    \item \textbf{Leibniz} : La preuve d'Anselme/Descartes est valide \emph{à condition} que le concept de Dieu soit \emph{possible} (non contradictoire). Il tente de démontrer cette possibilité : les perfections sont simples, positives, illimitées et compatibles entre elles.
    \item \textbf{Kurt Gödel} (1970) : Formalisation en logique modale (système S5). Si la propriété « être divin » est positive et possible, alors Dieu existe nécessairement dans tous les mondes possibles. Vérifié par ordinateur (2013), mais la validité dépend de l'acceptation des axiomes modaux.
\end{itemize}

\begin{savaistu}[Gödel et la machine]
En 2013, les logiciens Christoph Benzmüller et Bruno Woltzenlogel Paleo ont encodé la preuve modale de Gödel dans les assistants de preuve \emph{Coq} et \emph{Isabelle}. L'ordinateur a validé la déduction logique : \emph{si} les axiomes sont acceptés, la conclusion suit nécessairement. Cela ne prouve pas Dieu, mais la rigueur formelle de l'argument.
\end{savaistu}

\paragraph{La critique kantienne décisive.}
\begin{propriete}[Théorème (Kant) : L'existence n'est pas un prédicat réel]
Dans la \emph{Critique de la raison pure} (Dialectique transcendantale), Kant montre que « l'existence » n'ajoute \emph{aucun contenu} au concept d'une chose. Cent thalers réels ne contiennent pas un centime de plus que cent thalers possibles. Le concept reste identique ; seule la \emph{position} de l'objet change (il est donné dans l'intuition). On ne peut donc pas passer du plan \emph{logique} (concept) au plan \emph{réel} (existence) par une simple analyse conceptuelle. La preuve ontologique commet un \emph{sophisme} : elle traite l'existence comme une propriété (prédicat) alors qu'elle est la condition de toute propriété.
\end{propriete}

\begin{attention}[Piège fréquent]
Dire « Kant a prouvé que Dieu n'existe pas » est faux. Kant a prouvé que la \emph{raison théorique} ne peut \emph{prouver} l'existence de Dieu. Il laisse la place à la foi rationnelle (postulat de la raison pratique, voir 3.4).
\end{attention}

\courssub{La preuve cosmologique : la chaîne des causes}

\paragraph{Intuition et énoncé.} Tout ce qui commence à exister a une cause. L'univers a commencé à exister (ou est contingent). Donc l'univers a une cause : Dieu. C'est l'argument de la \emph{cause première} ou de l'\emph{être nécessaire}.

\begin{definition}[Être nécessaire vs Être contingent]
Un \cle{être nécessaire} existe par sa propre essence ; il est impossible qu'il n'existe pas (sa non-existence implique contradiction). Un \cle{être contingent} tire son existence d'un autre ; il peut ne pas exister (sa non-existence est concevable sans contradiction). La preuve cosmologique leibnizienne repose sur le \emph{Principe de raison suffisante} : rien n'est sans raison ; la raison totale de l'univers contingent doit se trouver dans un être nécessaire.
\end{definition}

\paragraph{Les grandes versions.}
\begin{itemize}
    \item \textbf{Aristote / Thomas d'Aquin (2\textsuperscript{e} voie)} : Dans le monde sensible, rien n'est cause efficiente de soi-même. La régression à l'infini des causes efficientes est impossible (pas de cause première = pas d'effet intermédiaire). Donc il faut une \emph{cause première non causée} : Dieu.
    \item \textbf{Leibniz} : Pourquoi y a-t-il quelque chose plutôt que rien ? La raison suffisante de l'univers contingent ne peut être trouvée dans la série des choses contingentes (régression infinie). Elle doit être dans un \emph{être nécessaire} qui porte en lui sa propre raison d'être.
    \item \textbf{Argument Kalam (théologie islamique, W.L. Craig)} : 1. Tout ce qui a un commencement a une cause. 2. L'univers a un commencement (Big Bang). 3. Donc l'univers a une cause (personnelle, intemporelle, immatérielle).
\end{itemize}

\begin{popfigure}
\begin{tikzpicture}[
    node distance=1.2cm and 1.5cm,
    box/.style={draw=popblue, thick, rounded corners, fill=popblue!10, minimum width=2.5cm, minimum height=1cm, align=center},
    arrow/.style={-Stealth, thick, popdark}
]
\node[box, align=center] (contingent){Êtres contingents\\ (Univers, nous)};
\node[box, below=of contingent, align=center] (regression){Régression des causes\\ (Infinie ?)};
\node[box, below=of regression, fill=poporange!15, draw=poporange, align=center] (necessaire){Être Nécessaire\\ (Cause non causée)};
\node[box, right=of necessaire, fill=popgreen!15, draw=popgreen] (dieu) {Dieu};

\draw[arrow] (contingent) -- (regression) node[midway, right] {dépend de};
\draw[arrow] (regression) -- (necessaire) node[midway, right] {exige};
\draw[arrow, poporange] (necessaire) -- (dieu) node[midway, above] {identifié comme};
\end{tikzpicture}
\legende{Schéma de la preuve cosmologique leibnizienne : de la contingence du monde à la nécessité d'un fondement absolu.}
\end{popfigure}

\paragraph{Critiques majeures.}
\begin{itemize}
    \item \textbf{Hume} (\emph{Dialogues sur la religion naturelle}) : 1. L'idée de cause nécessaire n'a pas de sens empirique ; nous ne voyons que des \emph{conjonctions constantes}. 2. Pourquoi l'univers ne serait-il pas sa propre cause ou une « brute fact » (fait brut) ? 3. La régression à l'infini n'est pas logiquement impossible (série infinie de nombres).
    \item \textbf{Russell} : « L'univers est juste là, et c'est tout » (\emph{brute fact}). Le principe de causalité s'applique \emph{dans} l'univers, pas \emph{à} l'univers comme totalité.
    \item \textbf{Physique contemporaine} : Le Big Bang marque un début \emph{temporel} (t=0), mais les modèles quantiques (Hartle-Hawking, Vilenkin, fluctuations du vide) proposent des descriptions sans singularité initiale ni cause externe. La causalité classique s'effondre à l'échelle de Planck.
\end{itemize}

\begin{exemplebox}[Big Bang vs Cause divine]
L'argument Kalam utilise le Big Bang comme prémisse empirique (« l'univers a commencé »). Mais la cosmologie quantique montre que « début temporel » $\neq$ « création ex nihilo par un agent ». Le vide quantique n'est pas le « rien » philosophique ; il a des lois, de l'énergie, une structure. La physique décrit le \emph{comment}, pas le \emph{pourquoi} métaphysique.
\end{exemplebox}

\courssub{La preuve physico-théologique (finaliste) : l'horloger et le fine-tuning}

\paragraph{Intuition et énoncé.} La nature manifeste un ordre, une complexité, une finalité (télos) qui ne peuvent être le fruit du hasard. Comme une montre implique un horloger, l'univers implique une Intelligence ordonnatrice.

\paragraph{Versions classiques et modernes.}
\begin{itemize}
    \item \textbf{Thomas d'Aquin (5\textsuperscript{e} voie)} : Les corps naturels agissent pour une fin (régularité des astres, des corps). Ils n'ont pas d'intelligence. Donc ils sont dirigés par un Être intelligent : Dieu.
    \item \textbf{William Paley (1802)} : Analogie de la montre trouvée dans la lande. Complexité adaptée $\rightarrow$ Dessein $\rightarrow$ Desseinateur.
    \item \textbf{\emph{Intelligent Design} (Behe, Dembski)} : « Complexité irréductible » (ex. flagelle bactérien, cascade de coagulation). Impossible par étapes darwiniennes successives.
    \item \textbf{Fine-tuning (Ajustement fin)} : Les constantes physiques (constante de structure fine $\alpha$, constante cosmologique $\Lambda$, masse du proton, force nucléaire forte) sont « ajustées » avec une précision extrême pour permettre la vie (chimie du carbone, stabilité des étoiles).
\end{itemize}

\begin{exemplebox}[Chiffres du Fine-tuning]
La \cle{constante de structure fine} $\alpha \approx 1/137,036$.
\begin{itemize}
    \item Si $\alpha$ variait de $\pm 4\%$ : la résonance de Hoyle dans le carbone-12 disparaît $\rightarrow$ pas de synthèse du carbone dans les étoiles $\rightarrow$ pas de chimie organique $\rightarrow$ pas de vie.
    \item Constante cosmologique $\Lambda$ : précision $\sim 10^{-120}$ (1 chance sur $10^{120}$).
\end{itemize}
Argument théiste (Polkinghorne, Collins) : Best explanation = Dessein. Argument naturaliste (Susskind, Weinberg) : \emph{Multivers} (paysage des vacua de la théorie des cordes) + principe anthropique (nous ne pouvons observer que l'univers compatible avec nous).
\end{exemplebox}

\paragraph{Critiques.}
\begin{itemize}
    \item \textbf{Darwin / Néo-darwinisme} : La sélection naturelle + mutations aléatoires + temps profond expliquent la complexité adaptée \emph{sans} dessein. La « complexité irréductible » est réfutée par l'exaptation (cooption de structures existantes pour de nouvelles fonctions).
    \item \textbf{Hume} (anticipant Darwin) : 1. Analogie faible (univers $\neq$ artefact). 2. Mal/imperfection dans la nature (parasites, extinction, souffrances) $\rightarrow$ Desseinateur incompétent ou cruel. 3. Pluralité de dieux possibles (équipe de concepteurs). 4. L'ordre pourrait être immanent (matière auto-organisatrice).
    \item \textbf{Épistémologie : « Dieu des interstices » (God of the gaps)}. Invoquer Dieu pour combler les lacunes scientifiques actuelles est une stratégie perdante : l'histoire des sciences comble ces lacunes (foudre $\rightarrow$ électricité ; mouvement planétaire $\rightarrow$ gravitation ; vie $\rightarrow$ évolution).
\end{itemize}

\begin{attention}[Erreur de catégorie]
Dire « La science a réfuté Dieu » ou « Le Fine-tuning prouve Dieu » commet la même erreur : \emph{confusion de niveaux} (category mistake). La science explique les \emph{mécanismes immanents} (comment la structure émerge). La métaphysique questionne le \emph{fondement transcendant} (pourquoi il y a des lois, pourquoi ces lois permettent la vie). Les deux discours sont orthogonaux, pas concurrents sur le même plan.
\end{attention}

\courssub{La preuve morale : postulat de la raison pratique}

\paragraph{Intuition et énoncé.} Kant (\emph{Critique de la raison pratique}) renonce à prouver Dieu \emph{théoriquement}. Mais la \emph{raison pratique} (morale) exige le \emph{Souverain Bien} : la conjonction de la Vertu (faire son devoir) et du Bonheur (proportionné à la vertu). Or, dans l'expérience, le vertueux n'est pas toujours heureux, le méchant prospère. Pour que l'exigence morale ne soit pas illusoire, il faut \emph{postuler} : 1. L'immortalité de l'âme (temps infini pour progresser vers la sainteté). 2. L'existence de Dieu (juge suprême qui assure la distribution du bonheur proportionné à la vertu). 3. La liberté (condition de la responsabilité).

\begin{propriete}[Preuve morale (Kant)]
Ce n'est pas une preuve \emph{théorique} (connaissance), mais un \cle{postulat} \emph{pratique} (foi rationnelle). « Il est moralement nécessaire d'assumer l'existence de Dieu. » Newman ajoutera le « sens moral » (conscience) comme voix de Dieu.
\end{propriete}

\paragraph{Critique.} Elle ne convainc que celui qui partage l'exigence morale absolue. Elle ne prouve pas le Dieu des philosophes (créateur, tout-puissant), mais un « Dieu moral » (juge, garant). Elle laisse intact le problème du mal (voir 3.5).

\courssub{Le problème du mal : l'épreuve de la théodicée}

\paragraph{Formulation logique (Épicure / Mackie).}
\begin{quote}
« Dieu veut-il empêcher le mal, mais ne le peut pas ? Il est impuissant.
Peut-il, mais ne le veut pas ? Il est méchant.
Veut-il et le peut ? D'où vient le mal ?
Peut ni ne veut ? Pourquoi l'appeler Dieu ? »
\end{quote}
\cle{Mackie} (1955) : L'ensemble \{Dieu tout-puissant, Dieu tout-bon, Mal existe\} est \emph{logiquement incohérent}.

\paragraph{Formulation probabiliste (Rowe).} L'existence de maux « gratuits », « intenses », « apparemment sans finalité rédemptrice » (ex. faon brûlé vif dans un incendie de forêt, souffrances d'enfants) rend l'existence de Dieu \emph{improbable}, même si non logiquement impossible.

\begin{popfigure}
\begin{tikzpicture}[
    node distance=1.5cm and 2cm,
    attr/.style={draw=popblue, thick, ellipse, fill=popblue!15, minimum width=2.2cm, minimum height=1.2cm, align=center},
    evil/.style={draw=poporange, thick, rectangle, rounded corners, fill=poporange!15, minimum width=2.5cm, minimum height=1.2cm, align=center},
    sol/.style={draw=popgreen, thick, diamond, fill=popgreen!15, aspect=2, align=center},
    arrow/.style={-Stealth, thick, popdark},
    solarrow/.style={-Stealth, thick, popgreen}
]
\node[attr] (omnip) {Omnipotence};
\node[attr, right=of omnip] (omnib) {Omnibonté};
\node[attr, below=of omnip] (omnis) {Omniscience};
\node[evil, below=of omnib, align=center] (mal){Réalité du Mal\\ (Souffrance, Injustice)};

\draw[arrow] (omnip) -- (mal);
\draw[arrow] (omnib) -- (mal);
\draw[arrow] (omnis) -- (mal);

\node[sol, left=of omnip, xshift=-1cm, align=center] (libre){Libre Arbitre\\(Plantinga)};
\node[sol, right=of omnib, xshift=1cm, align=center] (ame){Âme-Valeur\\(Hick)};
\node[sol, below=of mal, align=center] (privatio){Mal = Privation\\ (Augustin/Leibniz)};
\node[sol, below=of privatio, align=center] (inconnu){Inconnaissance\\ Humaine (Skeptical Theism)};

\draw[solarrow] (libre) -- (mal);
\draw[solarrow] (ame) -- (mal);
\draw[solarrow] (privatio) -- (mal);
\draw[solarrow] (inconnu) -- (mal);
\end{tikzpicture}
\legende{Le triangle du problème du mal : les attributs divins classiques face à la réalité du mal, et les principales théodicées (réponses).}
\end{popfigure}

\paragraph{Les grandes théodicées (justifications de Dieu).}
\begin{itemize}
    \item \textbf{Libre arbitre (Plantinga, Défense du Libre Arbitre)} : Le mal moral (guerre, crime) est le prix nécessaire de la liberté authentique. Un monde avec des créatures libres (capables d'aimer) est meilleur qu'un monde d'automates, même s'ils choisissent parfois le mal. \emph{Défense logique} : Dieu ne peut pas créer des agents libres \emph{et} déterminer qu'ils fassent toujours le bien (contradiction). $\rightarrow$ Résout l'incohérence \emph{logique} (Mackie).
    \item \textbf{Âme-valeur / Soul-making (Hick, Irénée)} : Le mal naturel (séismes, maladies) est nécessaire pour le développement moral et spirituel (courage, compassion, solidarité). Ce monde est une « vallée de fabrication d'âmes », pas un paradis confortable.
    \item \textbf{Mal comme privation (Augustin, Leibniz)} : Le mal n'est pas une « chose » créée par Dieu, mais une \emph{privation} de bien (aveuglement = privation de vue, injustice = privation de justice). Dieu crée l'être (bon) ; le mal est parasite du bien.
    \item \textbf{Théisme sceptique (Bergmann, Howard-Snyder)} : Nos facultés cognitives sont limitées. Nous ne sommes pas en position de juger qu'un mal est « gratuit » (sans raison divine). L'inférence « Je ne vois pas de raison $\rightarrow$ Il n'y a pas de raison » est invalide.
\end{itemize}

\begin{methode}[Évaluer une preuve de l'existence de Dieu]
\begin{enumerate}
    \item \textbf{Formaliser} le raisonnement sous forme de syllogisme (Prémisse 1, Prémisse 2, Conclusion).
    \item \textbf{Vérifier la validité logique} : La conclusion suit-elle nécessairement des prémisses ? (Modus ponens, modus tollens, etc.).
    \item \textbf{Examiner la vérité des prémisses} : Sont-elles empiriques (vérifiables) ? Métaphysiques (intuitives, axiomatiques) ? Théologiques (révélées) ?
    \item \textbf{Identifier le type de conclusion} : Dieu des philosophes (Cause première, Horloger, Être nécessaire) ou Dieu de la foi (Personnel, Amour, Trinité) ?
    \item \textbf{Tester les contre-arguments} : Y a-t-il des explications alternatives naturalistes ? Des objections logiques (Kant, Hume) ? Des faits récalcitrants (Mal, Fine-tuning vs Multivers) ?
\end{enumerate}
\end{methode}

\begin{exoresolu}[Analyse de l'argument du Fine-tuning]
\textbf{Énoncé :} « Les constantes physiques sont ajustées pour la vie avec une précision extrême (ex. $\alpha \approx 1/137,036$). Cet ajustement ne peut être dû au hasard. Donc, un Ajusteur intelligent (Dieu) existe. »

\textbf{Solution détaillée :}
\begin{enumerate}
    \item \textbf{Formalisation :}
    \begin{itemize}
        \item P1 : L'univers permet la vie (observation).
        \item P2 : La permission de la vie dépend de constantes dans un intervalle extrême (fait physique).
        \item P3 : Soit Hasard, Soit Nécessité physique, Soit Dessein (Multivers exclu ou inclus dans Hasard/Nécessité).
        \item P4 : Hasard $\approx 0$ (probabilité $< 10^{-100}$). Nécessité physique non prouvée (pas de Théorie du Tout).
        \item C : Donc Dessein (Dieu) est la meilleure explication (Inférence à la meilleure explication / Abduction).
    \end{itemize}
    \item \textbf{Validité :} Valide \emph{si} l'alternative est exhaustive (Hasard / Nécessité / Dessein) et si P4 est vrai.
    \item \textbf{Critique des prémisses :}
    \begin{itemize}
        \item P3 : Faux dilemme ? Le \emph{Multivers} (paysage de vacua de la théorie des cordes, inflation éternelle) rend P4 caduc : si $10^{500}$ univers existent avec constantes variables, l'observation de $\alpha \approx 1/137,036$ dans \emph{notre} univers est biaisée par le principe anthropique (sélection d'observation). Ce n'est plus du « hasard » au sens improbable, mais une certitude statistique dans l'ensemble.
        \item P4 : « Nécessité physique » : une future théorie unifiée (M-théorie ?) pourrait déduire $\alpha$ de manière unique.
    \end{itemize}
    \item \textbf{Portée :} Même si l'argument tient, il prouve un \emph{Desseinateur} (Ingénieur cosmique), pas nécessairement le Dieu biblique (Amour, Providence, Révélation). C'est un « Dieu des physiciens », pas le « Dieu d'Abraham, d'Isaac et de Jacob » (Pascal).
\end{enumerate}
\end{exoresolu}

\begin{aretenir}[Synthèse : Les preuves et leurs limites]
\begin{itemize}
    \item \textbf{Ontologique} : Échec théorique (Kant : existence $\neq$ prédicat). Valeur heuristique : clarifie le concept de « nécessaire ».
    \item \textbf{Cosmologique} : Force intuitive (Principe de raison suffisante). Faiblesse : application du principe de causalité \emph{hors} de son domaine d'expérience (univers total), physique quantique (vide, fluctuations).
    \item \textbf{Physico-théologique} : Historique (Paley) réfuté par Darwin. Version moderne (Fine-tuning) : argument abductif fort, mais contourné par Multivers (métaphysique non testable) ou Nécessité future.
    \item \textbf{Morale} : Ne prouve pas, \emph{postule} (Kant). Lie éthique et espérance.
    \item \textbf{Mal} : Objection la plus forte contre le Dieu théiste (Omnipotence + Bonté). Théodicées (Libre arbitre, Âme-valeur) défendent la \emph{cohérence logique}, mais peinent sur le mal \emph{probabiliste} (horreurs gratuites).
\end{itemize}
\emph{Conclusion provisoire :} La raison \emph{ne peut ni prouver, ni réfuter} définitivement l'existence de Dieu. Elle montre la \emph{raisonnabilité} de l'hypothèse théiste (meilleure explication de l'ordre, de la conscience, de la morale) tout en laissant l'espace pour le choix existentiel (Foi / Athéisme / Agnosticisme).
\end{aretenir}

\begin{exercice}[Sujet de dissertation type Bac Tech]
\textbf{Sujet :} « Les preuves de l'existence de Dieu sont-elles encore pertinentes aujourd'hui ? »
\textbf{Consignes :} 1. Analyser les termes. 2. Construire un plan dialectique (Thèse / Antithèse / Synthèse). 3. Mobiliser au moins deux preuves classiques et leurs critiques (Kant, Hume, Darwin, Fine-tuning). 4. Conclure sur le statut de la raison face au religieux.
\end{exercice}

\begin{corrige}[Exercice n°1 - Éléments de corrigé]
\textbf{Introduction :} Accroche (Pascal : « Dieu d'Abraham... non des philosophes »). Définition : Preuve = démonstration rationnelle. Pertinence = capacité à convaincre l'esprit moderne (science, critique). Problématique : La raison théorique peut-elle atteindre l'Absolu ou la foi relève-t-elle d'un autre registre ?
\textbf{Thèse (Pertinence maintenue) :} Preuve cosmologique : Big Bang $\rightarrow$ Cause première (Kalam). Preuve physico-théologique : Fine-tuning $\rightarrow$ Dessein (meilleure explication vs Multivers spéculatif). Preuve morale : Exigence de justice $\rightarrow$ Postulat nécessaire (Kant).
\textbf{Antithèse (Obsolescence / Insuffisance) :} Critique kantienne : Raison théorique incompétente (noumène inconnaissable). Critique humienne : Causalité, analogie, mal. Science : Darwin (complexité sans dessein), Cosmologie quantique (univers sans cause). Preuves = « Dieu des interstices » (recul constant).
\textbf{Synthèse :} Distinction des registres. Preuves = \emph{chemins} (chemins de raison), pas \emph{démonstrations} mathématiques. Elles rendent la foi \emph{raisonnable} (non irrationnelle), mais ne la contraignent pas. Le « pari » (Pascal) ou l'« option fondamentale » (Ricœur) prennent le relais. La pertinence est \emph{existentielle}, pas seulement logique.
\textbf{Conclusion :} Ouverture : Laïcité (État neutre) garantit la liberté de ce choix rationnel et existentiel au Cameroun.
\end{corrige}

\coursec{La religion : définition, formes et valeur sociale}

Après avoir examiné les arguments rationnels pour ou contre l'existence de Dieu (\emph{section 3}), nous devons maintenant déplacer le regard de la \emph{théologie} vers la \emph{sociologie} et l'\emph{anthropologie}. Dieu existe-t-il pour le croyant ? Oui, mais surtout \emph{comment} cette croyance s'organise-t-elle collectivement ? La religion n'est pas seulement une opinion sur le divin ; c'est un \cle{fait social total} (Marcel Mauss) qui structure le temps, l'espace, la morale et le lien social. Au Cameroun, où le Préambule de la Constitution de 1996 proclame la laïcité tout en garantissant la liberté de culte, comprendre la religion comme phénomène social est une nécessité civique.

\courssub{Définir la religion : le face-à-face fonctionnaliste vs substantialiste}

Comment définir un phénomène aussi divers que le christianisme, l'islam, le bouddhisme, les traditions africaines ou le shintoïsme ? Deux approches majeures s'affrontent.

\begin{definition}[Définition fonctionnaliste (Émile Durkheim, \emph{Les Formes élémentaires de la vie religieuse}, 1912)]
« Une religion est un \cle{système solidaire de croyances et de pratiques relatives à des choses sacrées, c'est-à-dire séparées, interdites, qui unissent en une même communauté morale appelée \textbf{Église}, tous ceux qui y adhèrent. »}
\end{definition}

\noindent Pour Durkheim, l'essentiel n'est pas \emph{ce} qu'on croit (le contenu), mais \emph{ce que ça fait} (la fonction). Trois piliers indissociables :
\begin{itemize}
    \item \textbf{Les Croyances :} Représentations collectives qui classent le monde en \cle{sacré} (interdit, pur, puissant) et \cle{profane} (quotidien, impur, banal). Le sacré ne se touche pas sans rites.
    \item \textbf{Les Rites :} Règles d'action déterminées qui régissent les conduites vis-à-vis du sacré (rites positifs : sacrifices, prières ; rites négatifs : interdits, tabous, jeûnes).
    \item \textbf{La Communauté (Église) :} Le groupe moral uni par la foi commune. Pas de religion solitaire : « L'individu ne fait la religion que parce qu'il est en société. »
\end{itemize}

\begin{exemplebox}[Application camerounaise : Le culte des ancêtres chez les Bamiléké]
Chez les Bamiléké (Ouest Cameroun), la \textbf{croyance} affirme que les ancêtres (\emph{meu}) veillent sur la lignée. Le \textbf{rite} du versement de libations (vin de palme, huile) sur le sol ou au pied du \emph{tchoua} (arbre sacré) maintient le lien. La \textbf{communauté} est la famille élargie (le \emph{su}) rassemblée autour du chef de lignée. Cela remplit parfaitement la définition durkheimienne sans « Dieu » unique transcendant.
\end{exemplebox}

\begin{definition}[Définition substantialiste (Max Weber, Rudolf Otto)]
Centrée sur le \textbf{contenu} de la croyance : la relation à un \cle{Salut} (délivrance du mal, de la souffrance, du cycle des réincarnations) et à l'\cle{Au-delà}. Pour Rudolf Otto (\emph{Le Sacré}, 1917), l'essence du religieux est le \cle{\emph{Mysterium Tremendum et Fascinans}} : une expérience irrationnelle de la toute-puissance (tremendum) qui attire pourtant (fascinans).
\end{definition}

\begin{attention}[Piège : L'ethnocentrisme de la définition substantialiste]
Définir la religion par la « foi en un Dieu personnel » ou la « quête du salut » exclut de facto le \textbf{bouddhisme} (non-théiste, but = Nirvâna), le \textbf{confucianisme} (éthique sociale sans dogme surnaturel) et nombre de \textbf{traditions africaines} (relation aux forces vitales/ancêtres, pas « salut » individuel). La définition fonctionnaliste (Durkheim) est plus universelle pour l'analyse scientifique.
\end{attention}

\courssub{Distinctions cruciales : Religion, Magie, Superstition, Spiritualité, Secte}

Il est vital de ne pas confondre des phénomènes voisins. Le diagramme de Venn ci-dessous synthétise les critères sociologiques (institutionnalisation, rationalité, finalité).

\begin{popfigure}
\begin{tikzpicture}[
    set/.style={draw=popdark, thick, fill=popblue!15, opacity=0.6, circle, minimum width=5.5cm},
    labelstyle/.style={font=\small\bfseries, text=popdark},
    critstyle/.style={font=\tiny, text=popmuted, align=center}
]
% Cercles
\begin{scope}[blend group=soft light]
    \node[set] (R) at (0,0) {};
    \node[set] (M) at (3.5,2) {};
    \node[set] (S) at (3.5,-2) {};
    \node[set] (Sp) at (-3.5,0) {};
    \node[set] (Se) at (0,-4) {};
\end{scope}

% Labels principaux
\node[labelstyle] at (-4.5, 1.8) {Religion};
\node[labelstyle] at (5.5, 3.2) {Magie};
\node[labelstyle] at (5.5, -3.2) {Superstition};
\node[labelstyle] at (-4.5, -1.8) {Spiritualité};
\node[labelstyle] at (0, -6.2) {Secte};

% Critères (simplifiés pour lisibilité)
\node[critstyle, align=center] at (-4.5, 0.5){Institutionnalisée\\Communauté (Église)\\Salut / Sens global};
\node[critstyle, align=center] at (5.5, 1.8){Technique / Mécanique\\Contrainte des forces\\Utilité immédiate};
\node[critstyle, align=center] at (5.5, -1.8){Croyance erronée\\Irrationnelle\\Jugement de valeur};
\node[critstyle, align=center] at (-4.5, -0.5){Intériorité\\Déinstitutionnalisée\\Bien-être personnel};
\node[critstyle, align=center] at (0, -5.0){Rupture sociale\\Emprise mentale\\Gourou charismatique};

% Intersections clés (texte)
\node[font=\tiny\bfseries, text=poppurple, align=center] at (1.5, 0.5){Magie\\religieuse};
\node[font=\tiny\bfseries, text=poppurple, align=center] at (-1.5, 0.5){Religion\\vivante};
\node[font=\tiny\bfseries, text=poppurple, align=center] at (1.5, -1.5){Superstition\\populaire};
\node[font=\tiny\bfseries, text=poppurple] at (-1.5, -1.5) {Mystique};
\node[font=\tiny\bfseries, text=poporange, align=center] at (0, -3.0){Dérive\\sectaire};
\end{tikzpicture}
\legende{Diagramme de Venn : Intersections entre Religion, Magie, Superstition, Spiritualité et Secte. Critères : Institutionnalisation (forte pour Religion/Secte, faible pour Spiritualité/Magie), Relation au sacré (relationnelle pour Religion, mécanique pour Magie), Rationalité (dogmatique vs technique vs irrationnelle), Finalité (Salut global vs Utilité immédiate vs Bien-être).}
\end{popfigure}

\begin{propriete}[Distinctions opératoires]
\begin{itemize}
    \item \textbf{Religion vs Magie (Mauss, Durkheim) :} La magie est \cle{technique} : l'opérateur (mage/sorcier) contraint mécaniquement des forces impersonnelles (mana) par des formules/gestes exacts pour un résultat immédiat (guérir, nuire, faire pleuvoir). La religion est \cle{relationnelle} : le croyant \emph{supplie}, adore, se soumet à une volonté personnelle (Dieu, Esprits) ; l'efficacité n'est pas garantie (mystère de la grâce).
    \item \textbf{Religion vs Superstition :} La superstition est un \emph{jugement de valeur} (pégoratif) porté par un groupe dominant sur les croyances d'un groupe dominé ou sur des survivances archaïques (ex: « passer sous une échelle porte malheur »). Sociologiquement, c'est une \cle{croyance erronée en une causalité magique} sans fondement empirique ni institutionnalisation.
    \item \textbf{Religion vs Spiritualité :} La spiritualité moderne (New Age, développement personnel) privilégie l'\cle{intériorité}, l'\cle{autonomie} et le \cle{bien-être} subjectif, rejetant l'institution, le dogme et l'autorité cléricale. « Je suis spirituel mais pas religieux. »
    \item \textbf{Religion vs Secte (Dérives sectaires) :} Une secte n'est pas « une petite religion ». C'est un groupe à \cle{rupture sociale} (isolement familial/professionnel), \cle{emprise mentale} (contrôle de l'information, pensée unique), \cle{exploitation financière}, \cle{discours apocalyptique} et \cle{gourou charismatique} non contrôlé. \textbf{Préférer le terme juridique « dérive sectaire » (MIVILUDES).}
\end{itemize}
\end{propriete}

\courssub{Les fonctions sociales de la religion (Durkheim) : le ciment de la société}

Pour Durkheim, « Dieu, c'est la société divinisée ». La religion est le culte que la société se rend à elle-même.

\begin{aretenir}[Les 4 fonctions majeures (Durkheim / Parsons)]
\begin{enumerate}
    \item \textbf{Cohésion / Intégration (Identité collective) :} Les rites (fêtes, pèlerinages) créent l'\cle{effervescence collective} : fusion des consciences individuelles dans une conscience commune. Ex: La fête du Ngondo (Sawa, Littoral) ou le Mpo'o (Bassa) reforgent l'unité ethnique.
    \item \textbf{Légitimation de l'ordre social et moral :} Les normes ne sont pas seulement humaines, elles sont \cle{sacralisées} (volonté divine, ordre des ancêtres). « Tu ne tueras point » a plus de force si c'est un commandement divin. Le pouvoir politique s'en réclame (droit divin, serment sur la Bible/Coran).
    \item \textbf{Contrôle social :} La surveillance divine (œil de Dieu, ancêtres voyants) intériorise la norme. La culpabilité religieuse précède la sanction judiciaire. Confession, pénitence, peur de l'enfer/karmique = police intérieure peu coûteuse.
    \item \textbf{Support psychologique (Gestion de l'angoisse) :} Réponse à l'\cle{angoisse de la mort}, du hasard, de l'injustice, de la maladie. Rites de passage (naissance, puberté, mariage, mort) cadrent les crises existentielles. « La religion est le soupir de la créature opprimée » (Marx) mais aussi un anxiolytique social puissant.
\end{enumerate}
\end{aretenir}

\courssub{Dysfonctions, critiques et « retour du religieux »}

La religion n'est pas toujours fonctionnelle. Les classiques de la « méfiance » (Marx, Nietzsche, Freud) et la sociologie moderne pointent les \cle{dysfonctions}.

\begin{exoresolu}[Analyse critique : La religion, opium ou moteur de l'histoire ?]
\textbf{Énoncé :} « La religion est l'opium du peuple » (Marx). « La religion est une névrose obsessionnelle universelle » (Freud). Confrontez ces deux critiques à la réalité camerounaise actuelle.

\textbf{Solution détaillée :}
\begin{itemize}
    \item \textbf{Marx (Critique matérialiste) :} La religion console de la misère \emph{ici-bas} en promettant le bonheur \emph{là-haut}, désarmant la révolte sociale. \emph{Au Cameroun :} Certains discours prônent la résignation (« Dieu pourvoira », « C'est la volonté de Dieu ») face au chômage, à la corruption, aux routes impraticables. Cela peut freiner l'exigence citoyenne de droits et de services publics.
    \item \textbf{Freud (Critique psychanalytique) :} La religion est une illusion née du désir (père protecteur), régression infantile, névrose collective (rites = actes obsessionnels). \emph{Au Cameroun :} La prolifération des « églises de réveil » centrées sur la prospérité immédiate (miracles financiers, guérisons spectaculaires) peut relever d'une demande psychique de maîtrise magique du réel (proche de la magie pour Freud).
    \item \textbf{Nuance (Weber, Berger) :} La religion peut être \emph{révolutionnaire} (Théologie de la libération, rôle des églises dans la transition démocratique des années 90 au Cameroun, lutte anti-apartheid, droits civiques US). Elle fournit un \cle{capital social} (réseaux d'entraide, écoles, hôpitaux) que l'État défaillant ne fournit pas.
\end{itemize}
\textbf{Conclusion :} La religion est \emph{ambivalente} : elle peut aliéner (opium) ou libérer (prophétisme), diviser (fondamentalisme) ou unir (dialogue interreligieux).
\end{exoresolu}

\begin{propriete}[Weber : Le « Désenchantement du monde » (\emph{Entzauberung})]
La rationalisation occidentale (science expérimentale, droit codifié, bureaucratie, capitalisme) chasse les esprits, la magie, le mystère. Le monde devient \cle{calculable}, \cle{prévisible}, \cle{technique}. La religion perd sa prise sur le réel (explication de la foudre, de la maladie, des saisons) et se réfugie dans la sphère privée (éthique, sens de la vie). C'est le cœur de la thèse de la \cle{sécularisation}.
\end{propriete}

\begin{definition}[Sécularisation (sens sociologique : Berger, Wilson, Casanova)]
\cle{Processus historique de perte d'influence sociale des institutions religieuses} au profit d'institutions séculières (État, Marché, Science, Éducation, Santé).
\begin{itemize}
    \item \textbf{Niveau macro :} Différenciation des sphères (politique, droit, économie, science deviennent autonomes). L'État ne se réclame plus de Dieu.
    \item \textbf{Niveau méso :} Déclin de l'autorité religieuse sur l'éducation, la santé, la famille (mariage civil, divorce, bioéthique).
    \item \textbf{Niveau micro :} Baisse de la pratique dominicale, du vocabulaire religieux, de la croyance dogmatique (Europe de l'Ouest).
\end{itemize}
\textbf{Attention :} La sécularisation \textbf{ne signifie pas la disparition de la croyance individuelle} (« croire sans appartenir » - Grace Davie).
\end{definition}

\begin{savaistu}[Le paradoxe africain : « Retour du religieux » et vitalité]
Contrairement à l'Europe, l'Afrique (et le Cameroun) connaît une \textbf{explosion religieuse} (pentecôtisme, islam réformiste, réveil catholique charismatique). Thèse de la \textbf{régulation religieuse} (Finke, Stark) : là où l'État est faible (échec services publics), le « marché religieux » est concurrentiel et dynamique. Les églises offrent santé, éducation, réseau d'emplois, encadrement jeunesse. La religion remplit les vides de l'État.
\end{savaistu}

\courssub{Laïcité et pluralisme religieux au Cameroun : un modèle de « coopération »}

Le Cameroun n'applique ni la laïcité « stricte » à la française (séparation 1905, invisibilité religieuse dans l'espace public) ni le confessionnalisme (religion d'État). C'est une \cle{laïcité de coopération} (ou « laïcité positive »).

\begin{popfigure}
\begin{tikzpicture}[
    node distance=1.2cm and 1.5cm,
    box/.style={draw=popblue, thick, rounded corners, fill=popblue!10, minimum width=3.2cm, minimum height=1.2cm, align=center, font=\small},
    arrow/.style={-Stealth, thick, popdark},
    leg/.style={font=\tiny, text=popmuted, align=center}
]
% Niveau 1 : Cadre Constitutionnel
\node[box, fill=popgold!20, draw=popgold, align=center] (const){Cadre Constitutionnel\\Préambule 1996 : Laïcité,\\Liberté de conscience/culte};

% Niveau 2 : Lois Cadres
\node[box, below left=of const, align=center] (loi1){Loi 90/053\\Liberté d'association};
\node[box, below right=of const, align=center] (loi2){Loi 2000/001\\Associations (reconnaissance)};

% Niveau 3 : Acteurs Étatiques
\node[box, below=of loi1, align=center] (minat){MINAT\\Administration territoriale,\\Reconnaissance légale};
\node[box, below=of loi2, align=center] (minjec){MINJEC\\Jeunesse, Éducation civique,\\Aumôneries scolaires};
\node[box, right=of minjec, align=center] (minedub){MINEDUB/MINESEC\\Enseignement confessionnel\\sous contrat (subventions)};

% Niveau 4 : Organes Concertation
\node[box, below=of minat, fill=popgreen!15, draw=popgreen, align=center] (cnc){Concertations État-\\Confessions (MINAT)\\\& Plateformes interreligieuses};
\node[box, below=of minjec, fill=popgreen!15, draw=popgreen, align=center] (plateforme){Plateforme Interreligieuse\\Cameroun (CEC, CPC,\\Muslim Council)};

% Niveau 5 : Réalisations Concrètes
\node[box, below=of cnc, fill=poppink!15, draw=poppink, align=center] (fet){Fêtes chômées légales\\Noël, Ascension, Assomption,\\Aïd el-Fitr, Aïd al-Kabir,\\Fête Nationale 20 mai};
\node[box, below=of plateforme, fill=poppink!15, draw=poppink, align=center] (aum){Aumôneries\\Armée, Police, Prisons,\\Hôpitaux, Universités};
\node[box, right=of aum, fill=poppink!15, draw=poppink, align=center] (mariage){Mariages\\Coutumier / Religieux / Civil\\Art. 52 Code Civil :\\Célébration par officier d'état civil};

% Flèches
\draw[arrow] (const) -- (loi1);
\draw[arrow] (const) -- (loi2);
\draw[arrow] (loi1) -- (minat);
\draw[arrow] (loi2) -- (minjec);
\draw[arrow] (loi2) -- (minedub);
\draw[arrow] (minat) -- (cnc);
\draw[arrow] (minjec) -- (plateforme);
\draw[arrow] (minedub) -- (plateforme);
\draw[arrow] (cnc) -- (fet);
\draw[arrow] (plateforme) -- (aum);
\draw[arrow] (plateforme) -- (mariage);

% Légende
\node[leg] at (0, -9.5) {Flux : Cadre juridique $\to$ Acteurs étatiques $\to$ Concertation $\to$ Gestion concrète du fait religieux};
\end{tikzpicture}
\legende{Carte mentale : Gestion publique du fait religieux au Cameroun. Modèle de « laïcité de coopération » : l'État neutre régule, finance (écoles confessionnelles, pèlerinages) et concertation avec les leaders religieux pour la paix sociale.}
\end{popfigure}

\begin{exemplebox}[Chiffres clés : Laïcité de coopération en acte (Estimations budgétaires)]
\begin{itemize}
    \item \textbf{Pèlerinages :} L'État subventionne le Hadj (La Mecque) et les pèlerinages chrétiens (Jérusalem, Rome, Marial) : \textbf{plusieurs milliards de FCFA/an} (billets d'avion, logistique, encadrement médical).
    \item \textbf{Écoles confessionnelles :} \textasciitilde 30\% des élèves du secondaire dans le privé confessionnel (catholique, protestant, islamique). L'État paie une partie des salaires (enseignants « mis à disposition » ou subventions de fonctionnement).
    \item \textbf{Fêtes chômées :} 5 fêtes religieuses officielles (3 chrétiennes : Ascension, Assomption, Noël ; 2 musulmanes : Aïd el-Fitr, Aïd al-Kabir) sur \textasciitilde 9 fêtes chômées totales (avec 1er jan, 11 fév, 1er mai, 20 mai). Équilibre symbolique.
\end{itemize}
\end{exemplebox}

\begin{attention}[Erreur fréquente au Bac : Laïcité $\neq$ Laïcisme]
\begin{itemize}
    \item \textbf{Laïcité (Principe constitutionnel camerounais) :} Séparation de l'État et des religions. L'État est \cle{neutre} (pas de religion d'État), il \cle{protège} la liberté de conscience et de culte, il \cle{coopère} avec les confessions (intérêt général : éducation, santé, paix sociale).
    \item \textbf{Laïcisme (Idéologie) :} Volonté d'\emph{exclure} le religieux de l'espace public (interdiction du voile à l'école, suppression fêtes religieuses, laïcité « combattante »). \textbf{Ce n'est PAS le modèle camerounais.}
\end{itemize}
Dans une dissertation sur « La laïcité au Cameroun », ne pas confondre les deux. Citer le Préambule de la Constitution 1996 : « L'État est laïc... garantit le libre exercice des cultes. »
\end{attention}

\courssub{Atelier d'application : Qualifier un groupe religieux}

\begin{methode}[Comment analyser un « fait religieux » dans un sujet de dissertation ou commentaire ?]
\begin{enumerate}
    \item \textbf{Identifier le type de définition pertinente :} Fonctionnelle (lien social, rites, communauté) ou Substantialiste (contenu dogmatique, salut). \emph{Choisir selon le sujet.}
    \item \textbf{Distinguer le niveau d'analyse :} Individuel (croyance, expérience mystique) vs Collectif (institution, pouvoir, contrôle social).
    \item \textbf{Repérer les fonctions ET dysfonctions :} Ne pas faire l'apologie ni le procès. Cohésion \emph{et} exclusion. Apaisement \emph{et} fanatisme.
    \item \textbf{Situer dans l'espace public :} Quel rapport à la laïcité ? Soumission, négociation, conflit, indifférence ? Contexte camerounais obligatoire si le sujet l'implique.
    \item \textbf{Mobiliser les auteurs :} Durkheim (fonction), Weber (sens/désenchantement), Marx/Freud (critique), Berger (sécularisation/retour), auteurs africains (Mbiti, Nyamnjoh, Fancello) pour le contexte.
\end{enumerate}
\end{methode}

\begin{exercice}[Sujet type Bac Tech : Dissertation]
\textbf{« La religion est-elle encore le ciment de la société camerounaise ? »}
\textbf{Pistes de plan (dialectique) :}
\begin{enumerate}
    \item \textbf{Thèse : Oui, ciment majeur.} Fonctions durkheimiennes : Cohésion ethnique (rites traditionnels), légitimation morale (serments, mariages), substitut à l'État défaillant (écoles, santé, microcrédit via tontines paroissiales), régulateur de paix (Plateforme interreligieuse, appel au calme en période électorale).
    \item \textbf{Antithèse : Non, facteur de division et de fragilisation.} Dysfonctions : Conflits inter/intra religieux (radicalisation Boko Haram au Nord, tensions chrétiens/musulmans, guerres de leadership dans églises), justification de l'oppression (femmes, minorités sexuelles, sorcellerie), détournement fonds (prospérité), emprise sectaire, affaiblissement de la laïcité (instrumentalisation politique du religieux).
    \item \textbf{Synthèse : Un ciment fissuré, à réinventer.} La religion reste un \emph{capital social} indispensable mais insuffisant. Le vrai ciment moderne doit être \textbf{civique} (citoyenneté, droit, institutions républicaines fortes). La laïcité de coopération doit évoluer vers une \textbf{régulation transparente} (financement, formation des pasteurs/imams, lutte contre dérives) pour que le religieux serve le \emph{bien commun} et non l'entre-soi communautaire.
\end{enumerate}
\end{exercice}

\begin{corrige}[Exercice n°1 : Sujet type Bac Tech]
\textbf{Attendus du correcteur :}
\begin{itemize}
    \item Définition claire (Durkheim) en intro.
    \item Exemples camerounais \textbf{précis et datés} (Ngondo, Hadj subventionné, crise NOSO rôle églises, Boko Haram, loi 2000/001).
    \item Articulation fonctions/dysfonctions équilibrée.
    \item Notion de laïcité correctement manipulée (coopération vs exclusion).
    \item Conclusion ouverte sur le politique (citoyenneté vs identité religieuse).
\end{itemize}
\end{corrige}

\coursec{Foi, Raison et Laïcité : le défi camerounais}

Nous avons vu, dans la section précédente, que la religion est un fait social total : elle structure les communautés, fonde des normes et répond à un besoin de sens. Mais une question politique et philosophique demeure : que faire de la religion dans un État moderne, pluraliste, où coexistent chrétiens, musulmans, adeptes des religions traditionnelles et non-croyants ? Comment la \cle{foi} peut-elle dialoguer avec la \cle{raison} ? Et quel statut juridique accorder aux normes religieuses face à la loi civile ? C'est tout l'enjeu de la \cle{laïcité}, que nous abordons maintenant à travers le cas camerounais.

\courssub{Foi et Raison : conflit, hiérarchie ou complémentarité ?}

\begin{definition}[Foi et Raison]
La \cle{foi} est l'adhésion confiante à des vérités révélées, qui ne peuvent être démontrées (Dieu, le salut, les dogmes). La \cle{raison} est la faculté de connaître par l'argumentation et la démonstration. Le problème philosophique est le suivant : la foi a-t-elle besoin de la raison ? Peuvent-elles se contredire ?
\end{definition}

Quatre positions classiques structurent le débat :

\begin{enumerate}
\item \textbf{Le fidéisme} (Tertullien, II\textsuperscript{e}--III\textsuperscript{e} siècle ; Kierkegaard, XIX\textsuperscript{e} siècle) : la foi seule suffit, elle n'a pas besoin de la raison. La tradition attribue à Tertullien le paradoxe célèbre : \emph{Credo quia absurdum} (« Je crois parce que c'est absurde »), formule qui résume son idée que les vérités de la foi dépassent, voire choquent, la raison humaine. Pour Kierkegaard, la foi est un « saut » existentiel au-delà de toute preuve : croire parce qu'on a démontré, ce ne serait plus croire.
\item \textbf{Le rationalisme théologique} (Thomas d'Aquin, XIII\textsuperscript{e} siècle) : la raison peut démontrer certaines vérités religieuses (l'existence de Dieu par les « cinq voies »), tandis que les mystères (la Trinité) relèvent de la foi. Foi et raison sont harmonieuses car elles viennent du même Dieu.
\item \textbf{La critique kantienne} (Kant, XVIII\textsuperscript{e} siècle) : la raison théorique ne peut ni prouver ni réfuter Dieu, qui dépasse les limites de toute expérience possible. Kant « limite le savoir pour faire place à la foi » : Dieu devient un postulat de la raison \emph{pratique} (morale), non un objet de science.
\item \textbf{La philosophie contemporaine de la religion} (Alvin Plantinga, XX\textsuperscript{e}--XXI\textsuperscript{e} siècle) : la croyance en Dieu peut être une « croyance de base » (\emph{properly basic belief}), rationnellement justifiée (\emph{warrant}) sans passer par des preuves, comme la croyance en l'existence d'autrui ou du monde extérieur.
\end{enumerate}

\begin{aretenir}[L'essentiel]
Le débat foi/raison n'oppose pas « croyants irrationnels » et « savants rationnels ». Il porte sur la \emph{nature} de la justification : la foi est-elle un savoir à démontrer (Thomas), un saut au-delà du savoir (Kierkegaard), une limite du savoir (Kant) ou une croyance de base légitime (Plantinga) ?
\end{aretenir}

\courssub{Le principe de laïcité : séparation, neutralité, liberté, égalité}

\begin{definition}[Laïcité]
La \cle{laïcité} est le principe politique qui organise la relation entre l'État et les religions selon quatre piliers : la \textbf{séparation} des institutions (l'État ne dépend d'aucune Église, les Églises ne gouvernent pas l'État), la \textbf{neutralité} (l'État ne favorise ni ne défavorise aucune religion), la \textbf{liberté} de conscience et de culte (croire, ne pas croire, changer de religion, pratiquer), et l'\textbf{égalité} (aucune discrimination fondée sur la religion).
\end{definition}

Il n'existe pas un mais plusieurs modèles de laïcité :

\begin{center}
\begin{tabularx}{\linewidth}{|l|X|X|}\hline
\rowcolor{popblueL} \textbf{Modèle} & \textbf{Principe} & \textbf{Exemple} \\ \hline
Français (strict) & Séparation rigide ; la religion relève de la sphère privée ; aucun signe religieux ostensible à l'école publique. & Loi de 1905 \\ \hline
Américain & Séparation stricte mais forte visibilité publique de la religion (« \emph{In God We Trust} »). & États-Unis \\ \hline
Allemand & Coopération : l'État collabore avec les Églises (cours de religion, impôt cultuel). & Allemagne \\ \hline
Camerounais & Coopération régulée : conventions avec les écoles et hôpitaux confessionnels. & Cameroun \\ \hline
\end{tabularx}
\end{center}

\begin{definition}[Laïcité de « coopération » (modèle camerounais/allemand)]
L'État reconnaît l'utilité sociale des religions et leur délègue des missions de service public (éducation, santé, action sociale) par des \textbf{conventions}, des \textbf{subventions} et le paiement de personnels, tout en gardant le contrôle : programmes officiels, normes d'hygiène, inspection. La religion devient ainsi un \emph{partenaire} de l'État, non un concurrent ni un ennemi.
\end{definition}

\begin{exemplebox}[Un partenariat chiffré]
Au Cameroun, des milliers d'écoles confessionnelles (catholiques, protestantes, islamiques) scolarisent une part importante des élèves — environ un tiers des effectifs du secondaire selon les estimations courantes. L'État paie une partie des salaires des enseignants contractuels et subventionne le fonctionnement de ces établissements, qui appliquent les programmes officiels MINESEC/MINESUP. Sans ce partenariat, le système éducatif public serait saturé.
\end{exemplebox}

\begin{attention}[Piège fréquent]
Écrire « le Cameroun est un pays chrétien » (ou « musulman ») est une erreur juridique et philosophique. Le Cameroun est un \textbf{État laïc à majorité croyante} : le préambule de la Constitution de 1996 proclame la laïcité de l'État et garantit la liberté de conscience et de culte. La démographie religieuse ne fonde pas la légitimité juridique de l'État : même à 99\,\% croyant, un État laïc reste tenu à la neutralité et à l'égalité de tous les citoyens, y compris les athées et les minorités.
\end{attention}

\courssub{Les enjeux camerounais actuels : quand les normes entrent en conflit}

\textbf{1. Éducation.} Le statut des écoles confessionnelles sous contrat pose la question de la sous-traitance du service public : l'État délègue, mais doit garantir l'égalité des programmes. Le port du voile ou du hijab à l'école publique cristallise le conflit entre liberté religieuse de l'élève et neutralité de l'espace scolaire.

\textbf{2. Droit de la famille.} Le droit camerounais reconnaît aussi bien le mariage monogamique que le mariage polygamique (l'option est choisie lors du mariage civil), mais les successions opposent souvent la coutume (qui peut exclure la fille de l'héritage) au droit écrit et aux conventions internationales. Le mariage forcé ou précoce, justifié par certaines traditions, viole le droit de l'enfant.

\textbf{3. Santé.} Le refus de transfusion sanguine par les Témoins de Jéhovah oppose la liberté religieuse au devoir de sauver une vie (surtout celle d'un mineur). Les réticences vaccinales d'inspiration religieuse et le recours exclusif à la médecine traditionnelle posent le même problème : jusqu'où la conviction personnelle peut-elle prévaloir sur la santé publique ?

\textbf{4. Extrémisme violent.} Dans l'Extrême-Nord, Boko Haram illustre la perversion de la religion en idéologie meurtrière. La réponse camerounaise combine sécurité, déradicalisation et programmes DDR (Désarmement, Démobilisation, Réinsertion). Philosophiquement, cela rappelle que la liberté de conscience a des limites : elle ne protège jamais l'appel à la violence.

\begin{methode}[Analyser un cas concret : le voile à l'école]
\begin{enumerate}
\item \textbf{Identifier les droits en jeu} : liberté religieuse de l'élève contre neutralité de l'école publique et égalité entre élèves.
\item \textbf{Identifier le niveau} : école publique (neutralité stricte possible) ou école privée sous contrat (règlement intérieur propre, dans le respect de la loi).
\item \textbf{Appliquer le test de proportionnalité} : la restriction poursuit-elle un but légitime (vivre-ensemble, égalité filles/garçons) ? Est-elle nécessaire ? Est-elle proportionnée (interdiction totale ou encadrement) ?
\item \textbf{Conclure} en distinguant ce qui relève de la loi, du règlement et de la pédagogie du dialogue.
\end{enumerate}
\end{methode}

\courssub{Le dialogue interreligieux : la PCRC et la cohésion nationale}

\begin{savaistu}[La PCRC, acteur de paix]
La Plateforme des Confessions Religieuses du Cameroun (PCRC), qui réunit autorités catholiques, protestantes et musulmanes, a joué un rôle de médiation dans la crise anglophone (depuis 2016) et dans la gestion du Covid-19 (sensibilisation des fidèles, acceptation des vaccins). Elle intervient aussi dans l'observation citoyenne des élections et la prévention des conflits.
\end{savaistu}

Le dialogue interreligieux a une \textbf{valeur politique} : il montre que des traditions différentes peuvent s'accorder sur des exigences éthiques communes (la paix, la dignité, la justice). Mais il a des \textbf{limites} : risque d'instrumentalisation politique (les chefs religieux récupérés par le pouvoir ou l'opposition), et exclusion des minorités — religions traditionnelles, libres penseurs, athées — qui ne siègent pas à la table des « grandes » confessions. Une laïcité authentique doit protéger aussi ceux qui ne sont représentés par aucune Église.

\courssub{Éthique laïque et éthiques religieuses : la loi civile prime-t-elle ?}

Les conflits de normes sont particulièrement vifs en \textbf{bioéthique} (avortement, procréation médicalement assistée, fin de vie), sur les \textbf{droits des personnes homosexuelles} (l'article 347-1 du Code pénal camerounais réprime les relations homosexuelles, en tension avec les standards internationaux des droits humains) et sur l'\textbf{égalité homme/femme} (héritage, tutelle). La question philosophique est : dans un État laïc, la loi civile prime-t-elle la norme religieuse ?

La réponse de principe est oui : la loi vaut pour tous, indépendamment des convictions. Mais cette primauté doit être \emph{justifiée} — et c'est là qu'intervient la philosophie politique contemporaine.

\begin{propriete}[Le devoir de civilité (Rawls)]
Selon John Rawls, les citoyens et surtout les députés sont tenus par un \textbf{devoir de civilité} (\emph{duty of civility}) : ils doivent pouvoir justifier les lois coercitives par des raisons \textbf{publiques}, c'est-à-dire accessibles à la raison commune de tous les citoyens, et non par la seule autorité d'un texte sacré. Dire « la loi doit interdire X parce que la Bible (ou le Coran) le dit » n'est pas une raison publique ; dire « la loi doit protéger la dignité égale de chaque personne » en est une.
\end{propriete}

Jürgen Habermas prolonge Rawls dans sa théorie de la \textbf{« société post-séculière »} : les traditions religieuses portent des « apports sémantiques » précieux (le sens de la dignité, du pardon, de la solidarité) que l'État laïc ne doit pas mépriser mais \textbf{traduire} en un langage universellement accessible. La raison publique fonctionne alors comme une machine de traduction entre la sphère privée des convictions et la sphère publique de la loi.

\begin{popfigure}
\begin{center}
\begin{tikzpicture}[
box/.style={draw=popblue, thick, rounded corners, fill=popblueL, text width=3.4cm, align=center, minimum height=1.6cm, font=\small},
trad/.style={draw=poporange, thick, rounded corners, fill=poporangeL, text width=2.6cm, align=center, minimum height=1.4cm, font=\small},
fleche/.style={-{Stealth[length=3mm]}, very thick, popdark}]
\node[box, align=center] (prive) at (0,0){\textbf{Sphère privée}\\ Raisons religieuses (textes sacrés, traditions, convictions)};
\node[trad, align=center] (trad) at (5.6,0){\textbf{Traduction}\\ Raison publique (Rawls/Habermas)};
\node[box, align=center] (pub) at (11,0){\textbf{Sphère publique}\\ Raisons accessibles à tous les citoyens};
\node[box, fill=popgreenL, draw=popgreen, align=center] (loi) at (11,-2.6){\textbf{Loi légitime}\\ Valable pour tous, croyants et non-croyants};
\draw[fleche] (prive) -- (trad);
\draw[fleche] (trad) -- (pub);
\draw[fleche] (pub) -- (loi);
\draw[fleche, poporange, dashed] (pub.south west) to[bend right=25] node[below, font=\scriptsize, text width=4.5cm, align=center]{Débat démocratique : la sphère publique interroge et enrichit les traditions} (prive.south east);
\end{tikzpicture}
\end{center}
\legende{Le modèle de la « raison publique » : les convictions religieuses, légitimes dans la sphère privée, doivent être traduites en raisons accessibles à tous pour fonder une loi légitime. Le flux est bidirectionnel : le débat public enrichit aussi les traditions.}
\end{popfigure}

Prenons le cas de l'héritage des filles. La coutume, le droit musulman, le Code civil et la CEDAW (Convention de l'ONU sur l'élimination de toutes les formes de discrimination à l'égard des femmes, ratifiée par le Cameroun en 1994) peuvent prescrire des solutions différentes. Comment l'État arbitre-t-il ?

\begin{popfigure}
\begin{center}
\begin{tikzpicture}[
norme/.style={draw=poppurple, thick, rounded corners, fill=poppurpleL, text width=2.5cm, align=center, minimum height=1.1cm, font=\scriptsize},
inst/.style={draw=popblue, thick, rounded corners, fill=popblueL, text width=2.9cm, align=center, minimum height=1.2cm, font=\small},
fleche/.style={-{Stealth[length=2.5mm]}, thick, popdark}]
\node[norme] (cout) at (0,2.6) {Coutume};
\node[norme] (char) at (0,1.3) {Droit musulman};
\node[norme] (civ) at (0,0) {Code civil};
\node[norme] (cedaw) at (0,-1.3) {CEDAW (ONU)};
\node[inst, align=center] (trib) at (4.6,0.6){\textbf{Tribunal}\\ Cas concret : héritage d'une fille};
\node[inst, align=center] (sup) at (8.9,0.6){\textbf{Cour Suprême}\\ Jurisprudence};
\node[inst, fill=popgreenL, draw=popgreen, align=center] (harm) at (8.9,-1.8){\textbf{Harmonisation législative}\\ Réforme de la loi};
\draw[fleche] (cout.east) -- (trib.west);
\draw[fleche] (char.east) -- (trib.west);
\draw[fleche] (civ.east) -- (trib.west);
\draw[fleche] (cedaw.east) -- (trib.west);
\draw[fleche] (trib) -- (sup);
\draw[fleche] (sup) -- (harm);
\draw[fleche, poporange, dashed] (harm.west) to[bend left=20] node[below, font=\scriptsize]{la loi réformée unifie les normes} (civ.south);
\end{tikzpicture}
\end{center}
\legende{Gestion d'un conflit de normes (exemple : l'héritage des femmes). Quatre sources de droit concourent ; le juge tranche le cas concret, la Cour Suprême fixe la jurisprudence, et le législateur peut harmoniser la loi. Dans un État laïc, le droit civil et les engagements internationaux priment en dernier ressort.}
\end{popfigure}

\begin{exoresolu}[Sujet de dissertation : « La religion a-t-elle sa place dans l'État ? »]
\textbf{Énoncé.} Un candidat écrit : « La religion doit être totalement bannie de la vie publique, car elle est source de division. » Discutez cette affirmation.
\tcblower
\textbf{Corrigé (plan détaillé).}
\begin{enumerate}
\item \textbf{Thèse} : la religion menace l'unité politique. Arguments : les guerres de religion, le terrorisme (Boko Haram), les conflits de normes (polygamie, voile). La laïcité stricte (modèle français) protège la paix civile en reléguant la religion au privé.
\item \textbf{Antithèse} : mais la religion est aussi une ressource sociale et morale. Elle scolarise une part importante des élèves camerounais, médie les conflits (PCRC dans la crise anglophone), porte des valeurs de solidarité. Habermas : priver l'espace public des « apports sémantiques » religieux appauvrirait la démocratie.
\item \textbf{Synthèse} : la solution n'est ni l'exclusion ni la soumission, mais la \emph{traduction} (Rawls/Habermas) : les convictions religieuses sont légitimes dans le débat, à condition d'être formulées en raisons publiques. Le modèle camerounais de laïcité de coopération illustre cette voie médiane : partenariat avec les religions, primauté de la loi civile.
\end{enumerate}
\textbf{Conclusion} : la religion a sa place dans l'État comme \emph{interlocuteur}, jamais comme \emph{législateur}.
\end{exoresolu}

\begin{exercice}[Pour s'entraîner]
1. Définis la laïcité et distingue ses quatre piliers. 2. Explique la différence entre le fidéisme de Kierkegaard et le rationalisme de Thomas d'Aquin. 3. Applique la méthode d'analyse de cas au refus de transfusion sanguine d'un mineur par ses parents Témoins de Jéhovah. 4. Dissertation : « La loi civile doit-elle primer la norme religieuse ? »
\end{exercice}

\begin{corrige}[Exercice 5]
1. La laïcité organise la relation État/religions par la séparation des institutions, la neutralité de l'État, la liberté de conscience et de culte, et l'égalité de tous devant la loi. 2. Kierkegaard soutient que la foi est un saut existentiel au-delà de toute preuve (fidéisme) ; Thomas d'Aquin soutient que la raison peut démontrer certaines vérités religieuses, foi et raison étant harmonieuses. 3. Droits en jeu : liberté religieuse des parents contre droit à la vie de l'enfant ; le juge applique le test de proportionnalité : sauver une vie est un but légitime, la transfusion est nécessaire et proportionnée ; conclusion : le droit à la vie du mineur prime, la loi autorise le juge à ordonner la transfusion. 4. Plan attendu : thèse (la loi civile prime car elle vaut pour tous — neutralité, égalité) ; antithèse (la norme religieuse fonde le sens moral de nombreux citoyens, la mépriser serait illégitime) ; synthèse (primauté de la loi civile, mais justifiée par des raisons publiques et ouverte à la traduction des apports religieux — Rawls, Habermas ; exemple camerounais de l'héritage et de la CEDAW).
\end{corrige}

\coursec{Atelier méthodologique : Dissertation et Explication de texte (Bac Tech)}

Après avoir exploré les enjeux contemporains de la laïcité et du fait religieux au Cameroun (section 5), nous devons maintenant outiller l'élève pour l'épreuve terminale. La philosophie au Bac Technique ne sanctionne pas une culture générale, mais une \cle{compétence argumentative} : savoir transformer un sujet en problème, structurer une démonstration rigoureuse et mobiliser des concepts comme des outils de pensée. Cette section est votre « boîte à outils » pour l'examen.

\courssub{Analyse du sujet type Bac : repérage, concepts et tension}

Tout sujet de dissertation est un \emph{piège} tendu à l'élève qui répondrait « oui » ou « non » trop vite. La première minute détermine la note.

\begin{methode}[Décortiquer un sujet en 4 temps (Brouillon)]
\begin{enumerate}
    \item \textbf{Isoler les concepts clés} : Les termes philosophiques porteurs de sens (ex: \cle{Religion}, \cle{Progrès}, \cle{Preuve}, \cle{Existence}, \cle{Laïcité}, \cle{Fait religieux}). Les définir \emph{avant} d'écrire l'introduction.
    \item \textbf{Identifier la tension (le « vs »)} : Un sujet philosophique oppose toujours deux valeurs, deux logiques ou deux ordres de réalité. « La religion est-elle un obstacle au progrès ? » $\rightarrow$ Tension : \emph{Fidélité au sacré / Tradition} vs \emph{Rupture scientifique / Émancipation}.
    \item \textbf{Délimiter le champ} : \textbf{Philosophie}, pas théologie (pas de citation biblique/coranique comme argument), pas sociologie pure (pas de stats sans analyse), pas opinion (« Je pense que... »).
    \item \textbf{Formuler la problématique} : Transformer le sujet en question philosophique précise qui guide le plan. Voir boîte \textbf{Méthode} ci-dessous.
\end{enumerate}
\end{methode}

\begin{definition}[Problématiser]
Transformer le sujet imposé en une \cle{question philosophique tensionnée} qui met en jeu des concepts opposés et ouvre un espace de débat rationnel. Ce n'est pas répéter le sujet, c'en est le \emph{moteur}.
\end{definition}

\begin{methode}[Problématiser = Mettre en tension]
\begin{itemize}
    \item \textbf{Sujet :} « La religion est-elle un obstacle au progrès ? »
    \item \textbf{Mauvaise problématique :} « Est-ce que la religion empêche le progrès ? » (Répétition, binaire).
    \item \textbf{Bonne problématique :} « La quête de sens et le lien au sacré (religion) sont-ils structurellement incompatibles avec l'exigence de rationalité critique et de transformation technique du monde (progrès), ou peuvent-ils en être le ferment éthique ? »
    \item \textbf{Sujet :} « Peut-on prouver l'existence de Dieu ? »
    \item \textbf{Bonne problématique :} « La raison humaine, limitée à l'expérience sensible, peut-elle atteindre par la démonstration (preuve \emph{a posteriori} ou \emph{a priori}) un être infiniment transcendant, ou l'accès à Dieu relève-t-il d'un autre registre (foi, pari, intuition) ? »
    \item \textbf{Sujet :} « L'État laïc doit-il ignorer le fait religieux ? »
    \item \textbf{Bonne problématique :} « La neutralité de l'État (laïcité) implique-t-elle l'invisibilité du religieux dans l'espace public, ou la gestion démocratique du pluralisme impose-t-elle une \emph{reconnaissance} institutionnelle du fait religieux sans confusion des sphères ? »
\end{itemize}
\end{methode}

\courssub{Architecture de la dissertation philosophique}

La dissertation n'est pas un « pour/contre ». C'est une \cle{dialectique} : chaque partie répond à la précédente, la nie ou la dépasse pour aboutir à une synthèse qui \emph{transforme} le problème initial.

\begin{popfigure}
\centering
\begin{tikzpicture}[
    node distance=1.2cm and 1.5cm,
    box/.style={draw=popblue, thick, rounded corners, fill=popblueL, text width=4.5cm, align=center, minimum height=2.2cm, font=\small},
    arrow/.style={-Stealth, thick, popdark},
    title/.style={font=\bfseries\large, text=popblue},
    subtitle/.style={font=\itshape\small, text=popmuted}
]
% Introduction - Entonnoir
\node[box, title, align=center] (intro){\textbf{INTRODUCTION}\\(Entonnoir)\\\textit{Accroche $\rightarrow$ Définitions $\rightarrow$ Problématique $\rightarrow$ Annonce du plan}};

% Development - 3 Colonnes
\node[box, below left=of intro.south west, xshift=-1.5cm, fill=popgreenL, draw=popgreen, align=center] (thesis){\textbf{I. THÈSE}\\(L'évidence / Le sens commun)\\\textit{Argument 1 + Exemple\\Argument 2 + Exemple\\$\rightarrow$ Transition : « Mais... »}};
\node[box, below=of intro, fill=poporangeL, draw=poporange, align=center] (antithesis){\textbf{II. ANTITHÈSE}\\(La critique / Le retournement)\\\textit{Argument 1 + Exemple\\Argument 2 + Exemple\\$\rightarrow$ Transition : « Toutefois... »}};
\node[box, below right=of intro.south east, xshift=1.5cm, fill=poppurpleL, draw=poppurple, align=center] (synthesis){\textbf{III. SYNTHÈSE}\\(Le dépassement / La nuance)\\\textit{Nouveau concept opératoire\\Réconciliation conditionnelle\\Réponse à la problématique}};

% Arrows from Intro
\draw[arrow] (intro.south west) -- ++(-0.5,-0.5) -| (thesis.north);
\draw[arrow] (intro.south) -- (antithesis.north);
\draw[arrow] (intro.south east) -- ++(0.5,-0.5) -| (synthesis.north);

% Links between parts (Dialectic)
\draw[arrow, dashed, popmuted] (thesis.east) to[bend left=20] node[above, sloped, pos=0.5, font=\tiny] {Négation / Limite} (antithesis.west);
\draw[arrow, dashed, popmuted] (antithesis.east) to[bend left=20] node[above, sloped, pos=0.5, font=\tiny] {Dépassement / \emph{Aufhebung}} (synthesis.west);

% Conclusion - Entonnoir inversé
\node[box, below=of antithesis, yshift=-0.5cm, fill=popgoldL, draw=popgold, title, align=center] (conclu){\textbf{CONCLUSION}\\(Entonnoir inversé)\\\textit{Bilan synthétique $\rightarrow$ Réponse définitive $\rightarrow$ Ouverture}};

\draw[arrow] (thesis.south) |- (conclu.north west);
\draw[arrow] (antithesis.south) -- (conclu.north);
\draw[arrow] (synthesis.south) |- (conclu.north east);

% Labels on sides
\node[left=0.5cm of thesis, rotate=90, font=\bfseries\small, text=popgreen] {Moment 1 : Affirmation};
\node[left=0.5cm of antithesis, rotate=90, font=\bfseries\small, text=poporange] {Moment 2 : Négation};
\node[left=0.5cm of synthesis, rotate=90, font=\bfseries\small, text=poppurple] {Moment 3 : Dépassement};
\end{tikzpicture}
\legende{Architecture standard de la dissertation philosophique : du particulier (introduction) au général (conclusion) par le mouvement dialectique à trois temps.}
\end{popfigure}

\begin{aretenir}[Structure type au Bac Tech]
\begin{itemize}
    \item \textbf{Introduction (1 page max)} : Accroche (citation, fait d'actualité camerounais, paradoxe) $\rightarrow$ \textbf{Définitions rigoureuses} des concepts clés (3 max) $\rightarrow$ \textbf{Problématique} (une phrase, point d'interrogation) $\rightarrow$ \textbf{Annonce du plan} (I, II, III avec verbes d'action : \emph{Analyser, Montrer, Dépasser}).
    \item \textbf{Développement (3 grandes parties $\approx$ 3--4 pages)} : Chaque partie = 2 à 3 sous-parties (A, B, C). \textbf{Un paragraphe = Une idée directrice}. Structure paragraphe : \textbf{Idée} $\rightarrow$ \textbf{Argument philosophique (concept/auteur)} $\rightarrow$ \textbf{Exemple concret (Cameroun si possible)} $\rightarrow$ \textbf{Lien/Transition}.
    \item \textbf{Conclusion (1/2 page)} : \textbf{Bilan} (reprendre le fil de I, II, III sans liste de courses) $\rightarrow$ \textbf{Réponse finale} à la problématique $\rightarrow$ \textbf{Ouverture} (autre sujet, autre auteur, actualité brûlante).
\end{itemize}
\end{aretenir}

\courssub{L'art du paragraphe argumentatif : l'auteur comme outil}

\begin{attention}[Citer $\neq$ Argumenter]
\textbf{Erreur fatale :} « Comme le dit Marx, la religion est l'opium du peuple, donc elle freine le progrès. » $\rightarrow$ C'est un \emph{argument d'autorité}, non philosophique.
\textbf{Bonne pratique :} « Marx analyse la religion comme une \cle{consolation illusoire} qui détourne l'homme de sa \cle{réalité matérielle} (argument conceptuel). Au Cameroun, certains mouvements de réveil promettent la prospérité par la prière seule, retardant l'engagement entrepreneurial ou syndical (exemple local). \emph{Cependant}, cette analyse néglige le rôle moteur de l'éthique protestante dans le capitalisme (Weber)... » $\rightarrow$ L'auteur est un \textbf{témoin expert} dont on \emph{explique le concept} pour faire avancer \emph{sa propre démonstration}.
\end{attention}

\begin{exemplebox}[Paragraphe type corrigé (Extrait Thèse : Religion $\rightarrow$ Progrès social)]
\textbf{Idée directrice (A) :} La religion structure le lien social et fonde une éthique du travail favorable au développement.
\textbf{Argument (Weber) :} Max Weber, dans \emph{L'Éthique protestante et l'esprit du capitalisme}, montre que l'\cle{ascétisme intramondain} calviniste (devoir de réussir dans son métier comme signe de grâce) a converti l'énergie religieuse en \cle{rationalisation économique}.
\textbf{Illustration camerounaise :} Les réseaux d'écoles et de santé confessionnels (CBC, EPC, Catholique, Islamique) forment l'élite technique et soignent 40\% de la population (chiffres MINSANTE), palliant les défaillances de l'État : la religion \emph{est} ici un opérateur de progrès humain.
\textbf{Transition (vers B) :} Mais cette contribution sociale ne nie-t-elle pas le risque d'aliénation critique pointé par Marx ou Freud ?
\end{exemplebox}

\courssub{L'explication de texte : méthode en 4 étapes}

L'explication n'est pas un résumé, ni un commentaire linéaire « au fil de l'eau ». C'est une \emph{reconstruction logique} de la pensée de l'auteur.

\begin{methode}[Les 4 temps de l'explication (Bac Tech -- 1h30 env.)]
\begin{enumerate}
    \item \textbf{Contextualisation (10 min)} : Auteur, Œuvre, Date, École/Contexte intellectuel, \textbf{Thèse générale} du texte en une phrase. \textit{Ex: « Extrait de \emph{Le Gai Savoir} (1882), Nietzsche, philosophe allemand, critique de la morale chrétienne. Thèse : La mort de Dieu libère l'homme pour une création de valeurs nouvelles. »}
    \item \textbf{Analyse structurée (45 min)} : \textbf{Segmentation} (2--3 mouvements logiques, pas des paragraphes). Pour chaque mouvement : \emph{Idée centrale} + \emph{Vocabulaire clé} (concepts techniques) + \emph{Enchaînement logique} (connecteurs : \emph{donc, or, mais, car}). \textbf{Ne pas paraphraser}, \emph{expliciter} le sens philosophique des termes.
    \item \textbf{Discussion critique (25 min)} : \textbf{Validité} : Où l'argument tient-il ? (Ex: La critique nietzschéenne de la « morale d'esclave » est pertinente sur le ressentiment). \textbf{Limites} : Angles morts, présupposés non dits, conséquences extrêmes (Ex: Risque de nihilisme, absence de fondement universel des droits de l'homme). \textbf{Confrontation} : Un autre auteur/école (Ex: Kant / Levinas / Ricoeur).
    \item \textbf{Conclusion synthétique (10 min)} : Reprendre la thèse du texte $\rightarrow$ Situer sa portée (apport à l'histoire de la philosophie) $\rightarrow$ Ouvrir sur une question actuelle (Ex: « Cette « mort de Dieu » annonce-t-elle la crise du sens dans nos sociétés techniciennes ? »).
\end{enumerate}
\end{methode}

\courssub{Exploitation des documents (Sujets à dossier)}

Le Bac Technique propose souvent un sujet \textbf{avec dossier} (2 à 4 documents : textes, graphiques, photos, cartes, stats).

\begin{propriete}[Règles d'or du dossier]
\begin{itemize}
    \item \textbf{Nature} : Identifier chaque doc (Texte philosophique ? Article de presse ? Graphique d'évolution ? Photo de manifestation ? Carte de répartition ?).
    \item \textbf{Croisement} : Ne pas traiter doc 1, puis doc 2. \textbf{Croiser} : « Le document 1 (texte de Locke) définit la tolérance ; le document 3 (graphique) montre que les pays à forte tolérance religieuse ont un IDH plus élevé ; le document 2 (photo) illustre la tension réelle au Nord-Cameroun. »
    \item \textbf{Extraction d'arguments} : Chaque doc doit nourrir \textbf{une partie ou une sous-partie} de votre dissertation. Le doc \emph{est} votre exemple ou votre contre-exemple.
    \item \textbf{Interdiction} : Faire un « commentaire composé » des documents les uns après les autres sans thèse directrice. La dissertation \emph{structure} les documents, pas l'inverse.
\end{itemize}
\end{propriete}

\courssub{Les erreurs fatales au Bac (Liste noire)}

\begin{attention}[Pièges éliminatoires (< 5/20 fréquent)]
\begin{enumerate}
    \item \textbf{Hors-sujet théologique/catéchétique} : Citer la Bible/Coran comme vérité absolue. $\rightarrow$ \textbf{Sanction :} Note éliminatoire. La philosophie examine les \emph{conditions de possibilité} et les \emph{implications} du croire, pas la vérité révélée.
    \item \textbf{Hors-sujet sociologique pur} : Décrire des rites, donner des stats sans analyse conceptuelle. $\rightarrow$ Manque de \emph{philosophicité}.
    \item \textbf{Le « Catalogue d'auteurs »} : « Platon dit ça, Aristote dit l'inverse, Kant tranche... » sans fil conducteur. $\rightarrow$ Absence de \emph{problématisation personnelle}.
    \item \textbf{L'opinion non argumentée} : « Je pense que Dieu existe car je le sens. » $\rightarrow$ Subjectivité non universalisable. Remplacer par : « L'argument de l'expérience religieuse (Schleiermacher) pose le sentiment de dépendance absolue comme fondement... »
    \item \textbf{Plan binaire « Pour / Contre » sans Synthèse} : I. Oui c'est un obstacle. II. Non ce n'est pas un obstacle. III. Conclusion. $\rightarrow$ Plan \emph{analytique} plat, pas dialectique. Il faut : I. Pourquoi ça en a l'air (Thèse). II. Pourquoi c'est plus complexe/le contraire (Antithèse). III. Comment on dépasse l'opposition (Synthèse : distinction des \emph{formes} de religion, des \emph{types} de progrès, du \emph{niveau} d'analyse).
    \item \textbf{Mauvaise gestion du temps} : Introduction de 2 pages, pas de conclusion. $\rightarrow$ S'entrainer au \emph{chrono} : 15 min brouillon (analyse + plan détaillé), 2h15 rédaction, 15 min relecture.
\end{enumerate}
\end{attention}

\courssub{Entraînement guidé : Sujet Zéro}

\textbf{Sujet :} \og La foi est-elle incompatible avec la raison ? \fg

\begin{exoresolu}[Sujet Zéro : La foi est-elle incompatible avec la raison ?]
\textbf{Sujet :} « La foi est-elle incompatible avec la raison ? »
\tcblower
\textbf{1. Concepts clés :} \cle{Foi} (adhésion libre à l'invisible, confiance, don de soi), \cle{Raison} (faculté de connaître par concepts, preuve, universalité), \cle{Incompatibilité} (contradiction logique, exclusion mutuelle).
\textbf{2. Tension :} \emph{Certitude subjective / Engagement existentiel} (Foi) vs \emph{Évidence objective / Démonstration universelle} (Raison).
\textbf{3. Problématique :} « L'adhésion du cœur (foi) exclut-elle nécessairement l'exigence du concept (raison), ou la raison trouve-t-elle dans la foi son horizon de sens, tandis que la foi a besoin de la raison pour se purifier de la superstition ? »
\textbf{4. Plan dialectique :}
\begin{itemize}
    \item \textbf{I. Thèse : L'incompatibilité structurelle (La foi comme sortie de la raison).}
    \begin{itemize}
        \item A. Tertullien / Kierkegaard : \emph{Credo quia absurdum} -- La foi commence où la raison s'arrête (paradoxe, saut qualitatif).
        \item B. Critique rationaliste (Voltaire, Russell) : La foi = croyance sans preuve = superstition / fanatisme. Ex: Refus des vaccins au nom de la foi (Nord Cameroun, Pakistan).
        \item C. La raison moderne (sciences) explique le monde sans hypothèse Dieu (Laplace) : la foi devient \emph{superflue} (Occam).
    \end{itemize}
    \item \textbf{II. Antithèse : L'harmonie possible / La raison au service de la foi (Fides et Ratio).}
    \begin{itemize}
        \item A. Thomas d'Aquin / scolastique : \cle{Préambules de la foi} (existence de Dieu) accessibles à la raison ; \cle{Mystères} (Trinité) au-dessus d'elle, mais non \emph{contraires}. \emph{Gratia non tollit naturam}.
        \item B. Descartes / Kant : La raison \emph{prouve} Dieu (preuve ontologique, morale) $\rightarrow$ La foi rationnelle. Kant : \emph{Religion dans les limites de la simple raison}.
        \item C. Exemple camerounais : Scientifiques croyants (ex: Pr. Mbua, médecins) qui voient dans la complexité cellulaire (ADN) une \emph{intelligibilité} renforçant leur foi. La raison éclaire la foi (théologie naturelle).
    \end{itemize}
    \item \textbf{III. Synthèse : Distinction des régimes / Complémentarité critique (Ricœur, Habermas, Jean-Paul II).}
    \begin{itemize}
        \item A. Distinction des \cle{registres} : La raison = \emph{logos} (explication, causalité, universel) ; La foi = \emph{pistis} (confiance, sens, salut, singularité). Pas même objet $\rightarrow$ pas contradiction logique (Wittgenstein : « formes de vie » différentes).
        \item B. \cle{Purification mutuelle} : La raison purge la foi de l'idolâtrie (magique, fondamentalisme) ; La foi rappelle à la raison ses limites (éthique de la techno-science, sens de la vie, dignité humaine). \emph{Fides quaerens intellectum} (Anselme).
        \item C. Au Cameroun (Laïcité 1996) : L'État de droit (raison juridique) garantit la liberté de culte (foi) ; Les leaders religieux (foi) appellent au respect des lois (raison) -- PCRC, Plateforme des confessions religieuses. \textbf{Incompatibilité ? Non. Tension féconde ? Oui.}
    \end{itemize}
\end{itemize}
\textbf{Introduction type :}
\textbf{Accroche :} « La foi et la raison sont comme les deux ailes qui permettent à l'esprit humain de s'élever vers la vérité » (Jean-Paul II, \emph{Fides et Ratio}, 1998). Cette image aérienne masque un conflit séculaire : Galilée condamné, Darwin contesté, Boko Haram interdisant l'école « occidentale » (\emph{Haram}).
\textbf{Définitions :} La \cle{foi} est une adhésion libre et totale à une vérité révélée, inaccessible à la démonstration (Saint Augustin : \emph{credere est assentire}). La \cle{raison} est la faculté de connaître par des concepts, d'argumenter et de prouver universellement (Descartes : \emph{bon sens}).
\textbf{Problématique :} L'adhésion du cœur (foi) exclut-elle nécessairement l'exigence du concept (raison), ou la raison trouve-t-elle dans la foi son horizon de sens, tandis que la foi a besoin de la raison pour se purifier de la superstition ?
\textbf{Annonce du plan :} Nous analyserons d'abord pourquoi la foi semble structurellement incompatible avec la raison (I), avant de montrer comment les grandes traditions philosophiques et théologiques ont pensé leur harmonie (II), pour enfin distinguer leurs registres respectifs et montrer leur complémentarité critique dans l'espace public camerounais (III).
\end{exoresolu}

\courssub{Outils de révision : Grille d'auto-évaluation}

Utilisez cette grille \emph{avant} de rendre votre copie (ou en vous corrigeant sur annales).

\begin{popfigure}
\centering
\begin{tabularx}{\linewidth}{|>{\bfseries}l|X|X|X|X|}\hline
\rowcolor{popblueL}
\textbf{Critère (Sur 4 pts)} & \textbf{4 -- Excellent} & \textbf{3 -- Bon} & \textbf{2 -- Moyen} & \textbf{1 -- Insuffisant} \\ \hline
\textbf{Compréhension du sujet} & Concepts définis, tension saisie, champ philo respecté & Concepts définis, tension vue, quelques glissements & Concepts flous, tension mal vue, hors-champ partiel & Hors-sujet total, définitions absentes, opinion pure \\ \hline
\textbf{Problématisation / Plan} & Problématique tensionnée, plan dialectique 3 temps clair, annonces précises & Problématique correcte, plan binaire ou analytique mais structuré & Problématique faible (reformulation), plan désorganisé (liste) & Pas de problématique, pas de plan visible \\ \hline
\textbf{Connaissances / Concepts} & 3--5 concepts opératoires maîtrisés, auteurs utilisés comme outils, ex. camerounais pertinents & Concepts présents, auteurs cités correctement, ex. généraux & Concepts confus, auteurs en « name-dropping », ex. vagues/absents & Aucune référence philosophique, connaissances erronées \\ \hline
\textbf{Argumentation / Rigueur} & Enchaînement logique fluide, transitions dialectiques, paragraphes types (Idée/Arg/Ex/Trans) & Argumentation suivie, quelques sauts logiques, transitions faibles & Arguments juxtaposés, assertions non prouvées, répétitions & Incohérence, contradictions, paralogismes \\ \hline
\textbf{Expression / Français} & Style académique, vocabulaire précis, syntaxe impeccable, 0 faute & Correct, quelques lourdeurs, rares fautes & Langage familier/approximatif, fautes récurrentes gênant la lecture & Illisible, charabia, fautes graves \\ \hline
\end{tabularx}
\legende{Grille d'évaluation type Bac Tech (Total /20). Visez 3--4 partout. Le point faible récurrent : Colonne 2 (Problématisation) et Colonne 4 (Argumentation).}
\end{popfigure}

\begin{aretenir}[Check-list finale avant l'examen (J-1)]
\begin{itemize}
    \item ✓ Je sais définir \textbf{5 concepts opératoires} du programme (Laïcité, Sacré, Preuve, Aliénation, Sécularisation, Tolérance, Fondamentalisme, Athéisme, Agnosticisme, Dieu \emph{philosophique} vs \emph{religieux}).
    \item ✓ Je connais \textbf{3 arguments « choc »} par grand thème (Preuves de Dieu : Ontologique, Cosmologique, Physico-théologique, Moral + Critiques Kant/Hume ; Religion/Progrès : Weber vs Marx vs Durkheim ; Laïcité : Séparation vs Reconnaissance vs Gestion).
    \item ✓ J'ai \textbf{3 exemples camerounais datés/précis} (Loi 2004/001, Constitution 1996 Art 1 \& 37, Écoles confessionnelles, Boko Haram / Radicalisation, Code de la famille, PCRC, Pèlerinage Maroua/Mokolo).
    \item ✓ Je maîtrise le \textbf{chrono} : 15' brouillon / 2h15 rédac / 15' relecture (Dissert) ; 10' / 45' / 25' / 10' (Explication).
    \item ✓ Je sais rédiger une \textbf{introduction en entonnoir} et une \textbf{conclusion en entonnoir inversé} sans bafouiller.
\end{itemize}
\end{aretenir}

\begin{savaistu}[Le savais-tu ? -- Statistiques Bac Philo Tech (Session 2023 -- estimations examinateurs)]
\begin{itemize}
    \item \textbf{Moyenne nationale :} $\approx$ 9/20.
    \item \textbf{Taux de réussite dissertation :} $<$ 40\% (note $\ge$ 10/20).
    \item \textbf{Causes majeures d'échec :} 1. Absence de problématique (sujet recopié). 2. Plan binaire « Pour/Contre ». 3. Hors-sujet théologique/sociologique. 4. Conclusion absente ou « En conclusion, on a vu que... ».
    \item \textbf{Levier de progression :} L'explication de texte est souvent mieux réussie (moyenne $\approx$ 11/20) car guidée. \textbf{Mais la dissertation compte double dans le coefficient final.} Travaillez la dissertation \textbf{en priorité} : 1 sujet par semaine en conditions réelles (chronométré).
\end{itemize}
\end{savaistu}

\begin{aretenir}[Astuce Bac Tech : L'ancrage camerounais]
Les correcteurs du Bac Technique valorisent l'intelligence de la \emph{situation locale}. Avoir \textbf{2--3 exemples camerounais précis, datés, analysés philosophiquement} vaut mieux que 10 noms d'auteurs mal compris.
\begin{center}
\begin{tabularx}{\linewidth}{|l|X|}\hline
\rowcolor{popblueL}
\textbf{Thème} & \textbf{Exemple « Clé en main » à maîtriser} \\ \hline
Laïcité & Constitution 1996 (Art. 1 : « Laïque » ; Art. 37 : Liberté conscience/culte) + Loi 2004/001 (Éducation : écoles confessionnelles sous contrat). Tension : Financement public / Autonomie religieuse. \\ \hline
Religion / Conflit & Boko Haram (Extrémisme violent, instrumentalisation de l'islam) vs Réponse militaire + Déradicalisation (Centre de Mérina) + Dialogue interreligieux (PCRC). \emph{Analyser :} Foi $\neq$ Fanatisme ; Laïcité $\neq$ Laïcisme. \\ \hline
Religion / Développement & Réseaux santé/éducation (CBC, Caritas, OIC) : 40\% couverture santé (MINSANTE). \emph{Analyser :} Subsidiarité, Service public, Éthique du care (Levinas/Tronto). \\ \hline
Foi / Raison & Scientifiques croyants camerounais (ex: Académie des Sciences) ; Débat sur « Science sans conscience » (Rabelais) appliqué aux biotechnologies (PMA, GPA, Fin de vie) au Cameroun. \\ \hline
\end{tabularx}
\end{center}
\end{aretenir}

\begin{exercice}[Entraînement individuel -- Sujets types Bac Tech 2024-2025]
Traitez \textbf{un} de ces sujets en conditions d'examen (2h30 chrono, brouillon + copie). Utilisez la grille d'auto-évaluation ci-dessus.
\begin{enumerate}
    \item \textbf{Dissertation :} « La tolérance religieuse est-elle la condition de la paix civile ? »
    \item \textbf{Dissertation :} « La science peut-elle se passer de toute référence à Dieu ? »
    \item \textbf{Explication de texte :} Extrait de \emph{Traité théologico-politique}, Spinoza (Chap. 20) : « Le droit souverain de l'individu s'étend jusqu'où s'étend sa puissance... » (Liberté de philosopher, séparation théologie/philosophie).
    \item \textbf{Sujet à dossier :} Doc 1 : Extrait \emph{Laïcité 1905 / 1996 Cameroun}. Doc 2 : Graphique « Taux de pratique religieuse par région (Cameroun, 2020) ». Doc 3 : Photo « Cérémonie officielle œcuménique Palais de l'Unité ». \textbf{Sujet :} « L'État camerounais face au fait religieux : neutralité ou reconnaissance ? »
\end{enumerate}
\end{exercice}

\begin{corrige}[Exercice n -- Pistes pour Sujet 1 : Tolérance / Paix civile]
\textbf{Problématique :} La tolérance (support de l'autre \emph{malgré} le désaccord) suffit-elle à fonder la paix (ordre juste), ou la paix civile exige-t-elle une \emph{reconnaissance} juridique positive (droits) et une \emph{fraternité} politique au-delà de la simple tolérance (condescendance) ?
\textbf{Plan :} I. Tolérance comme condition \emph{nécessaire} (Arrêt guerres de religion, Locke, Voltaire, Laïcité 1996). II. Tolérance comme condition \emph{insuffisante} (Tolérance « négative », instable, pouvoir révocable -- Richter, Walzer ; Ex: Crises Nord-Cameroun, tolérance précaire). III. Vers une \emph{concordance} civile : Reconnaissance mutuelle (Honneth), Citoyenneté constitutionnelle (Droit commun), Dialogue interconvictionnel (PCRC) -- La paix comme construction politique active, pas seulement absence de guerre.
\end{corrige}

\newpage

\coursec{Activité d'intégration}

\courssub{Objectifs de l'activité}

Cette activité d'intégration prend la forme d'un \cle{débat réglementé} sur le thème : « La laïcité camerounaise face au port du voile intégral (niqab/burqa) dans l'espace public et scolaire ». Elle vise à mobiliser, dans une situation concrète, l'ensemble des acquis de la leçon sur Dieu et la religion.

À l'issue de cette activité, l'élève doit être capable de :
\begin{itemize}
\item \cle{Définir} et articuler les concepts de \cle{laïcité}, de \cle{liberté de conscience}, de \cle{neutralité} de l'espace public et de \cle{vivre-ensemble} ;
\item \cle{Appliquer} les notions de l'unité (conceptions de Dieu, valeur sociale de la religion, rapport foi/raison) à une situation concrète de la vie courante camerounaise ;
\item \cle{Construire} une argumentation structurée (thèse, arguments, exemples, réponse aux objections) conformément à la méthode de la dissertation ;
\item \cle{Distinguer} l'argument rationnel du sophisme (fallacie) et évaluer la pertinence des arguments adverses ;
\item \cle{Rédiger et justifier} une position personnelle nuancée, ni laxiste, ni liberticide.
\end{itemize}

\courssub{Situation}

\textbf{Communiqué de la communauté urbaine de Garoua — Mars 2026.}

Depuis la rentrée scolaire, plusieurs cas ont été signalés dans des établissements publics de la région du Nord : des élèves de classes de Première et de Terminale se présentent en cours le visage entièrement couvert par le niqab. Dans un lycée de Garoua, trois élèves ont refusé de retirer le voile intégral pour l'épreuve pratique du baccalauréat technique, rendant impossible leur identification par le jury. Le proviseur, invoquant la circulaire du \textbf{MINESEC} de 2018 relative aux tenues scolaires et au déroulement des examens, les a exclues de la salle d'examen. Les parents ont saisi le Conseil Islamique local, qui dénonce une atteinte à la liberté religieuse garantie par la Constitution.

Le débat enflamme les réseaux sociaux. Un sondage informel mené par une radio locale auprès de 500 auditeurs donne les résultats suivants :

\begin{center}
\begin{tabularx}{\linewidth}{|l|X|c|}\hline
\rowcolor{popblueL} \textbf{Position} & \textbf{Justification dominante} & \textbf{Part} \\ \hline
Interdiction totale dans l'espace public & Sécurité, identification, vivre-ensemble & 38 \% \\ \hline
Interdiction à l'école publique seulement & Neutralité du service public, protection des mineures & 34 \% \\ \hline
Liberté totale & Liberté de conscience, non-discrimination & 22 \% \\ \hline
Sans opinion & --- & 6 \% \\ \hline
\end{tabularx}
\end{center}

Le contexte est sensible : la région a connu entre 2014 et 2016 des attentats-suicides commis par des femmes dissimulées sous le voile intégral, ce qui avait conduit les autorités de l'Extrême-Nord à interdire temporairement son port par arrêté préfectoral. Mais les juristes rappellent que le Cameroun n'a pas de loi générale d'interdiction comparable à la loi française de 2010, validée par la Cour européenne des droits de l'homme dans l'arrêt \emph{S.A.S. c. France} (2014).

Le préfet, soucieux d'apaiser la tension, décide d'organiser une \textbf{concertation citoyenne} au lycée. Votre classe de Terminale technique y est invitée et devra produire, à l'issue d'un débat réglementé, un \emph{avis philosophique motivé} remis aux autorités.

\courssub{Consignes}

La classe est divisée en quatre groupes. Les groupes 1 à 3 préparent chacun une \textbf{plaidoirie de 5 minutes} structurée ainsi : thèse claire, deux arguments philosophiques, un exemple camerounais concret, et une réponse à l'objection principale adverse. Le groupe 4 constitue le \textbf{jury philosophique}.

\begin{enumerate}
\item \textbf{Groupe 1 — « Interdiction totale ».} Défendez l'interdiction du voile intégral dans tout l'espace public. Mobilisez les arguments de la sécurité, de la laïcité stricte et de l'égalité femmes/hommes. Votre plaidoirie doit répondre à l'objection : « Vous violez la liberté de conscience garantie par la Constitution ».
\item \textbf{Groupe 2 — « Interdiction à l'école publique seulement ».} Défendez une interdiction limitée aux services publics (écoles, administrations, hôpitaux). Mobilisez la neutralité du service public et la protection des mineures. Répondez à l'objection : « Une interdiction partielle est incohérente : soit on interdit partout, soit nulle part ».
\item \textbf{Groupe 3 — « Liberté totale ».} Défendez le droit de porter le voile intégral partout. Mobilisez la liberté de conscience, la non-discrimination et l'autonomie du choix individuel. Répondez à l'objection : « Le voile intégral est un instrument d'oppression de la femme et un danger sécuritaire ».
\item \textbf{Groupe 4 — Jury philosophique.} Pendant les plaidoiries, complétez la grille d'évaluation : chaque argument est-il défini, étayé, pertinent ? Relevez les sophismes (généralisation abusive, appel à la peur, argument d'autorité, attaque personnelle). Rédigez ensuite un \emph{avis consultatif motivé} de 15 lignes maximum.
\item \textbf{Questions de synthèse (tous les groupes).} 
\begin{enumerate}
\item La laïcité camerounaise est-elle identique à la laïcité française ? Justifiez à partir du préambule de la Constitution.
\item Peut-on limiter une liberté fondamentale ? À quelles conditions (légalité, nécessité, proportionnalité) ?
\item Le port du voile intégral relève-t-il d'un choix autonome ou d'une pression sociale ? Que répondre à cette alternative ?
\item Formulez en une phrase votre position personnelle nuancée, en distinguant le plan du \cle{droit} et le plan de l'\cle{éthique}.
\end{enumerate}
\end{enumerate}

\courssub{Ressources à mobiliser}

\begin{aretenir}[Concepts et repères de la leçon]
\begin{itemize}
\item \textbf{Laïcité} : principe de séparation des institutions religieuses et de l'État, garantissant la neutralité de l'espace public et la liberté de conscience de chacun.
\item \textbf{Liberté de conscience} : droit de croire ou de ne pas croire, de pratiquer ou non une religion ; elle est absolue dans le for intérieur, mais ses \emph{manifestations} peuvent être limitées par la loi pour l'ordre public.
\item \textbf{Neutralité} : les agents et les services publics ne manifestent pas d'appartenance religieuse, afin de traiter tous les usagers également.
\item \textbf{Valeur sociale de la religion} (Durkheim) : la religion est un ciment social, mais elle peut aussi entrer en tension avec les exigences du vivre-ensemble moderne.
\item \textbf{Foi et raison} : la croyance religieuse relève de la conviction personnelle ; la règle commune, dans l'espace public, relève de la discussion rationnelle.
\item \textbf{Principe de proportionnalité} : toute restriction d'une liberté doit être prévue par la loi, nécessaire à un but légitime (sécurité, égalité) et adaptée à ce but.
\end{itemize}
\end{aretenir}

\begin{attention}[Sophismes à éviter]
\begin{itemize}
\item \emph{Généralisation abusive} : « Toutes les femmes voilées sont opprimées » / « Tous les partisans de l'interdiction sont islamophobes ».
\item \emph{Appel à la peur} : « Si on autorise le niqab, ce sera le chaos ».
\item \emph{Argument d'autorité} : « Un chef religieux l'a dit, donc c'est vrai ».
\item \emph{Faux dilemme} : « Soit la liberté totale, soit la dictature ».
\end{itemize}
\end{attention}

\courssub{Démarche et corrigé commenté}

\begin{exoresolu}[Construction d'une position philosophique nuancée]
\textbf{Énoncé.} À partir du débat, rédigez la position de synthèse attendue : faut-il interdire le voile intégral dans l'espace public et scolaire camerounais ?

\tcblower
\textbf{Étape 1 — Clarifier les termes du problème.} Le débat ne porte pas sur l'existence de Dieu ni sur la vérité de l'islam, mais sur une \emph{manifestation extérieure} d'une croyance. Il faut distinguer trois plans : le plan \textbf{religieux} (le niqab est-il une prescription ? — les avis religieux divergent), le plan \textbf{juridique} (que dit le droit camerounais ?) et le plan \textbf{philosophique} (qu'est-ce qui est juste ?). Confondre ces plans est la première erreur du débat.

\textbf{Étape 2 — Peser les arguments des trois thèses.}
\begin{itemize}
\item L'\emph{interdiction totale} protège la sécurité et l'identification, mais elle est \textbf{disproportionnée} : elle pénalise toutes les femmes concernées, y compris dans la rue, là où aucune nécessité d'identification ne s'impose. Elle risque de stigmatiser une communauté et de violer la liberté de conscience.
\item La \emph{liberté totale} respecte l'autonomie individuelle, mais elle ignore deux contraintes réelles : l'impossibilité matérielle d'identifier une candidate à un examen (ce qui menace l'égalité de traitement et la fraude) et le contexte sécuritaire de l'Extrême-Nord.
\item L'\emph{interdiction limitée aux services publics} apparaît comme la position la plus conforme au principe de proportionnalité : elle restreint la liberté seulement là où un but légitime (identification, neutralité, protection des mineures) l'exige.
\end{itemize}

\textbf{Étape 3 — Appliquer le principe de proportionnalité.} Une restriction est légitime si elle remplit trois conditions : (a) elle est \emph{prévue par la loi} — or le Cameroun ne dispose que de circulaires et d'arrêtés locaux, d'où la nécessité d'un cadre légal clair ; (b) elle poursuit un \emph{but légitime} — sécurité, égalité des candidats, neutralité scolaire ; (c) elle est \emph{adaptée} — interdire le niqab à l'école et lors des contrôles d'identité suffit, sans interdiction générale dans la rue.

\textbf{Étape 4 — Distinguer droit et éthique.} Sur le plan du droit, l'État peut encadrer les manifestations religieuses dans les services publics. Sur le plan éthique, la société doit s'interroger : le port du voile intégral est-il un choix libre ? La réponse exige le dialogue, l'éducation et la lutte contre les pressions sociales — non l'humiliation des femmes concernées, qui seraient doublement victimes.

\textbf{Conclusion de synthèse.} Une position philosophique nuancée consiste à défendre : la liberté de conscience comme principe, l'interdiction du voile intégral dans les services publics et lors des examens comme restriction proportionnée, et le refus de toute stigmatisation comme exigence du vivre-ensemble. Cette position n'est ni laxiste (elle pose des limites claires) ni liberticide (elle préserve la sphère privée).
\end{exoresolu}

\courssub{Bilan de l'activité}

Ce débat réglementé a permis de mettre en évidence plusieurs acquis essentiels. D'abord, les notions abstraites de la leçon — laïcité, liberté de conscience, valeur sociale de la religion — prennent un sens concret dès qu'on les confronte à une situation réelle de la vie scolaire camerounaise : un concept philosophique n'est pas une formule à réciter, mais un outil pour penser un conflit. Ensuite, l'activité a montré qu'une bonne argumentation exige de \emph{définir ses termes}, de \emph{hiérarchiser ses arguments} et de \emph{répondre aux objections}, compétences directement transférables à la dissertation du baccalauréat technique. Le jury a en outre appris à débusquer les sophismes qui encombrent les débats publics, notamment sur les réseaux sociaux. Enfin, la synthèse a révélé l'exigence proprement philosophique de la \cle{nuance} : entre la répression générale et la liberté sans limites, la raison trace une voie proportionnée, qui respecte à la fois l'ordre public et la dignité des personnes. C'est là le cœur du métier de citoyen dans un État de droit pluraliste comme le Cameroun.

\newpage

\coursec{Méthodes et exercices résolus}

\courssub{Méthodes et exercices résolus : Dieu et la Religion}

\begin{methode}[Analyser l'origine du sentiment religieux dans une situation concrète]
\textbf{Objectif :} Identifier et distinguer les facteurs psychologiques (peur, merveilleux) et sociologiques (projection, cohésion) à l'œuvre dans un fait religieux donné.
\begin{enumerate}
    \item \textbf{Repérer le fait religieux :} Citer l'extrait, l'observation ou la situation (rite, prière, croyance, institution).
    \item \textbf{Isoler les mobiles subjectifs :} Chercher la \cle{peur} (de la mort, du hasard, de la nature), le \cle{merveilleux} (émerveillement face à l'ordre du monde, sentiment de dépendance), le \cle{besoin de sens} (réponse à l'absurde).
    \item \textbf{Identifier la fonction sociale (Durkheim, Freud) :} La religion fait-elle \cle{lien social} (cohésion, identité collective) ? Est-elle \cle{projection} des idéaux humains (Feuerbach) ou \cle{névrose obsessionnelle} (Freud) ?
    \item \textbf{Synthétiser :} Rédiger un paragraphe structuré : « Dans cette situation, le sentiment religieux prend sa source dans \ldots (facteur psychologique) et remplit la fonction de \ldots (facteur sociologique). »
\end{enumerate}
\textbf{Reflexe Bac :} Ne jamais réduire le religieux à \emph{une seule} cause. Le programme exige la confrontation des thèses (psychologiques vs sociologiques).
\end{methode}

\begin{exoresolu}[Analyse d'une situation : La procession de la fête du Ngondo à Douala]
\textbf{Énoncé :} « Chaque année, les peuples Sawa se réunissent sur les rives du Wouri pour le Ngondo. Un vase est immergé par un initié pour aller chercher le message des ancêtres au fond de l'eau. La foule, dans le silence puis la liesse, attend le retour du message qui scelle l'unité du peuple et assure la prospérité. »
\textbf{Consigne :} En vous appuyant sur les thèses étudiées (peur, merveilleux, projection sociale), analysez l'origine et la fonction du sentiment religieux manifesté dans cette cérémonie.

\tcblower
\textbf{Solution détaillée :}

\textbf{1. Repérage du fait religieux :} La situation décrit un \cle{rite} collectif (immersion du vase, attente du message), porté par une \cle{communauté} (peuples Sawa), visant une \cle{communication avec le sacré} (ancêtres, forces de l'eau) pour obtenir un résultat temporel (prospérité, unité).

\textbf{2. Analyse des mobiles psychologiques (Origine subjective) :}
\begin{itemize}
    \item \textbf{La peur / l'angoisse :} L'immersion dans l'eau profonde symbolise le danger, la mort, l'inconnu. L'initié risque sa vie. La foule retient son souffle. Le rite répond à l'angoisse face aux forces naturelles incontrôlables (le fleuve, les inondations) et à l'incertitude de l'avenir (prospérité). \emph{(Thèse de la peur : Lucrèce, Hobbes, Freud « angoisse face à la nature »)}.
    \item \textbf{Le merveilleux / le sentiment de dépendance :} Le silence, l'attente, la liesse finale traduisent un \cle{sentiment océanique} (Romain Rolland / Freud) ou le \cle{sentiment de dépendance absolue} (Schleiermacher). L'homme se sait contingent face à l'immensité du cosmos (le Wouri). L'émerveillement devant l'ordre du monde (le cycle des saisons, la vie) bascule dans le sacré.
\end{itemize}

\textbf{3. Analyse de la fonction sociologique (Origine objective / Fonction) :}
\begin{itemize}
    \item \textbf{Projection sociale (Feuerbach / Durkheim) :} Le « message des ancêtres » est une \cle{projection} des valeurs suprêmes de la société Sawa : l'unité, la paix, la fécondité. Dieu/Ancêtres = Société divinisée. « L'essence de la religion est l'essence de l'homme projetée hors de lui » (Feuerbach).
    \item \textbf{Cohésion et identité collective (Durkheim) :} Le rite crée l'\cle{effervescence collective}. Le silence partagé puis la liesse commune génèrent la \cle{conscience collective}. Le Ngondo \emph{fait} le peuple Sawa ; sans ce rite, l'unité se déliterait. La religion est « un système solidaire de croyances et de pratiques relatives à des choses sacrées [...] qui unissent en une même communauté morale appelée Église, tous ceux qui y adhèrent » (Durkheim, \emph{Les Formes élémentaires de la vie religieuse}).
    \item \textbf{Consolation / Illusion (Freud / Marx) :} Le message assure la prospérité future. C'est une \cle{consolation} face à la précarité économique (Marx : « l'opium du peuple », soupir de la créature opprimée) et une \cle{réponse illusoire} aux angoisses insolubles (Freud : « névrose obsessionnelle de l'humanité »).
\end{itemize}

\textbf{4. Synthèse rédactionnelle (Modèle de réponse attendue) :}
« Dans la cérémonie du Ngondo, le sentiment religieux trouve une double origine. \textbf{Subjectivement}, il naît de la \cle{peur} face à la puissance indomptée du fleuve Wouri et de la mort (l'immersion de l'initié), ainsi que du \cle{merveilleux} ou sentiment de dépendance absolue face à l'immensité cosmique. \textbf{Objectivement}, il remplit une fonction \cle{sociologique} majeure : selon Durkheim, le rite génère l'effervescence collective qui soude l'identité Sawa (cohésion) ; selon Feuerbach, le « message des ancêtres » est la projection des idéaux communautaires (unité, prospérité) ; selon Freud, il apaise l'angoisse par une illusion consolatrice. Le religieux est ainsi à la fois \emph{expérience intérieure} et \emph{fait social total}. »

\textbf{Conclusion :} L'analyse montre qu'aucune thèse unique n'épuise le phénomène : la peur et le merveilleux en sont le moteur psychique ; la cohésion sociale et la projection des idéaux en sont la fonction structurelle.
\end{exoresolu}

\begin{methode}[Cartographier les conceptions de Dieu : Classer et distinguer]
\textbf{Objectif :} Situer une citation, une croyance ou un auteur dans la topologie conceptuelle : Théisme, Déisme, Panthéisme, Athéisme, Agnosticisme.
\begin{enumerate}
    \item \textbf{Identifier les deux critères discriminants :}
    \begin{itemize}
        \item \textbf{Transcendance / Immanence :} Dieu est-il \emph{hors} du monde (créateur distinct) ou \emph{dans} le monde (identique à la nature) ?
        \item \textbf{Personnalité / Providencialité :} Dieu est-il une \cle{personne} (volonté, intelligence, amour, dialogue) ou une \cle{force impersonnelle} (ordre nécessaire, substance, destin) ? Intervient-il dans l'histoire (providence, miracles, révélation) ?
    \end{itemize}
    \item \textbf{Appliquer le tableau de décision :}
    \begin{center}
    \begin{tabularx}{\linewidth}{|l|X|X|}\hline
    \textbf{Conception} & \textbf{Transcendance / Immanence} & \textbf{Personnalité / Providence} \\ \hline
    \textbf{Théisme} (Judaïsme, Christianisme, Islam) & Transcendant (Créateur \emph{ex nihilo}) & \textbf{Personnel} (Père, Juge, Révélateur, Providence) \\ \hline
    \textbf{Déisme} (Voltaire, Horloger) & Transcendant (Créateur/Architecte) & \textbf{Impersonnel} (Horloger absent, pas de miracle, pas de révélation, religion naturelle) \\ \hline
    \textbf{Panthéisme} (Spinoza, Stoïciens, Hindouisme advaita) & \textbf{Immanent} (Dieu = Nature, \emph{Deus sive Natura}) & \textbf{Impersonnel} (Substance, Nécessité, pas de volonté libre, pas de providence) \\ \hline
    \textbf{Athéisme} (Nietzsche, Marx, Sartre, Matérialistes) & \textbf{Négation} de l'existence de Dieu & --- \\ \hline
    \textbf{Agnosticisme} (Kant, Huxley, Russell) & \textbf{Suspension} du jugement (Inconnaissable) & --- \\ \hline
    \end{tabularx}
    \end{center}
    \item \textbf{Rédiger la justification :} « Cette conception relève du [Théisme/Déisme...] car elle affirme que Dieu est [transcendant/immanent] et [personnel/impersonnel], ce qui se manifeste par [citation/mot-clé : révélation, miracle, loi naturelle, substance]. »
\end{enumerate}
\textbf{Formule à retenir :} \cle{Théisme = Transcendance + Personnalité} ; \cle{Déisme = Transcendance - Personnalité} ; \cle{Panthéisme = Immanence - Personnalité}.
\end{methode}

\begin{exoresolu}[Classification de conceptions de Dieu]
\textbf{Énoncé :} Rattachez chacune des propositions suivantes à l'une des cinq conceptions majeures (Théisme, Déisme, Panthéisme, Athéisme, Agnosticisme) et justifiez en une phrase en mobilisant les critères de transcendance/immanence et de personnalité/impersonnalité.
\begin{enumerate}
    \item « Dieu n'est pas un être parmi les êtres, il est l'Être même, la substance unique dont tout est mode. » (Spinoza, \emph{Éthique})
    \item « L'hypothèse Dieu est devenue inutile pour expliquer l'univers. La science suffit. » (Laplace / Athéisme scientifique contemporain)
    \item « Je ne sais pas si Dieu existe, et je pense qu'il est impossible de le savoir. La raison humaine bute sur ses limites. » (Témoignage agnostique)
    \item « Le monde est une machine parfaite créée par un Grand Architecte qui l'a mise en mouvement et ne s'en mêle plus. » (Voltaire, \emph{Dictionnaire philosophique})
    \item « Le Seigneur est mon berger, je ne manquerai de rien... Tu es avec moi. » (Psaume 23, Bible)
\end{enumerate}

\tcblower
\textbf{Solution détaillée :}

\begin{enumerate}
    \item \textbf{Panthéisme (Spinoza).} \textbf{Justification :} Dieu est \textbf{immanent} (« Dieu n'est pas un être parmi les êtres », Il \emph{est} le monde/nature) et \textbf{impersonnel} (Substance unique, Nécessité géométrique, pas de volonté libre, pas de « Tu » personnel, pas de providence).
    \item \textbf{Athéisme.} \textbf{Justification :} \textbf{Négation} pure et simple de l'existence de Dieu (l'hypothèse est « inutile »), sans substitut sacré. La matière/science s'autosuffit.
    \item \textbf{Agnosticisme.} \textbf{Justification :} \textbf{Suspension du jugement} épistémologique (« impossible de le savoir »). Ni affirmation (théisme), ni négation (athéisme). La raison atteint ses limites (Kant : Dieu est une idée régulatrice, non un objet de connaissance).
    \item \textbf{Déisme (Voltaire).} \textbf{Justification :} Dieu est \textbf{transcendant} (Créateur/Architecte extérieur à la machine-monde) mais \textbf{impersonnel} (pas de « Tu », pas de miracle, pas de révélation, pas de providence quotidienne : « ne s'en mêle plus »). Religion naturelle, morale universelle.
    \item \textbf{Théisme (Biblique / Judaïsme-Christianisme).} \textbf{Justification :} Dieu est \textbf{transcendant} (Seigneur, Créateur distinct de sa créature) et \textbf{personnel} (« Tu es avec moi », relation d'alliance, berger, providence aimante, réponse à la prière, révélation).
\end{enumerate}

\textbf{Conclusion :} La distinction opératoire repose sur le couple \cle{Transcendance/Immanence} (Où est Dieu ?) et \cle{Personnalité/Impersonnalité} (Comment se rapporte-t-Il à l'homme ?). Le théisme combine les deux pôles « relationnels » ; le panthéisme les nie tous deux ; le déisme ne garde que la transcendance.
\end{exoresolu}

\begin{methode}[Démontrer / Critiquer les preuves classiques de l'existence de Dieu]
\textbf{Objectif :} Reconstituer la structure logique d'une preuve (a priori / a posteriori) et appliquer la critique kantienne (ou humienne) pertinente.
\begin{enumerate}
    \item \textbf{Identifier le type de preuve :}
    \begin{itemize}
        \item \textbf{A priori (Déductive) :} \cle{Preuve ontologique} (Anselme, Descartes, Leibniz, Gödel). Parte du \emph{concept} de Dieu (Être parfait, Être suprême) pour déduire son \emph{existence}. « Dieu existe parce que l'existence fait partie de son essence. »
        \item \textbf{A posteriori (Inductive / Causale) :}
        \begin{itemize}
            \item \cle{Preuve cosmologique} (Aristote, Thomas d'Aquin, Leibniz) : Parte du \emph{fait} monde (mouvement, causalité, contingence) $\rightarrow$ Cause première / Moteur immobile / Raison suffisante.
            \item \cle{Preuve physico-théologique / Téléologique} (Stoïciens, Thomas d'Aquin 5e voie, Paley, Derrida) : Parte de l'\emph{ordre/finalité} du monde (horloge, œil, constants physiques) $\rightarrow$ Intelligence ordonnatrice (Grand Architecte).
            \item \cle{Preuve morale} (Kant pratique) : Parte du \emph{devoir moral} (souverain bien = vertu + bonheur) $\rightarrow$ Postule Dieu comme garant de la convergence.
        \end{itemize}
    \end{itemize}
    \item \textbf{Appliquer la critique majeure (Kant, \emph{Critique de la raison pure}) :}
    \begin{itemize}
        \item \textbf{Contre l'ontologique :} « L'existence n'est pas un prédicat réel » (ne rajoute rien au concept). 100 thalers réels = 100 thalers possibles (même concept). On ne peut passer du \emph{logique} (concept) au \emph{réel} (existence) par analyse.
        \item \textbf{Contre la cosmologique :} Elle dépend de l'ontologique (pour passer de « cause nécessaire » à « Dieu = être nécessaire »). Elle applique la catégorie de causalité \emph{au-delà} de l'expérience possible (antinomie de la raison).
        \item \textbf{Contre la physico-théologique :} Elle ne prouve qu'un \emph{Architecte} (demiurge), pas un \emph{Créateur ex nihilo}. L'ordre peut s'expliquer par la nature (évolution, auto-organisation). Analogie faible (monde $\neq$ machine).
    \end{itemize}
    \item \textbf{Conclure :} Les preuves théoriques échouent à établir l'existence comme \emph{connaissance scientifique}. La foi relève de l'\cle{intérêt pratique} (Kant) ou du \cle{pari} (Pascal) ou du \cle{saut subjectif} (Kierkegaard).
\end{enumerate}
\textbf{Reflexe Bac :} Ne pas dire « La preuve de Thomas d'Aquin est fausse ». Dire : « Elle est logiquement valide \emph{si} on accepte ses prémisses (causalité universelle, impossibilité de la régression à l'infini), mais Kant montre qu'elle ne constitue pas une \emph{connaissance théorique} car elle sort de l'usage légitime de l'entendement. »
\end{methode}

\begin{exoresolu}[Analyse critique de la preuve ontologique cartésienne]
\textbf{Énoncé :} « De ce que je ne puis penser une montagne sans une vallée, il ne s'ensuit pas qu'il y ait une montagne ou une vallée en la terre ; mais seulement que la montagne et la vallée, soit qu'elles soient, soit qu'elles ne soient pas, ne peuvent l'une être sans l'autre. Mais de ce que je pense Dieu, il s'ensuit nécessairement qu'il existe. » (Descartes, \emph{Méditations métaphysiques}, V).
\textbf{Consigne :} 1. Restituez le raisonnement de Descartes. 2. Expliquez la distinction qu'il opère avec la montagne/la vallée. 3. Présentez et évaluez la critique de Kant (« L'existence n'est pas un prédicat réel »).

\tcblower
\textbf{Solution détaillée :}

\textbf{1. Restitution du raisonnement cartésien (Preuve ontologique) :}
\begin{itemize}
    \item \textbf{Prémisse 1 (Intuition intellectuelle) :} Je possède en moi l'idée claire et distincte d'un \cle{Être souverainement parfait} (Dieu).
    \item \textbf{Prémisse 2 (Lien essence/existence) :} L'\cle{existence} est une \cle{perfection}. Un être parfait qui \emph{manquerait} d'existence serait imparfait (contradiction).
    \item \textbf{Conclusion :} Donc, l'existence \textbf{appartient nécessairement à l'essence} de Dieu. Je ne peux \emph{pas} concevoir Dieu \emph{sans} existence, comme je conçois un triangle sans ses trois angles. L'existence de Dieu est \cle{nécessaire} (déduite \emph{a priori} du concept).
\end{itemize}

\textbf{2. La distinction Montagne/Vallée (Réponse à l'objection de Gassendi/Caterus) :}
\begin{itemize}
    \item \textbf{Montagne/Vallée :} Relation \textbf{conditionnelle / hypothétique}. « \emph{Si} montagne existe, \emph{alors} vallée existe ». L'essence ne \emph{contient} pas l'existence. C'est une liaison entre deux concepts \emph{si} l'objet est posé comme réel.
    \item \textbf{Dieu :} Relation \textbf{nécessaire / inconditionnelle}. « Dieu existe » n'est pas « \emph{Si} Dieu existe, \emph{alors} Il existe ». L'existence est \emph{incluse} dans le concept même (analytique). Je \emph{ne puis pas} penser Dieu sans existence sans contradiction.
    \item \textbf{Portée :} Descartes veut montrer que pour Dieu \emph{seul}, l'essence \emph{implique} l'existence (Dieu = \emph{Ens realissimum}, seul être dont l'essence est d'exister).
\end{itemize}

\textbf{3. La critique de Kant (\emph{Critique de la raison pure}, Dialectique, Idéal de la raison pure) :}
\begin{itemize}
    \item \textbf{Thèse centrale :} « \cle{L'existence n'est pas un prédicat réel} » (un \emph{prédicat réel} ajoute une détermination au concept : « rouge », « lourd », « triangle »).
    \item \textbf{Argument :} Dire « Dieu est omnipotent, omniscient, \emph{existant} » n'ajoute \emph{aucune} détermination nouvelle au concept de Dieu par rapport à « Dieu est omnipotent, omniscient ». 100 thalers \emph{réels} contiennent \emph{exactement la même définition} (concept) que 100 thalers \emph{possibles}. La différence n'est pas dans le \emph{concept} (logique), mais dans la \emph{position} (réel) : l'objet est \emph{donné} dans l'intuition.
    \item \textbf{Conséquence :} Le passage du \emph{concept} (logique) à l'\emph{existence} (réel) ne peut se faire par \emph{analyse} (déduction ontologique). Il faut une \emph{intuition} (expérience) pour affirmer l'existence. Or Dieu n'est pas objet d'expérience possible.
    \item \textbf{Évaluation :} La critique est \textbf{décisive contre la preuve comme \emph{connaissance théorique}}. Elle ne prouve pas que Dieu \emph{n'existe pas}, mais que son existence ne peut être \emph{démontrée par la raison pure}. Descartes confond \emph{nécessité logique} (liaison de concepts) et \emph{nécessité ontologique} (réalité de l'être).
\end{itemize}

\textbf{Conclusion :} La preuve ontologique échoue comme démonstration scientifique. Elle reste toutefois pertinente comme \emph{éclaircissement du concept de Dieu} (Dieu \emph{ne peut pas ne pas exister} \emph{si} Il est pensable comme Être nécessaire), ce qui rejoint la distinction modale (contingent/nécessaire) mais ne vaut pas preuve d'existence réelle.
\end{exoresolu}

\begin{methode}[Définir la religion et analyser sa valeur sociale (Durkheim, Marx, Freud, Weber)]
\textbf{Objectif :} Construire une définition opératoire (croyance + rite + communauté) et évaluer la fonction sociale (cohésion / aliénation / consolation / éthique).
\begin{enumerate}
    \item \textbf{La définition minimale (Durkheim) :} « Un système solidaire de \textbf{croyances} (représentations sacrées) et de \textbf{pratiques} (rites) relatives à des choses \textbf{sacrées} (séparées, interdites), qui \textbf{unissent en une même communauté morale (Église)} tous ceux qui y adhèrent. »
    \item \textbf{Distinction sacré/profane :} C'est le \cle{critère spécifique} du religieux (Durkheim). Le sacré = mis à part, protégé par des interdits ; le profane = quotidien, utilitaire.
    \item \textbf{Grille d'analyse de la valeur sociale (Tableau synoptique) :}
    \begin{center}
    \begin{tabularx}{\linewidth}{|l|X|X|}\hline
    \textbf{Auteur / Thèse} & \textbf{Valeur positive (Fonction)} & \textbf{Valeur négative (Critique / Limite)} \\ \hline
    \textbf{Durkheim} (Fonctionnaliste) & \cle{Cohésion sociale}, identité collective, contrôle social, sacralisation des valeurs suprêmes (la société se divinise elle-même). & Risque de \cle{totémisme} : la société s'adorant elle-même, rigidité, exclusion des déviants. \\ \hline
    \textbf{Marx} (Critique matérialiste) & \cle{Consolation} : « Soupir de la créature opprimée », « Cœur d'un monde sans cœur », « Opium du peuple » (supporte la souffrance). & \cle{Aliénation} : Projection des forces humaines hors de soi $\rightarrow$ soumission. Masque l'exploitation réelle. « La critique de la religion est la condition de toute critique. » \\ \hline
    \textbf{Freud} (Psychanalyse) & \cle{Protection} : « Névrose obsessionnelle universelle » qui protège de l'angoisse (nature, mort, pulsions). Structure le Surmoi (morale). & \cle{Illusion} : « Illusion » (réalisation de désirs), infantile. Entrave à la maturité (autonomie, science). « L'avenir d'une illusion ». \\ \hline
    \textbf{Weber} (Sociologie compréhensive) & \cle{Moteur historique} : « Éthique protestante » $\rightarrow$ capitalisme (ascétisme intramondain). Rationalisation de la conduite de vie. & \cle{Désenchantement du monde} : La rationalisation religieuse conduit à la sécularisation et à la « cage d'acier » bureaucratique. \\ \hline
    \end{tabularx}
    \end{center}
    \item \textbf{Rédiger une évaluation nuancée :} « La religion a une valeur \emph{fonctionnelle} indéniable (lien, sens, éthique, consolation) mais porte en elle des risques d'\emph{aliénation} (soumission, illusion, fanatisme) que la raison critique doit surveiller. »
\end{enumerate}
\textbf{Formule à retenir :} \cle{Religion = Croyances (Sacré) + Rites (Pratiques) + Communauté (Église)}. \cle{Valeur = Fonction (Durkheim) vs Aliénation (Marx/Freud) + Rationalisation (Weber)}.
\end{methode}

\begin{exoresolu}[Dissertation : « La religion est-elle un facteur de cohésion ou de division pour la société ? »]
\textbf{Énoncé :} Sujet de dissertation type Bac. « La religion est-elle un facteur de cohésion ou de division pour la société ? »

\tcblower
\textbf{Solution détaillée : Plan dialectique structuré}

\textbf{Introduction :}
\begin{itemize}
    \item \textbf{Amorce :} Depuis l'origine des sociétés humaines, le fait religieux structure le temps (fêtes), l'espace (lieux saints) et le lien social (rituels). Au Cameroun, la coexistence de l'islam, du christianisme, des religions traditionnelles et de la laïcité d'État en fait un laboratoire vivant.
    \item \textbf{Définition :} Religion = système de croyances/rites unissant une communauté morale autour du sacré (Durkheim).
    \item \textbf{Problématique :} Si le religieux \emph{fait lien} (cohésion interne), il \emph{fait aussi frontière} (division externe/interne). La religion unit-elle \emph{forcément} ou divise-t-elle \emph{structurellement} ?
    \item \textbf{Annonce du plan :} 1) La religion comme ciment social (cohésion). 2) La religion comme source de fractures (division). 3) Dépassement : Laïcité et dialogue interreligieux comme conditions d'une cohésion \emph{inclusive}.
\end{itemize}

\textbf{Développement :}

\textbf{I. La religion, facteur de cohésion sociale (Thèse : Durkheim, Fonctionnalisme)}
\begin{itemize}
    \item \textbf{A. Le sacré unit la communauté :} Le rite (effervescence collective) génère la \cle{conscience collective}. Ex : Ngondo (Sawa), Ngoun (Bamoun), Pèlerinage de Maroua (Musulmans), Fête de l'Assomption (Marianne). Le sacré = \cle{symbolisation de la société elle-même} (Durkheim). « Dieu, c'est la société divinisée. »
    \item \textbf{B. Régulation morale et contrôle social :} Les commandements religieux (Décalogue, Charia, interdits ancestraux) intériorisent les normes. « Tu ne tueras point » protège la vie sociale. La religion légitime l'ordre politique (Droit divin, Mandat du ciel, Califat).
    \item \textbf{C. Solidarité concrète :} Zakat (aumône islamique), Dîme/Quête (Églises), Tontines paroissiales, Caritas, hôpitaux confessionnels. Gestion de la précarité là où l'État faillit.
\end{itemize}

\textbf{II. La religion, facteur de division et de conflit (Antithèse : Marx, Freud, Actualité géopolitique)}
\begin{itemize}
    \item \textbf{A. Division interne : Hétérodoxie / Hérétique / Laïque.} Exclusion des « impurs », femmes, minorités sexuelles, libre-penseurs. Guerre de succession religieuse (crises de chefferies au Cameroun). Fanatisme, rigorisme (Boko Haram au Nord-Cameroun : division mortelle au nom de l'islam).
    \item \textbf{B. Division externe : Guerres de religion / Choc des civilisations.} Histoire : Croisades, Guerres de religion (Europe), Partition Inde/Pakistan. Aujourd'hui : Conflits intercommunautaires (Centrafrique, Nord-Cameroun). « Nous vs Eux » : le sacré devient \cle{frontière identitaire} exclusive.
    \item \textbf{C. Aliénation et frein au progrès (Marx/Freud) :} « Opium du peuple » : console pour mieux faire accepter l'injustice (division de classe). « Illusion » : frein à l'autonomie rationnelle et scientifique (division sujet/objet de savoir).
\end{itemize}

\textbf{III. Dépassement : Laïcité, Dialogue et Spiritualité laïque (Synthèse)}
\begin{itemize}
    \item \textbf{A. La laïcité comme arbitre neutre (Loi 1905 France, Constitution Cameroun Art. 1) :} Séparation Églises/État. Neutralité de l'espace public. Garantie de la \cle{liberté de conscience} (croire ou ne pas croire). Transforme la division potentielle en \cle{pluralisme} géré par le droit.
    \item \textbf{B. Le dialogue interreligieux (Assise 1986, Plateforme interreligieuse Cameroun) :} Reconnaissance de l'altérité. Passage du « tolérance » (je supporte) au « respect » (je reconnais ta dignité). Valeurs communes : paix, charité, écologie (Laudato Si).
    \item \textbf{C. La « religion civile » (Rousseau / Bellah) :} Cérémonies républicaines (10 mai, 20 mai au Cameroun), hymne, drapeau, héros nationaux. Sacralise la \cle{patrie} et la \cle{loi} sans dogme théologique. Cohésion \emph{inclusive} (croyants + athées).
\end{itemize}

\textbf{Conclusion :}
\begin{itemize}
    \item \textbf{Bilan :} La religion est \emph{ambivalente} : elle est \emph{puissamment} cohésive \emph{à l'intérieur} du groupe (lien fort), mais \emph{structurellement} divisée \emph{vis-à-vis de l'extérieur} (frontière sacrée).
    \item \textbf{Ouverture :} Le défi camerounais (et mondial) n'est pas d'effacer le religieux, mais de l'inscrire dans un \cle{cadre juridique laïque} qui transforme la « guerre de dieux » en « dialogue des cultures ». La cohésion durable ne vient pas de l'uniformité religieuse, mais de la \cle{fraternité laïque}.
\end{itemize}
\end{exoresolu}

\begin{methode}[Articuler Foi, Raison et Laïcité dans le contexte camerounais]
\textbf{Objectif :} Maîtriser les trois régimes de rapport Foi/Raison et appliquer le principe de laïcité à des cas concrets camerounais.
\begin{enumerate}
    \item \textbf{Les trois positions classiques :}
    \begin{itemize}
        \item \textbf{Fidéisme / Irrationalisme (Tertullien, Kierkegaard, Pascal - cœur) :} « Credo quia absurdum ». La foi \emph{contredit} la raison ou la dépasse infiniment. \cle{Saut qualitatif}. Raison impuissante sur le salut.
        \item \textbf{Rationalisme théologique / Concordisme (Thomas d'Aquin, Anselme) :} \cle{Fides quaerens intellectum}. Foi et Raison \emph{ne peuvent se contredire} (même source : Dieu). Raison \emph{prépare} la foi (preambula fidei) et l'\emph{explicite}. Preuves de l'existence.
        \item \textbf{Séparation méthodologique / Laïcité (Kant, Comte, Jaurès, Loi 1905) :} \cle{Domaines distincts}. Raison = Savoir (science, morale autonome, droit). Foi = Croyance (salut, sens ultime, intime). « Croire n'est pas savoir ». L'État ne tranche pas la vérité théologique.
    \end{itemize}
    \item \textbf{Laïcité camerounaise (Constitution, Art. 1 \& 65) :} « La République du Cameroun est un État laïc. » $\rightarrow$ \textbf{Neutralité de l'État} (pas de religion d'État, pas de financement direct des cultes, école publique neutre) + \textbf{Liberté de conscience/culte} (pratique publique/privée, respect de l'ordre public).
    \item \textbf{Grille d'analyse d'un cas concret (Ex: Port du voile à l'école, Fêtes chômées, Prêt bancaire islamique, Gestion des chefferies) :}
    \begin{enumerate}
        \item Qualifier le fait : Acte de foi ? Pratique cultuelle ? Revendication identitaire ? Norme juridique ?
        \item Identifier le principe en tension : Liberté de conscience vs Neutralité du service public / Ordre public / Égalité femmes-hommes / Laïcité de l'instruction.
        \item Appliquer la jurisprudence / le droit positif : Décision du Conseil Constitutionnel, Circulaires ministérielles, Code du travail, Code pénal.
        \item Conclure : Solution d'équilibre (accommodement raisonnable vs interdiction proportionnée).
    \end{enumerate}
\end{enumerate}
\textbf{Reflexe Bac :} Ne pas confondre \emph{Laïcité} (séparation institutionnelle, neutralité de l'État) et \emph{Athéisme d'État} (promotion de l'incroyance, interdiction du religieux). Le Cameroun est laïc, pas athée.
\end{methode}

\begin{exoresolu}[Cas pratique : Laïcité et gestion du fait religieux au Cameroun]
\textbf{Énoncé :} « Dans un lycée public de Maroua, une élève voilée (hijab) est refusée en classe par le proviseur au nom de la laïcité et du règlement intérieur interdisant les « signes ostentatoires ». La famille saisit le tribunal administratif au nom de la liberté de conscience (Art. 1er Constitution). En tant que juriste/philosophe, analysez le conflit et proposez une solution argumentée. »

\tcblower
\textbf{Solution détaillée :}

\textbf{1. Qualification juridique et philosophique des faits :}
\begin{itemize}
    \item \textbf{Acte de l'élève :} Port du voile (hijab). Manifestation \emph{extérieure} d'une conviction religieuse intime (liberté de conscience, Art. 18 DUDH, Art. 1 Const. Cameroun). Signe \emph{religieux} (non purement culturel).
    \item \textbf{Acte du proviseur :} Interdiction basée sur règlement intérieur (« signes ostentatoires »). Invoque \textbf{laïcité de l'école publique} (neutralité du service public de l'éducation) et \textbf{ordre public scolaire} (propreté, discipline, mixité, protection contre le prosélytisme/pression de groupe).
    \item \textbf{Conflit :} \cle{Liberté de manifester sa religion} (Art. 9 CEDH / Constitution) vs \cle{Neutralité de l'institution publique / Laïcité}.
\end{itemize}

\textbf{2. Analyse des principes en tension :}
\begin{itemize}
    \item \textbf{Liberté de conscience (Absolue en for intérieur, relative en manifestation) :} L'État ne peut contraindre la croyance. La manifestation (culte, signe) peut être limitée \emph{si} : loi précise + but légitime (ordre public, droits d'autrui) + nécessité/proportionnalité.
    \item \textbf{Laïcité scolaire (Spécificité) :} L'école publique est le lieu de la \cle{formation du citoyen autonome}, à l'abri des pressions communautaires. La neutralité protège l'élève \emph{contre} l'emprise religieuse (famille, groupe) \emph{autant que} contre l'endoctrinement étatique.
    \item \textbf{Critère du « Signe ostentatoire » (Référence loi française 2004 / Jurisprudence CEDH \emph{Dahlab}, \emph{Sahin}) :} Interdiction des signes \emph{par nature} ostentatoires (voile, kippa, grande croix) qui \emph{affichent} l'appartenance religieuse de manière prosélyte ou identitaire forte, vs signes \emph{discrets} (petite croix, coran de poche, nom religieux) tolérés.
\end{itemize}

\textbf{3. Application au contexte camerounais (Droit positif) :}
\begin{itemize}
    \item \textbf{Constitution Art. 1 :} « L'État est laïc. » $\rightarrow$ Neutralité.
    \item \textbf{Loi n°98/004 du 14 avril 1998 (Orientation de l'éducation) :} Art. 4 : « L'éducation est laïque. » Art. 36 : « L'enseignement public est laïc. »
    \item \textbf{Règlements intérieurs :} Doivent être conformes à la loi. Une interdiction \emph{générale et absolue} de tout signe religieux pose problème (liberté). Une interdiction ciblée sur les \emph{signes ostentatoires/prosélytes} (voile intégral, turban, grande croix visible) est généralement validée par le Conseil Constitutionnel camerounais et la jurisprudence administrative pour préserver la neutralité du cadre d'apprentissage et l'égalité filles/garçons.
    \item \textbf{Réalité locale :} Au Nord, port du voile souvent culturel/identitaire (peul) + religieux. Risque de stigmatisation si interdiction brutale. Risque de communautarisme si tolérance totale.
\end{itemize}

\textbf{4. Solution argumentée (Accommodement raisonnable / Proportionnalité) :}
\begin{enumerate}
    \item \textbf{Rejeter l'exclusion pure et simple de l'élève :} Contraire au droit à l'éducation (Art. 26 DUDH, Const. Cameroun). L'exclusion est une sanction disproportionnée pour un vêtement.
    \item \textbf{Valider le principe d'interdiction des signes ostentatoires dans l'enceinte scolaire (salle de classe, cour) :} Le voile (hijab couvrant cheveux/cou) est un signe ostentatoire au sens de la jurisprudence. La neutralité de l'école justifie cette restriction \emph{proportionnée} à la liberté de manifester.
    \item \textbf{Proposer la médiation :} Dialogue Proviseur / Parents / Élève / Représentants cultuels (Plateforme interreligieuse).
    \item \textbf{Issu concret :} L'élève retire le voile \emph{dans l'enceinte scolaire} (salle de classe, cour) pendant les heures de cours. Elle le remet à la sortie. Respect de la laïcité de l'instruction + Respect de la liberté de conscience (elle reste musulmane, prie si besoin dans local prévu, porte voile hors école).
    \item \textbf{Sanction :} Si refus persistant après dialogue $\rightarrow$ Conseil de discipline $\rightarrow$ Mesure éducative (pas exclusion définitive immédiate).
\end{enumerate}

\textbf{Conclusion :} La laïcité camerounaise n'est pas la chasse au religieux, mais la \cle{neutralité du cadre commun} permettant le \cle{vivre-ensemble}. La solution réside dans la \textbf{distinction espace public (école = neutre) / espace privé (sortie = libre)} et le \textbf{dialogue}, non dans l'affrontement juridique stérile.
\end{exoresolu}

\begin{methode}[Réussir l'explication de texte philosophique (Méthode Bac Tech)]
\textbf{Objectif :} Structurer une explication de texte en 4 temps rigoureux (Contexte / Thèse / Structure / Commentaire critique) pour obtenir la note maximale.
\begin{enumerate}
    \item \textbf{Étape 1 : L'Introduction (5-7 lignes) — \cle{OBLIGATOIRE}}
    \begin{itemize}
        \item \textbf{Accroche / Contexte :} Auteur, Œuvre, Date, Courant philosophique.
        \item \textbf{Thèse générale du texte :} Une phrase résumant \emph{ce que l'auteur veut prouver}.
        \item \textbf{Problématique implicite :} La question à laquelle le texte répond.
        \item \textbf{Annonce du plan de l'explication :} « Nous suivrons le mouvement de la pensée de l'auteur en 3 temps : 1... 2... 3... »
    \end{itemize}
    \item \textbf{Étape 2 : L'Explication linéaire (Corps) — \cle{Suivre le fil du texte, paragraphe par paragraphe}}
    \begin{itemize}
        \item \textbf{Repérer les articulations logiques :} Connecteurs (Donc, Or, Mais, En effet, C'est pourquoi).
        \item \textbf{Pour chaque étape :} \textbf{Citer} (guillemets, ligne) $\rightarrow$ \textbf{Expliquer} (vocabulaire, concepts, références implicites) $\rightarrow$ \textbf{Analyser la logique} (argument, exemple, objection levée).
        \item \textbf{Ne pas faire de hors-texte} (pas de dissertation déguisée). Restez \emph{dans} le texte.
    \end{itemize}
    \item \textbf{Étape 3 : La Synthèse / Portée (1 paragraphe) — \cle{Recul critique}}
    \begin{itemize}
        \item Résumer la démonstration globale.
        \item Situer dans l'œuvre / le courant / l'histoire des idées.
        \item \textbf{Évaluation critique (Mesurée) :} Force de l'argumentation ? Limites ? Objections possibles (auteurs contraires) ? Actualité ?
    \end{itemize}
    \item \textbf{Étape 4 : La Conclusion (3-4 lignes) — \cle{Bilan + Ouverture}}
    \begin{itemize}
        \item Réponse directe à la problématique.
        \item Ouverture sur un auteur opposé, un débat actuel, une autre œuvre.
    \end{itemize}
\end{enumerate}
\textbf{Formules de liaison :} « L'auteur argue que... », « Il s'appuie sur l'exemple de... », « Cette distinction lui permet de... », « On peut toutefois objecter avec [Auteur] que... ».
\end{methode}

\begin{exoresolu}[Explication de texte : Kant, « La foi rationnelle » (Extrait simulé Bac)]
\textbf{Énoncé :} « La foi rationnelle ne peut être qu'une foi \emph{morale}, c'est-à-dire une foi qui a pour principe unique la loi morale en nous. Ce n'est pas une connaissance théorique de l'existence de Dieu, mais une \emph{certitude subjective} suffisante pour l'usage pratique de la raison. Je \emph{dois} faire mon devoir ; donc je \emph{puis} espérer le souverain bien ; donc l'existence de Dieu et l'immortalité de l'âme sont des \emph{postulats} nécessaires de la raison pure pratique. » (Kant, \emph{Critique de la raison pratique}, adapté).

\textbf{Consigne :} Expliquez ce texte en suivant la méthode : Introduction, Explication structurée (3 parties), Portée critique, Conclusion.

\tcblower
\textbf{Solution détaillée :}

\textbf{Introduction}
\begin{itemize}
    \item \textbf{Contexte :} Emmanuel Kant (1724-1804), philosophe allemand des Lumières, fondateur du criticisme. Texte issu de la \emph{Critique de la raison pratique} (1788), second grand ouvrage après la \emph{Critique de la raison pure} (1781).
    \item \textbf{Thèse :} Kant fonde la foi religieuse non sur la connaissance théorique (impossible), mais sur la \textbf{nécessité morale} : l'existence de Dieu et l'immortalité de l'âme sont des \textbf{postulats pratiques} indispensables pour que le devoir moral ait un sens (réalisation du souverain bien).
    \item \textbf{Problématique :} Comment la raison, incapable de \emph{prouver} Dieu théoriquement, peut-elle \emph{légitimer} la foi en Dieu pratiquement ?
    \item \textbf{Plan :} 1. La foi rationnelle = foi morale (principe interne). 2. Distinction cruciale : Certitude subjective pratique vs Connaissance théorique. 3. La déduction des postulats (Dieu, Immortalité) à partir de l'impératif catégorique.
\end{itemize}

\textbf{Développement (Explication linéaire)}

\textbf{Partie 1 : « La foi rationnelle ne peut être qu'une foi morale... principe unique la loi morale en nous » (L'origine subjective de la foi)}
\begin{itemize}
    \item \textbf{Analyse :} Kant qualifie la foi de « rationnelle » : elle n'est pas irrationnelle (fanatisme) ni révélée (hétérodoxie), elle émane de la \cle{raison pratique} (volonté rationnelle).
    \item \textbf{Concept clé :} « Principe unique = loi morale en nous ». La source n'est pas l'observation du monde (physique-théologie) ni l'idée innée (onto-théologie), mais le \cle{Fait de la raison} : la conscience immédiate du \cle{Devoir} (Impératif catégorique : « Agis uniquement d'après la maxime... »).
    \item \textbf{Portée :} La foi est \emph{autonome}. Je ne crois pas \emph{parce que} Dieu existe (connaissance), je « crois » (je postule) \emph{parce que} je dois être moral. La moralité est la \cle{condition de possibilité} de la foi rationnelle.
\end{itemize}

\textbf{Partie 2 : « Ce n'est pas une connaissance théorique... mais une certitude subjective suffisante pour l'usage pratique » (La distinction épistémologique fondamentale)}
\begin{itemize}
    \item \textbf{Analyse :} Référence directe à la \emph{Critique de la raison pure}. La raison \emph{théorique} (spéculative) ne connaît que les phénomènes (espace/temps/catégories). Dieu = Noumène $\rightarrow$ \textbf{Inconnaissable}. Les preuves (ontologique, cosmologique, physico-théo) échouent.
    \item \textbf{Concept clé :} \cle{Certitude subjective (Fürwahrhalten)} vs \cle{Connaissance objective (Wissen)}. \emph{Wissen} = Savoir objetivable, démontrable, universel (Science). \emph{Fürwahrhalten} = Tenir pour vrai, engagement personnel, subjectif, mais \emph{rationnellement nécessaire} pour l'agent moral.
    \item \textbf{Nuance :} « Suffisante pour l'usage pratique ». Je n'ai pas besoin de \emph{savoir} que Dieu existe pour \emph{agir moralement}, j'ai besoin d'y \emph{compter} (postuler) pour que mon action ait un horizon de sens (le Souverain Bien).
\end{itemize}

\textbf{Partie 3 : « Je dois... donc je puis... donc l'existence de Dieu... sont des postulats nécessaires » (La déduction pratique / Le syllogisme moral)}
\begin{itemize}
    \item \textbf{Structure déductive (Modus ponens pratique) :}
    \begin{enumerate}
        \item \textbf{Majeure morale :} « Je \emph{dois} faire mon devoir. » (Fait de la raison, inconditionné, catégorique).
        \item \textbf{Mineure téléologique :} Le devoir vise le \cle{Souverain Bien} (synthèse de la Vertu \emph{et} du Bonheur proportionné à la vertu).
        \item \textbf{Condition de possibilité :} « Je \emph{puis} espérer le souverain bien. » (Ought implies Can : « Tu dois, donc tu peux »).
        \item \textbf{Postulats requis :} Or, dans le monde sensible, vertu $\neq$ bonheur (souvent inversé). La finitude humaine empêche l'achèvement de la vertu. $\rightarrow$ Il \emph{faut} postuler :
        \begin{itemize}
            \item \textbf{Immortalite de l'âme} : Pour le progrès infini vers la sainteté (vertu parfaite).
            \item \textbf{Existence de Dieu} : Comme \cle{Auteur de la nature} capable d'accorder le bonheur à la vertu (cause intelligente du monde moral).
        \end{itemize}
    \end{enumerate}
    \item \textbf{Statut des postulats :} Non pas \emph{objets de connaissance}, mais \emph{conditions subjectives de possibilité} de l'action morale. « Je \emph{dois} $\rightarrow$ Je \emph{postule} ».
\end{itemize}

\textbf{Portée critique / Synthèse}
\begin{itemize}
    \item \textbf{Force :} Kant \textbf{sauve la religion} du scepticisme théorique (Hume) et du dogmatisme (Métaphysique wolffienne). Il la fonde sur l'\textbf{autonomie} morale (dignité humaine). Il sépare radicalement \emph{Savoir} (Science) et \emph{Croire} (Morale) $\rightarrow$ Fondement philosophique de la \textbf{laïcité}.
    \item \textbf{Limites / Objections :}
    \begin{itemize}
        \item \textbf{Hegel / Idéalisme :} La foi réduite à « moralisme » perd sa spécificité théologique (Grâce, Révélation, Amour/Trinité). Dieu devient un « commis de la morale » (garant du salaire).
        \item \textbf{Nietzsche / Athéisme :} « Dieu est mort ». Le postulat est une illusion, une projection de la volonté de puissance faible (morale d'esclaves). L'homme doit créer ses propres valeurs.
        \item \textbf{Freud :} Le « Surmoi » (intériorisation sociale) remplace l'impératif catégorique. Dieu = Père exalté. L'illusion persiste.
        \item \textbf{Problème du « Souverain Bien » :} Pourquoi le bonheur \emph{doit}-il couronner la vertu ? N'est-ce pas une attente infantile de justice cosmique ?
    \end{itemize}
\end{itemize}

\textbf{Conclusion}
\begin{itemize}
    \item \textbf{Bilan :} Kant opère un \textbf{tournant copernicien} en théologie : Dieu n'est plus l'objet d'une science (théorique), mais le \emph{postulat} d'une volonté morale (pratique). La foi devient \emph{raisonnable} sans être \emph{savante}.
    \item \textbf{Ouverture :} Cette distinction « Savoir / Croire » structure la modernité : la \textbf{laïcité de l'État} (Kant, \emph{Le Conflit des facultés}) repose sur ce partage : l'État gère le \emph{savoir} (droit, science, éducation), la société civile gère le \emph{croire} (églises, conscience). Au Cameroun, cette philosophie sous-tend l'article 1er de la Constitution.
\end{itemize}
\end{exoresolu}

\begin{methode}[Réussir la dissertation philosophique (Méthode Bac Tech - 4h)]
\textbf{Objectif :} Produire une copie structurée, problématisée, argumentée, illustrée (exemples camerounais), en 4 heures.
\begin{enumerate}
    \item \textbf{Analyse du sujet (30 min) :} \textbf{Sur brouillon.}
    \begin{itemize}
        \item Définir \emph{tous} les termes (dictionnaire + sens philosophique).
        \item Identifier la \textbf{tension} (opposition, paradoxe, « ou bien / ou bien »).
        \item Formuler la \textbf{problématique} (Question unique, ouverte, qui engage la réflexion).
        \item Trouver le \textbf{plan dialectique} (Thèse / Antithèse / Synthèse) ou \textbf{analytique} (Causes / Conséquences / Sens) selon le sujet.
    \end{itemize}
    \item \textbf{Rédaction de l'Introduction (20 min) :} \textbf{Au net.}
    \begin{itemize}
        \item \textbf{Amorce} (1 phrase : actualité, citation, fait de société camerounais).
        \item \textbf{Définitions} des concepts clés.
        \item \textbf{Problématique} (en question directe).
        \item \textbf{Annonce du plan} (3 parties, verbes d'action : Analyser, Montrer, Dépasser).
    \end{itemize}
    \item \textbf{Développement (2h30) :} \textbf{Au net.}
    \begin{itemize}
        \item \textbf{Structure par partie :} 1 Idée directrice (titre de paragraphe) = 1 Grand paragraphe.
        \item \textbf{Architecture du paragraphe :} \textbf{Thèse} (phrase affirmative) $\rightarrow$ \textbf{Argument} (raisonnement, concept, auteur) $\rightarrow$ \textbf{Exemple/Illustration} (concret, camerounais si possible) $\rightarrow$ \textbf{Transition} (vers l'argument suivant ou la partie suivante).
        \item \textbf{Auteurs obligatoires :} Citer \emph{noms + concepts} (Durkheim / conscience collective, Marx / opium, Kant / postulat, Pascal / pari, Nietzsche / mort de Dieu, etc.).
    \end{itemize}
    \item \textbf{Conclusion (15 min) :} \textbf{Au net.}
    \begin{itemize}
        \item \textbf{Réponse synthétique} à la problématique.
        \item \textbf{Nuance finale} (Ce qui reste en suspens).
        \item \textbf{Ouverture} (Auteur, autre discipline, actualité, question nouvelle).
    \end{itemize}
    \item \textbf{Relecture (15 min) :} Orthographe, syntaxe, cohérence, respect du sujet, noms propres exacts.
\end{enumerate}
\textbf{Formule magique :} \cle{1 Idée = 1 Paragraphe = 1 Auteur/Concept + 1 Exemple concret + 1 Lien logique}.
\end{methode}

\begin{exoresolu}[Dissertation type : « Le dialogue interreligieux est-il la condition de la paix sociale ? »]
\textbf{Énoncé :} Sujet complet. Traiter en 4h. (Ici : Plan détaillé + Introduction/Conclusion rédigées + Arguments clés).

\tcblower
\textbf{Solution détaillée : Plan de travail et rédaction des parties clés}

\textbf{1. Analyse \& Problématique (Brouillon) :}
\begin{itemize}
    \item \textbf{Termes :} Dialogue (échange, écoute, non conversion), Interreligieux (entre religions différentes), Condition (moyen nécessaire ? suffisant ?), Paix sociale (absence violence, justice, vivre-ensemble, cohésion).
    \item \textbf{Tension :} Le dialogue apaise (thèse) mais les religions portent des vérités \emph{absolues} et \emph{exclusives} qui rendent le dialogue difficile, voire impossible sans relativisme (antithèse). La paix ne vient-elle pas plutôt du \emph{droit/laïcité} (cadre neutre) que du dialogue théologique ?
    \item \textbf{Problématique :} Le dialogue interreligieux est-il le \emph{moteur premier} de la paix, ou seulement un \emph{adjuvant} nécessaire mais insuffisant sans un cadre politique laïque et juste ?
    \item \textbf{Plan dialectique :}
    \begin{enumerate}
        \item \textbf{Thèse :} Le dialogue interreligieux, par la reconnaissance mutuelle, désarme la violence identitaire et construit un « nous » commun (Esprit d'Assise, Plateforme Cameroun).
        \item \textbf{Antithèse :} Le dialogue a des limites structurelles (exclusivisme dogmatique, instrumentalisation politique) ; la paix sociale relève d'abord du \textbf{Droit et de la Laïcité} (cadre neutre, justice distributive).
        \item \textbf{Synthèse :} Le dialogue est \emph{nécessaire} pour l'éthique du vivre-ensemble, mais \emph{insuffisant} sans \textbf{laïcité institutionnelle} et \textbf{justice sociale}. Articulation : Dialogue des cœurs + Contrat citoyen.
    \end{enumerate}
\end{itemize}

\textbf{2. Introduction (Rédaction au net) :}
« "Il n'y a pas de paix entre les nations sans paix entre les religions. Il n'y a pas de paix entre les religions sans dialogue entre les religions." Cette parole de Hans Küng, théologien suisse, résonne avec une acuité particulière au Cameroun, pays de "l'Afrique en miniature" où cohabitent chrétiens (catholiques, protestants, réveil), musulmans (sunnites, chiites minoritaires) et adeptes des spiritualités traditionnelles (ancêtres, forces de la nature), dans un cadre constitutionnel laïc (Art. 1). Si le \textbf{dialogue interreligieux} — échange sincère entre croyants de traditions différentes sans renoncer à sa foi — apparaît comme une \textbf{condition morale} évidente de la \textbf{paix sociale} — ordre juste et non-violent —, il se heurte à l'absolu des convictions et à la réalité des inégalités structurelles. \textbf{Le dialogue interreligieux est-il la condition \emph{suffisante} de la paix, ou n'en est-il qu'une \emph{composante indispensable} articulée à la laïcité politique et à la justice économique ?} Nous montrerons d'abord que le dialogue désarme la haine religieuse (I), puis qu'il reste impuissant sans le cadre juridico-politique de la laïcité (II), pour conclure que la paix durable naît de l'\textbf{articulation} entre fraternité spirituelle et contrat citoyen (III). »

\textbf{3. Développement (Arguments clés pour rédaction) :}

\textbf{I. Le dialogue interreligieux : désamorçage de la violence et construction du "Nous" (Thèse)}
\begin{itemize}
    \item \textbf{Arg 1 : Déconstruction de l'ennemi.} L'ignorance de l'autre alimente le préjugé (\"Islam = terrorisme\", \"Christianisme = croisade\"). Le dialogue (théologique de la vie, de l'action) humanise l'adversaire. \emph{Ex :} Rencontres « Vivre ensemble » au Centre Protestant de Yaoundé, Journées Portes Ouvertes des mosquées/églises.
    \item \textbf{Arg 2 : Valeurs communes (Éthique universelle).} Au-delà des dogmes, convergence sur : dignité humaine, charité (Zakat/Dîme), famille, écologie (Laudato Si / Khilafa), paix. \emph{Auteur :} Hans Küng (\emph{Projet d'éthique mondiale}), Pape François (\emph{Fratelli Tutti}), Cheikh Bentounes (Soufisme).
    \item \textbf{Arg 3 : Mécanisme de prévention/résolution (Cas camerounais).} \textbf{Plateforme des Confessions Religieuses du Cameroun (PCRC)}. Rôle clé en 1990 (transition), 2018 (crise anglophone : appel au calme), lutte contre Boko Haram (fatwas contre l'extrémisme). \emph{Concept :} « Capital social religieux » (Putnam) $\rightarrow$ Confiance intergroupes.
\end{itemize}

\textbf{II. Les limites du dialogue : L'exclusivisme dogmatique et la primauté du politique (Antithèse)}
\begin{itemize}
    \item \textbf{Arg 1 : Vérité absolue vs Relativisme.} « Hors de l'Église point de salut » / « Coran incréé, dernier sceau ». Dialoguer = risquer de trahir sa vérité ou tomber en syncrétisme. \emph{Auteur :} Karl Barth (Non au dialogue théologique), Carl Schmitt (Le politique = distinction Ami/Ennemi, le théologique structure le politique).
    \item \textbf{Arg 2 : Instrumentalisation politique.} Les élites utilisent la religion pour mobiliser (clientélisme ethnico-religieux). Le dialogue des chefs religieux ne stoppe pas la manipulation des foules. \emph{Ex :} Crise NOSO (Nord-Ouest/Sud-Ouest) : discours de haine relayés par certains leaders locaux malgré appels PCRC.
    \item \textbf{Arg 3 : La paix sociale = Justice + Droit (Laïcité).} Marx : « La critique de la religion est la condition de toute critique. » La paix ne vient pas de l'entente des prêtres mais de la \textbf{redistribution des richesses}, de l'\textbf{État de droit}, de l'\textbf{éducation laïque}. \emph{Ex :} Conflits fonciers (Nord) : dialogue imams/pasteurs inutile sans cadastre juste et administration impartiale. La \cle{laïcité} (neutralité État) protège les minorités (athées, minorités religieuses) que le dialogue interreligieux \emph{entre majorités} peut marginaliser.
\end{itemize}

\textbf{III. Articulation : Dialogue des cœurs + Contrat citoyen = Paix durable (Synthèse)}
\begin{itemize}
    \item \textbf{Arg 1 : Deux niveaux distincts mais complémentaires.}
    \begin{itemize}
        \item Niveau \textbf{Éthique/Subjectif} : Dialogue interreligieux $\rightarrow$ \cle{Reconnaissance}, \cle{Pardon}, \cle{Fraternité}. (Ricoeur : « Le passage de la peur à la confiance »).
        \item Niveau \textbf{Juridique/Objectif} : Laïcité / Constitution / Lois $\rightarrow$ \cle{Égalité des droits}, \cle{Neutralité}, \cle{Sanction des violences}.
    \end{itemize}
    \item \textbf{Arg 2 : Modèle camerounais visé : « Laïcité de collaboration » (pas séparation stricte à la française).} État subventionne écoles confessionnelles (service public), fête nationales multireligieuses (20 mai), mais garantit liberté de conscience. Le dialogue informe la loi (ex: Code de la famille, bioéthique), la loi protège le dialogue.
    \item \textbf{Arg 3 : Éducation à la citoyenneté interculturelle.} École publique laïque + Enseignement moral civique (EMC) incluant « Éducation à la diversité religieuse ». Former le \cle{citoyen croyant ou non} capable de \cle{raisonnement critique} et d'\cle{empathie}.
\end{itemize}

\textbf{4. Conclusion (Rédaction au net) :}
« Au terme de cette analyse, il apparaît que le dialogue interreligieux est une \textbf{condition \emph{nécessaire} mais non \emph{suffisante}} de la paix sociale. Nécessaire, car il touche au cœur de l'identité profonde des peuples camerounais : sans réconciliation des mémoires et des symboles (le dialogue), la paix n'est qu'une trêve armée ou une cohabitation froide. Insuffisant, car la paix sociale est une \textbf{construction politique et économique} qui exige un État de droit impartial (laïcité), une justice redistributive et une école émancipatrice. Le défi camerounais n'est pas de choisir entre la foi et la loi, mais d'inventer une \textbf{« laïcité hospitalière »} (Ricoeur) : un espace public où le croyant apporte sa motivation éthique (issue du dialogue) pour servir le bien commun, et où le non-croyant trouve sa place pleine et entière. La paix, au Cameroun comme ailleurs, ne se décrète pas au ciel, elle se construit sur la terre, par des lois justes et des cœurs désarmés. »

\textbf{Ouverture :} La question du « retour du religieux » dans l'espace public mondial (Habermas : « société post-séculière ») : comment penser la contribution des religions à la raison publique sans renoncer à la neutralité de l'État ?
\end{exoresolu}

\begin{attention}[Les 5 erreurs qui coûtent des points au Bac (Philosophie - Dieu/Religion)]
\begin{enumerate}
    \item \textbf{Confondre « Preuve » et « Argument » / « Foi » et « Savoir ».} \rightarrow \textit{Erreur :} « La preuve ontologique \emph{prouve} scientifiquement Dieu. » \textit{Correct :} « Elle \emph{tente} de démontrer logiquement la nécessité de l'existence à partir du concept, mais Kant montre qu'elle ne donne pas de \emph{connaissance théorique}. » \textbf{Distinction épistémologique cruciale.}
    \item \textbf{Réduire la religion à une seule cause (Réductionnisme).} \rightarrow \textit{Erreur :} « La religion, c'est la peur de la mort (point). » ou « C'est seulement l'opium du peuple. » \textit{Correct :} « Elle a une dimension \emph{psychologique} (angoisse/merveilleux), \emph{sociologique} (cohésion/projection - Durkheim/Feuerbach), \emph{anthropologique} (sacré/profane) et \emph{existentielle} (sens). » \textbf{Le programme exige la pluralité des approches.}
    \item \textbf{Confondre Laïcité et Athéisme / Laïcité et Anticléricalisme.} \rightarrow \textit{Erreur :} « L'État laïc doit interdire le voile pour imposer l'athéisme. » \textit{Correct :} « La laïcité (Art. 1 Const. Cameroun) = \textbf{Séparation} (État neutre, pas de religion d'État) + \textbf{Liberté} (culte libre, conscience). Elle protège le croyant \emph{et} le non-croyant. » \textbf{Neutralité $\neq$ Hostilité.}
    \item \textbf{Dissertation : « Catalogue d'auteurs » sans problématisation ni articulation.} \rightarrow \textit{Erreur :} « Durkheim dit ça. Marx dit ça. Freud dit ça. » (Liste de courses). \textit{Correct :} « \textbf{Thèse :} Durkheim montre la cohésion. \textbf{Antithèse :} Marx répond que cette cohésion masque l'aliénation. \textbf{Synthèse :} Ricoeur/Habermas articulent reconnaissance symbolique et justice structurelle. » \textbf{Dialogue des idées, pas défilé de noms.}
    \item \textbf{Explication de texte : Hors-texte / Paraphrase / Dissertation déguisée.} \rightarrow \textit{Erreur :} Ne pas citer le texte, expliquer le mot « Dieu » pendant 10 lignes avec son cours général, donner son avis personnel (« Je pense que... »). \textit{Correct :} \textbf{Citer (l. 3-4) $\rightarrow$ Expliquer le sens précis dans la phrase $\rightarrow$ Montrer le rôle logique dans l'argumentation de l'auteur.} \textbf{Rester DANS le texte.}
\end{enumerate}
\textbf{Dernier conseil :} \cle{Relisez-vous.} Une copie propre, structurée (paragraphes visibles, intro/conclu claires), sans fautes grossières sur les noms propres (Kant, Durkheim, Marx, Pascal, Nietzsche, Coran, Bible, Constitution), vaut 2 à 4 points de plus qu'une copie brouillonne mais savante.
\end{attention}

\newpage

\coursec{Exercices}

\courssub{Je vérifie mes connaissances}

\begin{exercice}[QCM : notions fondamentales]
Pour chaque question, une seule réponse est correcte. Entoure la lettre correspondante.

\textbf{1.} La preuve ontologique de l'existence de Dieu repose sur :
\begin{itemize}
\item[a)] l'observation de l'ordre du monde ;
\item[b)] l'analyse du concept même de Dieu comme être parfait ;
\item[c)] le témoignage des prophètes ;
\item[d)] les miracles rapportés par les textes sacrés.
\end{itemize}

\textbf{2.} Pour Freud, dans \emph{L'Avenir d'une illusion}, la religion est :
\begin{itemize}
\item[a)] une vérité révélée ;
\item[b)] une illusion répondant au besoin infantile de protection ;
\item[c)] une erreur scientifique sans fonction psychologique ;
\item[d)] une invention des prêtres pour dominer le peuple.
\end{itemize}

\textbf{3.} Selon Durkheim, le critère de distinction entre religion et magie est :
\begin{itemize}
\item[a)] la croyance en un dieu unique ;
\item[b)] l'existence d'une Église, c'est-à-dire d'une communauté morale unie ;
\item[c)] la pratique de rites secrets ;
\item[d)] la croyance aux esprits.
\end{itemize}

\textbf{4.} Le mot \emph{Entzauberung} chez Max Weber signifie :
\begin{itemize}
\item[a)] la conversion des peuples au christianisme ;
\item[b)] le désenchantement du monde par la rationalisation moderne ;
\item[c)] l'interdiction de la magie par l'Église ;
\item[d)] le retour du sacré dans les sociétés contemporaines.
\end{itemize}
\end{exercice}

\begin{exercice}[Vrai ou faux justifié]
Pour chacune des affirmations suivantes, indique si elle est vraie ou fausse, puis justifie ta réponse en deux ou trois phrases.

\textbf{1.} « La preuve cosmologique conclut à l'existence de Dieu à partir de l'existence du monde et du principe de causalité. »

\textbf{2.} « Selon Kant, l'existence est un prédicat qui ajoute une propriété au concept d'une chose. »

\textbf{3.} « Le préambule de la Constitution camerounaise de 1996 proclame la laïcité de l'État et garantit la liberté de conscience et de culte. »

\textbf{4.} « Pour Feuerbach, Dieu est la projection de l'essence humaine idéalisée : l'homme a créé Dieu à son image. »

\textbf{5.} « Le déisme affirme que Dieu intervient constamment dans l'histoire humaine par des miracles. »
\end{exercice}

\begin{exercice}[Définitions]
Définis en trois à cinq lignes chacune des notions suivantes, en citant au moins un auteur associé :
\begin{enumerate}
\item Le \cle{théisme} ;
\item L'\cle{agnosticisme} ;
\item Le \cle{panthéisme} (Spinoza) ;
\item La \cle{laïcité} ;
\item Le \cle{sentiment religieux} (Schleiermacher).
\end{enumerate}
\end{exercice}

\begin{exercice}[Analyse de syllogisme]
On considère la preuve morale de l'existence de Dieu chez Kant, dont on donne la structure simplifiée :
\begin{itemize}
\item Prémisse 1 : La loi morale commande de rechercher le souverain bien (accord de la vertu et du bonheur).
\item Prémisse 2 : Or le souverain bien est impossible dans la nature si aucune cause suprême n'harmonise vertu et bonheur.
\item Conclusion : Il faut donc postuler l'existence de Dieu comme condition de possibilité du souverain bien.
\end{itemize}
\begin{enumerate}
\item Identifie la forme logique de ce raisonnement (déduction, induction, analogie ?).
\item Explique pourquoi Kant dit que Dieu est ici un \emph{postulat} de la raison pratique et non un objet de connaissance.
\item Formule une objection possible à la prémisse 2.
\end{enumerate}
\end{exercice}

\courssub{Je m'entraîne}

\begin{exercice}[Preuves et critiques : le tableau récapitulatif]
Recopie et complète le tableau suivant en t'aidant du cours.

\begin{center}
\begin{tabularx}{\linewidth}{|l|X|X|X|}
\hline
\rowcolor{popblueL} \textbf{Preuve} & \textbf{Auteur de référence} & \textbf{Principe du raisonnement} & \textbf{Critique principale} \\
\hline
Ontologique & & & \\
\hline
Cosmologique & & & \\
\hline
Physico-théologique (par le dessein) & & & \\
\hline
Morale & & & \\
\hline
\end{tabularx}
\end{center}

\begin{enumerate}
\item Pour chaque preuve, cite l'auteur classique qui la défend (Anselme de Cantorbéry, Thomas d'Aquin, Kant...).
\item Résume le raisonnement en une phrase.
\item Indique la critique correspondante (Kant, Hume, la science moderne...).
\end{enumerate}
\end{exercice}

\begin{exercice}[Cartographie des conceptions de Dieu]
À l'aide des notions étudiées en cours, classe les positions suivantes dans un schéma (tableau ou arbre conceptuel) : théisme, déisme, panthéisme, athéisme, agnosticisme, polythéisme.
\begin{enumerate}
\item Définis chaque position en une phrase.
\item Associe à chacune un auteur ou une tradition (Spinoza, Voltaire, les religions traditionnelles africaines, Sartre...).
\item Un élève de ta classe déclare : « Je ne sais pas si Dieu existe, et personne ne peut le savoir. » À quelle position cette déclaration correspond-elle ? Justifie.
\item Un chef traditionnel de l'Ouest camerounais affirme que les ancêtres veillent sur le village. Comment qualifier cette croyance ? Peut-on la réduire à l'une des catégories précédentes ? Discute.
\end{enumerate}
\end{exercice}

\begin{exercice}[Lecture de données : la pratique religieuse au Cameroun et en France]
On donne le tableau suivant, inspiré des enquêtes Afrobarometer et des enquêtes européennes (valeurs arrondies, en pourcentage de personnes déclarant que « la religion est très importante dans leur vie »).

\begin{center}
\begin{tabularx}{\linewidth}{|l|X|X|X|X|}
\hline
\rowcolor{popblueL} \textbf{Pays} & \textbf{2002} & \textbf{2008} & \textbf{2015} & \textbf{2022} \\
\hline
Cameroun & 88 & 90 & 91 & 89 \\
\hline
France & 25 & 22 & 20 & 18 \\
\hline
\end{tabularx}
\end{center}

\begin{enumerate}
\item Calcule le taux de variation (en points de pourcentage) entre 2002 et 2022 pour chaque pays.
\item Compare les deux évolutions : que constates-tu ?
\item Ces données confirment-elles la thèse weberienne du \emph{désenchantement du monde} ? Nuance ta réponse.
\item La théorie de Norris et Inglehart relie religiosité et sentiment de sécurité existentielle (conditions de vie, accès aux soins, stabilité). Comment cette théorie permet-elle d'expliquer l'écart observé ?
\end{enumerate}
\end{exercice}

\begin{exercice}[Cas pratique : religion et école]
Dans un lycée technique de Douala, une élève de Terminale, de confession musulmane, refuse de participer à une séance de dissection en SVT, invoquant ses convictions religieuses. L'enseignant lui attribue la note de 0/20 pour insoumission. Les parents saisissent le chef d'établissement.
\begin{enumerate}
\item Quels sont les deux principes en conflit dans cette situation ? (liberté de conscience / obligations scolaires)
\item Rappelle ce que dit le droit camerounais sur la liberté de conscience et de culte.
\item Propose une solution d'« aménagement raisonnable » qui respecte à la fois l'élève et les exigences du programme. Justifie ta proposition par un argument philosophique.
\end{enumerate}
\end{exercice}

\courssub{Je me prépare au Bac}

\begin{exercice}[Dissertation : Dieu et science]
\textbf{Sujet :} \emph{L'existence de Dieu est-elle une question scientifique ?}

\textbf{Consignes.} Rédige une dissertation complète (introduction avec problématique, développement en trois parties, conclusion). Tu mobiliseras obligatoirement :
\begin{itemize}
\item les preuves classiques de l'existence de Dieu et la critique kantienne de la théologie rationnelle ;
\item la distinction entre ordre de la foi et ordre de la démonstration (Pascal, Kant) ;
\item la notion de « Dieu des interstices » et le débat contemporain sur le réglage fin (\emph{fine-tuning}) de l'univers ;
\item au moins deux exemples concrets (débat créationnisme/évolution, attitudes des croyants face à la science au Cameroun).
\end{itemize}

\textbf{Questions préparatoires :}
\begin{enumerate}
\item Analyse le sujet : quel est le sens de « question scientifique » ? Quel présupposé du sujet faut-il interroger ?
\item Construis un plan détaillé (thèse, antithèse, dépassement) avec les arguments et références de chaque partie.
\item Rédige l'introduction complète (accroche, définition des termes, problématique, annonce du plan).
\end{enumerate}
\end{exercice}

\begin{exercice}[Dissertation : laïcité et liberté religieuse]
\textbf{Sujet :} \emph{La laïcité garantit-elle la liberté religieuse ?}

\textbf{Consignes.} Rédige une dissertation complète. Tu t'appuieras sur :
\begin{itemize}
\item la définition de la laïcité (neutralité de l'État, séparation des Églises et de l'État, liberté de conscience) ;
\item le préambule de la Constitution camerounaise de 1996 et la place des écoles confessionnelles dans le système éducatif camerounais ;
\item les tensions contemporaines : port de signes religieux à l'école, fêtes religieuses et service public ;
\item la distinction entre laïcité d'exclusion et laïcité de reconnaissance.
\end{itemize}

\textbf{Questions préparatoires :}
\begin{enumerate}
\item Distingue la laïcité comme principe juridique et l'athéisme comme position philosophique. Pourquoi cette distinction est-elle décisive pour traiter le sujet ?
\item La laïcité peut-elle, dans certains cas, \emph{limiter} la liberté religieuse ? Donne un exemple et analyse-le.
\item Rédige la conclusion de ta dissertation : elle doit répondre clairement à la problématique et proposer une ouverture.
\end{enumerate}
\end{exercice}

\begin{exercice}[Explication de texte : Freud]
\textbf{Texte.} « Les idées religieuses [...] sont des illusions, la réalisation des désirs les plus anciens, les plus forts, les plus pressants de l'humanité ; le secret de leur force est la force de ces désirs. [...] Le terrifiant sentiment d'impuissance de l'enfant suscite le besoin d'être protégé — protégé par amour —, besoin auquel le père pourvoit ; la reconnaissance que cette impuissance dure toute la vie fit que l'homme s'accrocha à un père, mais à un père cette fois beaucoup plus puissant. » (Freud, \emph{L'Avenir d'une illusion}, 1927, adapté.)

\textbf{Questions :}
\begin{enumerate}
\item Dégage la thèse du texte et montre comment Freud définit la religion comme \emph{illusion}.
\item Explique les concepts clés : \emph{illusion}, \emph{impuissance}, \emph{réalisation de désir}. En quoi l'illusion se distingue-t-elle, selon Freud, de l'erreur ?
\item Reconstitue l'argumentation : quel rôle joue la comparaison entre l'enfant et l'adulte croyant ?
\item Évalue la portée critique de cette thèse : réfute-t-elle l'existence de Dieu, ou seulement l'origine psychologique de la croyance ? La critique de Freud s'applique-t-elle aux religions traditionnelles africaines ?
\item Actualité : au Cameroun, la religion répond-elle seulement à un besoin de protection, ou remplit-elle d'autres fonctions (sociales, morales, politiques) ? Argumente avec des exemples.
\end{enumerate}
\end{exercice}

\newpage

\coursec{Corrigés des exercices}

\coursec{Leçon 12 : Dieu et la religion}
\courssub{Exercices corrigés type Bac Tech}

\begin{corrige}[Exercice 1 — QCM : notions fondamentales]

\textbf{1. Réponse b).} La preuve ontologique (Anselme de Cantorbéry, \emph{Proslogion}) raisonne à partir du seul \emph{concept} de Dieu : Dieu est « l'être dont on ne peut rien penser de plus grand » ; or un être qui n'existe pas serait moins parfait qu'un être existant ; donc Dieu existe. Elle ne s'appuie ni sur l'observation (a) ni sur la révélation (c, d).

\textbf{2. Réponse b).} Pour Freud, la religion est une \emph{illusion} (et non une simple erreur) : elle réalise par la pensée le désir infantile de protection paternelle face à l'impuissance humaine. Elle a donc bien une fonction psychologique (ce qui exclut c), et Freud ne la réduit pas à une machination de prêtres (d).

\textbf{3. Réponse b).} Dans \emph{Les Formes élémentaires de la vie religieuse}, Durkheim définit la religion comme un « système solidaire de croyances et de pratiques relatives à des choses sacrées qui unissent en une même communauté morale, appelée Église, tous ceux qui y adhèrent ». La \emph{communauté} (l'Église) distingue la religion de la magie, qui n'a pas d'Église durable.

\textbf{4. Réponse b).} \emph{Entzauberung} = « désenchantement du monde » : avec la rationalisation scientifique et bureaucratique moderne, le monde n'est plus expliqué par des forces magiques ou sacrées ; il devient calculable et démystifié.
\end{corrige}
\begin{attention}[Points de vigilance]
Ne pas confondre \emph{illusion} freudienne (croyance motivée par un désir, dont la vérité n'est pas tranchée) et \emph{erreur} (proposition fausse). Ne pas confondre désenchantement (perte du sens magique du monde) et disparition de la religion.
\end{attention}

\begin{corrige}[Exercice 2 — Vrai ou faux justifié]

\textbf{1. Vrai.} La preuve cosmologique (Thomas d'Aquin, première et deuxième voies) part d'un fait d'expérience — le monde existe, il y a du mouvement et des causes — et, par le principe de causalité (tout effet a une cause, et on ne peut remonter indéfiniment la chaîne des causes), conclut à l'existence d'une \emph{cause première incausée}, que l'on nomme Dieu.

\textbf{2. Faux.} C'est exactement l'inverse : dans la \emph{Critique de la raison pure}, Kant affirme que « l'existence n'est pas un prédicat réel ». Dire « Dieu existe » n'ajoute aucune propriété au concept de Dieu, pas plus que cent thalers réels ne contiennent une unité de plus que cent thalers possibles. C'est la critique décisive de la preuve ontologique.

\textbf{3. Vrai.} Le préambule de la Constitution du 18 janvier 1996 proclame « l'État laïc » et affirme que « la liberté de religion et le libre exercice du culte sont garantis ». La neutralité de l'État coexiste donc avec la garantie positive de la liberté de conscience.

\textbf{4. Vrai.} Dans \emph{L'Essence du christianisme} (1841), Feuerbach montre que les attributs divins (toute-puissance, bonté, sagesse infinies) sont les perfections de l'espèce humaine, détachées de l'homme et projetées dans un être imaginaire : « l'homme a créé Dieu à son image », et non l'inverse.

\textbf{5. Faux.} C'est le \emph{théisme} qui affirme un Dieu personnel intervenant dans l'histoire (providence, miracles, révélation). Le \emph{déisme} (Voltaire, Rousseau) affirme au contraire un Dieu « horloger » : il crée le monde et ses lois, puis n'intervient plus. Le déisme rejette les miracles comme contraires à la perfection des lois naturelles.
\end{corrige}
\begin{attention}[Points de vigilance]
Une justification doit citer un auteur ou un texte et expliciter le concept (ici : prédicat, cause première, projection, déisme/théisme). Une simple phrase « c'est faux car c'est le contraire » ne vaut aucun point au Bac.
\end{attention}

\begin{corrige}[Exercice 3 — Définitions]

\textbf{1. Le théisme.} Conception de Dieu comme être \emph{personnel}, créateur du monde, distinct de lui et intervenant dans l'histoire (providence, miracles, révélation). C'est la position des grandes religions monothéistes (judaïsme, christianisme, islam) ; Thomas d'Aquin en fournit la formulation philosophique classique.

\textbf{2. L'agnosticisme.} Position qui suspend le jugement sur l'existence de Dieu : elle affirme que la question dépasse les capacités de la connaissance humaine (terme forgé par T. H. Huxley, 1869). L'agnostique ne dit ni « Dieu existe » ni « Dieu n'existe pas », mais « nous ne pouvons pas le savoir ». Kant, en limitant la connaissance aux phénomènes, fournit le fondement théorique de cette position.

\textbf{3. Le panthéisme (Spinoza).} Conception selon laquelle Dieu n'est pas distinct du monde : \emph{Deus sive Natura}, « Dieu, c'est-à-dire la Nature ». Il n'y a qu'une seule substance, infinie, dont les hommes et les choses sont des modes. Dieu ne crée pas le monde par volonté : tout découle nécessairement de sa nature (Spinoza, \emph{Éthique}, 1677).

\textbf{4. La laïcité.} Principe juridique et politique de séparation des Églises \& de l'État : l'État est neutre en matière religieuse, ne reconnaît ni ne finance de religion d'État, et garantit à tous la liberté de conscience et de culte. Elle est inscrite au préambule de la Constitution camerounaise de 1996. Elle n'est ni l'athéisme ni l'hostilité aux religions.

\textbf{5. Le sentiment religieux (Schleiermacher).} Pour Schleiermacher (\emph{Discours sur la religion}, 1799), l'essence de la religion n'est ni un savoir ni une morale, mais un \emph{sentiment} : le « sentiment de dépendance absolue », intuition immédiate de l'infini dans le fini. La religion est d'abord une expérience vécue du cœur, antérieure aux dogmes.
\end{corrige}
\begin{attention}[Points de vigilance]
Bien distinguer les définitions philosophiques (théisme, déisme, panthéisme) des positions subjectives (agnosticisme, athéisme). Pour la laïcité, citer le texte constitutionnel camerounais est un atout majeur.
\end{attention}

\begin{corrige}[Exercice 4 — Analyse de syllogisme (Kant)]

\textbf{1. Forme logique.} Il s'agit d'une \emph{déduction} : d'une définition de la loi morale (prémisse 1) et d'une analyse de ses conditions de possibilité (prémisse 2), on tire nécessairement une conclusion. Plus précisément, c'est un raisonnement par \emph{conditions de possibilité} (raisonnement transcendantal) : si le souverain bien doit être possible, alors il faut postuler une cause qui le rende possible.

\textbf{2. Dieu comme postulat.} Un \emph{postulat de la raison pratique} est une affirmation que la raison doit poser pour que l'action morale ait un sens, sans que nous puissions la \emph{connaître} théoriquement. Dieu n'est pas un objet d'expérience ni une démonstration scientifique : c'est une exigence de la moralité. Kant ne prouve pas que Dieu existe ; il montre que l'agent moral \emph{doit croire} en Dieu pour espérer que le bonheur soit proportionné à la vertu. D'où la formule : « J'ai dû limiter le savoir pour faire place à la foi. »

\textbf{3. Objection à la prémisse 2.} On peut contester que le souverain bien soit impossible sans Dieu :
\begin{itemize}
\item la vertu peut être recherchée \emph{pour elle-même}, sans espoir de récompense (stoïcisme ; Kant lui-même dit que le devoir ne dépend pas du bonheur) ;
\item l'harmonie entre mérite et bonheur pourrait être réalisée progressivement par les institutions humaines (justice sociale, État de droit) plutôt que par une cause suprême ;
\item enfin, postuler Dieu parce que ce serait « souhaitable » ressemble au raisonnement par désir que Freud dénoncera : un besoin ne prouve pas l'existence de son objet.
\end{itemize}
\end{corrige}
\begin{attention}[Points de vigilance]
Bien distinguer : \emph{démontrer} (établir une vérité théorique) et \emph{postuler} (poser une croyance exigée par la pratique). La confusion entre les deux est l'erreur la plus fréquente dans les copies sur Kant.
\end{attention}

\begin{corrige}[Exercice 5 — Preuves et critiques : le tableau récapitulatif]

\begin{center}
\begin{tabularx}{\linewidth}{|l|X|X|X|}
\hline
\rowcolor{popblueL} \textbf{Preuve} & \textbf{Auteur de référence} & \textbf{Principe du raisonnement} & \textbf{Critique principale} \\
\hline
Ontologique & Anselme de Cantorbéry, Descartes & Dieu est l'être parfait ; l'existence appartient à la perfection ; donc Dieu existe par définition. & Kant : l'existence n'est pas un prédicat ; on ne peut pas passer du concept à la réalité. \\
\hline
Cosmologique & Thomas d'Aquin & Tout ce qui existe a une cause ; impossible de remonter à l'infini ; il faut une cause première incausée. & Hume : pourquoi la série des causes aurait-elle besoin d'une cause totale ? Le principe de causalité ne vaut que dans l'expérience. \\
\hline
Physico-théologique (par le dessein) & Thomas d'Aquin (5\textsuperscript{e} voie), Paley (l'horloger) & L'ordre et la finalité du monde supposent une intelligence ordonnatrice, comme la montre suppose l'horloger. & Hume, puis Darwin : l'ordre peut s'expliquer sans dessein (sélection naturelle, lois naturelles). \\
\hline
Morale & Kant & La loi morale exige le souverain bien ; il faut donc postuler Dieu comme garant de l'accord vertu/bonheur. & Ce n'est pas une preuve d'existence mais un postulat de la foi ; un besoin moral ne fonde pas une réalité. \\
\hline
\end{tabularx}
\end{center}
\end{corrige}
\begin{attention}[Points de vigilance]
Au Bac, il faut savoir que Kant récuse les trois preuves théoriques (ontologique, cosmologique, physico-théologique) comme illégitimes \emph{spéculativement}, tout en sauvegardant Dieu comme postulat pratique. Ne jamais présenter la preuve morale comme une démonstration scientifique.
\end{attention}

\begin{corrige}[Exercice 6 — Cartographie des conceptions de Dieu]

\textbf{1 et 2. Définitions et références.}

\begin{center}
\begin{tabularx}{\linewidth}{|l|X|X|}
\hline
\rowcolor{popblueL} \textbf{Position} & \textbf{Définition} & \textbf{Auteur / tradition} \\
\hline
Théisme & Dieu personnel, créateur, intervenant dans l'histoire. & Thomas d'Aquin ; christianisme, islam. \\
\hline
Déisme & Dieu créateur des lois du monde, mais n'intervenant plus ensuite. & Voltaire, Rousseau. \\
\hline
Panthéisme & Dieu identique à la Nature : une seule substance infinie. & Spinoza. \\
\hline
Polythéisme & Croyance en plusieurs dieux ou puissances divines. & Religions grecque et romaine antiques ; nombre de religions traditionnelles. \\
\hline
Athéisme & Négation de l'existence de tout dieu. & Sartre, Nietzsche, Feuerbach. \\
\hline
Agnosticisme & Suspension du jugement : la question dépasse notre connaissance. & Huxley ; fondement théorique chez Kant. \\
\hline
\end{tabularx}
\end{center}

\textbf{3.} La déclaration « Je ne sais pas si Dieu existe, et personne ne peut le savoir » relève de l'\emph{agnosticisme} : non pas un simple doute provisoire, mais la thèse que la connaissance de Dieu est \emph{en principe} impossible à l'homme (Dieu n'est pas un phénomène possible, dirait Kant).

\textbf{4.} La croyance aux ancêtres veillant sur le village relève du \emph{culte des ancêtres}, dimension centrale des religions traditionnelles africaines : les défunts restent membres de la communauté et interviennent dans la vie des vivants. On ne peut pas la réduire aux catégories précédentes :
\begin{itemize}
\item elle n'est pas un théisme classique (les ancêtres ne sont pas un Dieu unique créateur, même si beaucoup de traditions reconnaissent un Dieu suprême, par exemple \emph{Nyambé} chez les Beti) ;
\item elle n'est pas un simple polythéisme au sens antique, car les ancêtres sont des \emph{médiations}, non des dieux autonomes ;
\item elle montre la limite d'une cartographie forgée sur le monothéisme européen : les religions traditionnelles articulent Dieu suprême, esprits, ancêtres et forces vitales dans un système propre.
\end{itemize}
Conclusion attendue : les catégories conceptuelles doivent être appliquées avec prudence et sans ethnocentrisme.
\end{corrige}
\begin{attention}[Points de vigilance]
Ne pas plaquer les catégories occidentales (théisme/athéisme) sur les religions africaines sans nuance. La distinction « Dieu suprême / ancêtres médiateurs » est cruciale pour l'analyse anthropologique et philosophique.
\end{attention}

\begin{corrige}[Exercice 7 — Lecture de données : pratique religieuse au Cameroun et en France]

\textbf{1. Taux de variation (en points de pourcentage).}

Cameroun : $89 - 88 = +1$ point. France : $18 - 25 = -7$ points.

En pourcentage relatif (non exigé mais utile) : Cameroun : $\dfrac{89-88}{88} \approx +1,1\,\%$ ; France : $\dfrac{18-25}{25} = -28\,\%$.

\textbf{2. Comparaison.} Au Cameroun, la religiosité déclarée est très élevée et \emph{stable} (autour de 90 \% sur vingt ans). En France, elle est faible et en \emph{baisse continue} (de 25 \% à 18 \%). L'écart entre les deux pays se creuse : il passe de $88-25=63$ points en 2002 à $89-18=71$ points en 2022.

\textbf{3. Désenchantement du monde ?} Le cas français \emph{confirme} la thèse de Weber : la modernisation s'accompagne d'un recul de la religion dans la vie des individus. Mais le cas camerounais la \emph{nuance fortement} : la modernisation (urbanisation, scolarisation, technologies) n'y entraîne pas de sécularisation. La thèse du désenchantement n'est donc pas universelle : elle décrit une trajectoire européenne, non une loi de toutes les sociétés.

\textbf{4. Théorie de Norris et Inglehart.} Selon cette théorie, la religiosité dépend du \emph{sentiment de sécurité existentielle} : là où la vie est exposée aux risques (maladie, pauvreté, insécurité, faible accès aux soins), la religion fournit réconfort, explication et solidarité ; là où l'État providence et la médecine réduisent les risques, ce besoin diminue. L'écart Cameroun/France s'explique ainsi : la forte religiosité camerounaise correspond à une plus grande vulnérabilité existentielle, la faible religiosité française à une sécurité matérielle élevée.
\end{corrige}
\begin{attention}[Points de vigilance]
Distinguer \emph{points de pourcentage} (différence simple) et \emph{pourcentage relatif} (taux de variation). Au Bac, une lecture de données exige : calculs exacts, comparaison chiffrée, puis mise en relation avec une thèse du cours — jamais de philosophie « à côté » des données.
\end{attention}

\begin{corrige}[Exercice 8 — Cas pratique : religion et école]

\textbf{1. Les deux principes en conflit.} D'un côté, la \emph{liberté de conscience et de culte} de l'élève, droit fondamental ; de l'autre, les \emph{obligations scolaires} attachées au programme officiel et à l'égalité de traitement des candidats à l'examen. Le conflit oppose donc un droit individuel à une exigence d'intérêt général (la qualité et l'égalité de l'enseignement).

\textbf{2. Le droit camerounais.} Le préambule de la Constitution de 1996 affirme l'attachement du peuple camerounais aux libertés fondamentales, garantit « la liberté de religion et le libre exercice du culte » et proclame la laïcité de l'État : nul ne peut être inquiété pour ses convictions religieuses. Mais la liberté de culte n'est pas absolue : elle s'exerce dans le respect de l'ordre public, des lois et des droits d'autrui — et l'école publique reste tenue d'assurer le programme national.

\textbf{3. Aménagement raisonnable.} Proposition : dispenser l'élève de la manipulation elle-même, mais lui imposer la restitution des savoirs par un autre biais — compte rendu écrit sur la base d'une vidéo de dissection, schéma légendé, ou étude de préparations déjà réalisées. Justification philosophique :
\begin{itemize}
\item la liberté de conscience protège les \emph{convictions et les actes} (ne pas disséquer), non la dispense de \emph{connaître} : la connaissance scientifique d'un fait ne vaut pas adhésion morale à sa production ;
\item la laïcité n'oblige pas l'élève à renier sa foi, mais elle n'autorise pas non plus la foi à annuler le droit commun : l'aménagement respecte les deux exigences ;
\item la note de 0/20 pour « insoumission » est disproportionnée : elle sanctionne une conviction (protégée) et non un manquement aux savoirs (qui peuvent être évalués autrement).
\end{itemize}
\end{corrige}
\begin{methode}[Méthode du cas pratique]
1) Identifier les principes en tension. 2) Rappeler le cadre juridique. 3) Chercher une solution qui ne sacrifie aucun des deux principes. 4) Justifier par un argument philosophique explicite (ici : distinction entre conviction et connaissance).
\end{methode}
\begin{attention}[Points de vigilance]
La sanction (0/20) doit être analysée au prisme de la proportionnalité. L'aménagement raisonnable est la clé jurisprudentielle (Conseil constitutionnel, Cour européenne des droits de l'homme).
\end{attention}

\begin{corrige}[Exercice 9 — Dissertation : Dieu et science]

\textbf{1. Analyse du sujet.} Une « question scientifique » est une question tranchée par l'observation, l'expérience et la démonstration : elle porte sur des faits vérifiables et réfutables. Le présupposé à interroger : le sujet suppose que l'existence de Dieu \emph{pourrait} relever de ce type de traitement — or Dieu, par définition, n'est pas un objet d'expérience parmi d'autres. Tout l'enjeu est de savoir si la raison théorique peut atteindre Dieu, ou si la question relève d'un autre ordre (foi, morale, sentiment).

\textbf{2. Plan détaillé.}

\textbf{I. La raison peut prétendre démontrer Dieu (thèse).} Les grandes preuves rationnelles traitent l'existence de Dieu comme une conclusion démonstrative : preuve ontologique (Anselme, Descartes), cosmologique et physico-théologique (Thomas d'Aquin, l'horloger de Paley). Le réglage fin de l'univers (constantes physiques ajustées à la vie) est présenté aujourd'hui comme une version moderne de l'argument du dessein.

\textbf{II. Mais ces prétentions échouent devant la critique (antithèse).} Kant ruine la théologie rationnelle : l'existence n'est pas un prédicat (critique de l'ontologique), et le principe de causalité ne vaut que dans l'expérience (critique du cosmologique). Hume et Darwin dissolvent l'argument du dessein. Inversement, la science ne peut pas \emph{réfuter} Dieu : invoquer Dieu pour combler les lacunes de la science (« Dieu des interstices ») condamne la foi à reculer avec chaque découverte — le créationnisme face à l'évolution en est l'exemple. Dieu n'est ni démontrable ni réfutable : la question n'est donc pas scientifique.

\textbf{III. Dépassement : deux ordres distincts, un dialogue possible.} Pascal distingue l'ordre de la charité (foi) et l'ordre de la raison ; Kant fait de Dieu un \emph{postulat} de la raison pratique, non un objet de savoir. La science décrit le \emph{comment} du monde, la religion répond au \emph{pourquoi} du sens. Le débat sur le réglage fin peut alors être relu non comme une preuve, mais comme un lieu de dialogue : au Cameroun, nombre de croyants cultivés (médecins, ingénieurs) articulent foi et science sans les confondre.

\textbf{3. Introduction rédigée.}

« Lorsqu'un physicien calcule l'âge de l'univers et qu'un croyant affirme que Dieu l'a créé, parlent-ils de la même chose et au même titre ? La science moderne, en expliquant toujours davantage le monde sans recourir à l'hypothèse divine, semble avoir transformé une question métaphysique en question empirique. Par \emph{question scientifique}, on entend une question susceptible d'être tranchée par l'observation et la démonstration ; par \emph{existence de Dieu}, l'affirmation d'un être suprême, créateur et parfait, qui précisément échappe à toute expérience sensible. Le problème est donc le suivant : l'existence de Dieu relève-t-elle de la démonstration rationnelle, auquel cas la science pourrait la confirmer ou la réfuter, ou appartient-elle à un ordre — celui de la foi — que la science ne peut ni établir ni détruire ? Nous verrons d'abord que la raison a prétendu démontrer Dieu, puis que ces preuves n'ont pas résisté à la critique, enfin que la distinction des ordres permet de dépasser l'opposition stérile entre science et foi. »
\end{corrige}
\begin{attention}[Barème indicatif et vigilance]
Introduction (accroche, définitions, problématique, plan) : 4 pts. Trois parties équilibrées avec références précises : 12 pts. Conclusion répondant au sujet : 2 pts. Langue et cohérence : 2 pts. Vigilance : ne jamais écrire « la science a prouvé que Dieu n'existe pas » (c'est faux : la science ne statue pas sur Dieu) ; distinguer \emph{réfuter les preuves} et \emph{réfuter Dieu}.
\end{attention}

\begin{corrige}[Exercice 10 — Dissertation : laïcité et liberté religieuse]

\textbf{1. Laïcité et athéisme.} L'athéisme est une \emph{position philosophique} privée : la négation de l'existence de Dieu. La laïcité est un \emph{principe juridique} d'organisation de l'État : neutralité publique, séparation des Églises \& de l'État, garantie égale de toutes les libertés de conscience — y compris celle du croyant. Cette distinction est décisive : si la laïcité était un athéisme d'État, elle \emph{combattrait} la religion au lieu de garantir sa liberté. Or un État laïc protège aussi bien le musulman, le chrétien, l'animiste que l'athée.

\textbf{2. La laïcité peut-elle limiter la liberté religieuse ?} Oui, dans certains cas : l'interdiction de signes religieux ostensibles dans certains services publics, ou l'obligation de suivre l'intégralité du programme scolaire, restreignent l'\emph{expression} de la religion. Analyse : la liberté de conscience (croire) est absolue et illimitable ; la liberté de \emph{manifestation} (agir selon sa foi) peut être limitée au nom de l'ordre public, de l'égalité des usagers et de la neutralité du service public. La limite est légitime si elle est générale, proportionnée et non discriminatoire ; elle devient abusive si elle vise une religion particulière ou si elle empêche l'exercice du culte sans nécessité.

\textbf{3. Conclusion rédigée.}

« La laïcité garantit la liberté religieuse à condition d'être comprise comme neutralité protectrice et non comme hostilité : en séparant l'État des Églises, elle empêche toute religion de s'imposer par la loi et assure à chacun — croyant ou non — la même liberté de conscience. Le préambule de la Constitution camerounaise de 1996 l'illustre en proclamant ensemble la laïcité de l'État et la liberté de culte, tandis que la place reconnue aux écoles confessionnelles montre qu'une laïcité de \emph{reconnaissance} peut coopérer avec les religions sans se soumettre à elles. Certes, la laïcité limite parfois la manifestation publique des convictions ; mais limiter l'expression pour protéger l'égalité de tous, ce n'est pas nier la liberté, c'est la rendre possible pour chacun. Reste une question ouverte : dans une société profondément croyante comme la société camerounaise, la laïcité ne devrait-elle pas moins exclure le religieux de l'espace public qu'apprendre à le faire dialoguer avec la raison commune ? »
\end{corrige}
\begin{attention}[Barème indicatif et vigilance]
Distinction laïcité/athéisme clairement posée : 3 pts. Mobilisation du texte constitutionnel et d'exemples camerounais : 5 pts. Traitement dialectique (garantie / limite / dépassement) : 8 pts. Conclusion répondant au sujet avec ouverture : 2 pts. Langue : 2 pts. Vigilance : ne pas transformer la dissertation en plaidoyer pour ou contre la religion ; le sujet porte sur un \emph{principe politique}, non sur la vérité des croyances.
\end{attention}

\begin{corrige}[Exercice 11 — Explication de texte : Freud]

\textbf{1. Thèse du texte.} Freud soutient que les idées religieuses sont des \emph{illusions} : des croyances dont la force ne vient pas de preuves, mais de ce qu'elles réalisent en imagination les désirs les plus profonds de l'humanité — au premier rang desquels le désir d'être protégé. La religion prolonge à l'âge adulte la situation de l'enfant protégé par son père : Dieu est le « père beaucoup plus puissant » auquel l'homme s'attache parce que son impuissance dure toute la vie.

\textbf{2. Concepts clés.}
\begin{itemize}
\item \emph{Illusion} : croyance motivée par l'accomplissement d'un désir, indépendamment de sa vérification. Freud précise ailleurs que l'illusion n'est pas nécessairement fausse (une illusion peut, par hasard, se réaliser) — mais elle n'est pas tenue pour vraie en raison de preuves.
\item \emph{Impuissance} (\emph{Hilflosigkeit}) : le sentiment de détresse de l'enfant sans défense, que l'adulte retrouve face aux forces de la nature, à la maladie et à la mort.
\item \emph{Réalisation de désir} : comme le rêve, la croyance figure en images ce que le sujet souhaite — ici, un protecteur tout-puissant et bienveillant.
\end{itemize}
Distinction illusion/erreur : l'\emph{erreur} est une proposition fausse (croire que la Terre est plate) ; l'\emph{illusion} est définie par sa \emph{source} (le désir), non par sa fausseté. C'est pourquoi Freud peut critiquer la religion sans prétendre avoir démontré l'inexistence de Dieu.

\textbf{3. Structure argumentative.} Le raisonnement procède par \emph{analogie génétique} : (a) l'enfant, impuissant, est protégé par le père ; (b) l'adulte découvre que l'impuissance ne cesse jamais (nature hostile, mort) ; (c) il transpose donc la figure paternelle dans un « père beaucoup plus puissant » : Dieu. La comparaison enfant/adulte croyant fournit l'\emph{explication de l'origine} de la croyance : la religion est comprise comme la continuation, au niveau de l'humanité, d'un besoin psychologique infantile.

\textbf{4. Portée critique.} La thèse de Freud réfute l'\emph{origine psychologique} de la croyance, non l'\emph{existence} de Dieu : expliquer pourquoi l'on croit ne prouve pas que l'objet de la croyance n'existe pas (argument génétique ≠ argument de vérité). Freud lui-même reconnaît que l'illusion peut coïncider avec la réalité. Quant aux religions traditionnelles africaines, la critique ne s'y applique que partiellement : le culte des ancêtres répond certes à un besoin de protection, mais il structure aussi la parenté, la transmission des terres, la justice coutumière — fonctions sociales que le seul schéma « père protecteur » ne suffit pas à expliquer.

\textbf{5. Actualité au Cameroun.} La religion y remplit des fonctions multiples :
\begin{itemize}
\item \emph{sociale} : solidarité des paroisses, tontines, associations religieuses, soutien en cas de deuil ou de maladie ;
\item \emph{morale} : cadre de valeurs, éducation (écoles confessionnelles nombreuses et reconnues) ;
\item \emph{politique et identitaire} : les Églises \& interviennent dans le débat public (justice, paix, élections) ;
\item \emph{existentielle} : sens donné à la souffrance, espérance — fonction proche de celle que décrit Freud.
\end{itemize}
Conclusion : la fonction protectrice repérée par Freud est réelle, mais insuffisante ; la religion est aussi, comme le montraient Durkheim, un ciment social, et une instance de production de sens et de lien politique.
\end{corrige}
\begin{attention}[Barème indicatif et vigilance]
Thèse dégagée et concepts expliqués : 6 pts. Argumentation reconstituée fidèlement : 4 pts. Évaluation critique (distinction origine/vérité) : 5 pts. Mise en perspective camerounaise argumentée : 3 pts. Langue : 2 pts. Vigilance : ne pas faire dire à Freud qu'il « prouve que Dieu n'existe pas » ; citer le texte à l'appui de chaque analyse ; distinguer \emph{expliquer} une croyance et la \emph{réfuter}.
\end{attention}

\newpage

\coursec{Fiche bilan}

\coursec{Dieu et la religion}

\begin{popfigure}
\begin{center}\tcbox[colback=popink,colframe=popink,coltext=white,arc=9pt,boxsep=4pt,left=6pt,right=6pt]{\bfseries\small Dieu et la Religion}\end{center}
\vspace{-2pt}
\begin{tcbraster}[raster columns=3,raster equal height=rows,raster column skip=6pt,raster row skip=6pt,enhanced,blankest]
\begin{tcolorbox}[enhanced,colback=white,colframe=popgreen!70,boxrule=0.9pt,arc=3pt,title={1. Origine du sentiment religieux},fonttitle=\bfseries\scriptsize,coltitle=white,colbacktitle=popgreen!70,left=4pt,right=4pt,top=2pt,bottom=2pt,boxsep=1pt,toptitle=1.5pt,bottomtitle=1.5pt,halign=left]
\begin{itemize}[nosep,leftmargin=*,itemsep=1pt,font=\scriptsize]
\item Peur \emph{(Hume)}
\item Merveilleux \emph{(Otto)}
\item Projection sociale \emph{(Feuerbach)}
\item Illusion consolatrice \emph{(Freud)}
\item Opium du peuple \emph{(Marx)}
\end{itemize}
\end{tcolorbox}
\begin{tcolorbox}[enhanced,colback=white,colframe=poporange!70,boxrule=0.9pt,arc=3pt,title={2. Conceptions de Dieu},fonttitle=\bfseries\scriptsize,coltitle=white,colbacktitle=poporange!70,left=4pt,right=4pt,top=2pt,bottom=2pt,boxsep=1pt,toptitle=1.5pt,bottomtitle=1.5pt,halign=left]
\begin{itemize}[nosep,leftmargin=*,itemsep=1pt,font=\scriptsize]
\item Théisme (Personnel)
\item Déisme (Horloger)
\item Panthéisme (Nature)
\item Athéisme (Négation)
\item Agnosticisme (Inconnaissable)
\end{itemize}
\end{tcolorbox}
\begin{tcolorbox}[enhanced,colback=white,colframe=poppurple!70,boxrule=0.9pt,arc=3pt,title={3. Preuves de l'existence},fonttitle=\bfseries\scriptsize,coltitle=white,colbacktitle=poppurple!70,left=4pt,right=4pt,top=2pt,bottom=2pt,boxsep=1pt,toptitle=1.5pt,bottomtitle=1.5pt,halign=left]
\begin{itemize}[nosep,leftmargin=*,itemsep=1pt,font=\scriptsize]
\item Ontologique \emph{(Anselme, Descartes)}
\item Cosmologique \emph{(Thomas d'Aquin)}
\item Physico-théologique \emph{(Paley)}
\item Morale \emph{(Kant)}
\item Critiques \emph{(Kant, Hume)}
\end{itemize}
\end{tcolorbox}
\begin{tcolorbox}[enhanced,colback=white,colframe=poppink!70,boxrule=0.9pt,arc=3pt,title={4. Religion : Définition \& Valeur},fonttitle=\bfseries\scriptsize,coltitle=white,colbacktitle=poppink!70,left=4pt,right=4pt,top=2pt,bottom=2pt,boxsep=1pt,toptitle=1.5pt,bottomtitle=1.5pt,halign=left]
\begin{itemize}[nosep,leftmargin=*,itemsep=1pt,font=\scriptsize]
\item Lien social \emph{(Durkheim)}
\item Aliénation \emph{(Marx)}
\item Sens \& Salut
\item Dogmes, Rites, Église
\end{itemize}
\end{tcolorbox}
\begin{tcolorbox}[enhanced,colback=white,colframe=popgold!70,boxrule=0.9pt,arc=3pt,title={5. Foi, Raison, Laïcité},fonttitle=\bfseries\scriptsize,coltitle=white,colbacktitle=popgold!70,left=4pt,right=4pt,top=2pt,bottom=2pt,boxsep=1pt,toptitle=1.5pt,bottomtitle=1.5pt,halign=left]
\begin{itemize}[nosep,leftmargin=*,itemsep=1pt,font=\scriptsize]
\item Fidéisme (Foi seule)
\item Rationalisme (Raison seule)
\item Conciliation \emph{(Thomas d'Aquin)}
\item Limites du savoir \emph{(Kant)}
\item Laïcité camerounaise (Const. 1996, Loi 2004)
\end{itemize}
\end{tcolorbox}
\begin{tcolorbox}[enhanced,colback=white,colframe=popdark!70,boxrule=0.9pt,arc=3pt,title={6. Méthodologie Bac Tech},fonttitle=\bfseries\scriptsize,coltitle=white,colbacktitle=popdark!70,left=4pt,right=4pt,top=2pt,bottom=2pt,boxsep=1pt,toptitle=1.5pt,bottomtitle=1.5pt,halign=left]
\begin{itemize}[nosep,leftmargin=*,itemsep=1pt,font=\scriptsize]
\item Dissertation : Analyser, Plan, Rédiger
\item Explication : Contexte, Structure, Portée
\item Exemples concrets CM
\item Citations clés
\end{itemize}
\end{tcolorbox}
\end{tcbraster}

\legende{Carte mentale de synthèse : Leçon 12 -- Dieu et la Religion (Terminale Technique)}
\end{popfigure}

\begin{aretenir}[Synthèse des savoirs et méthodes -- Leçon 12]

\textbf{Définitions clés à connaître par cœur}
\begin{itemize}
    \item \cle{Religion} : (du latin \emph{religare}, relier) Ensemble de croyances, de rites et de dogmes unissant une communauté au sacré ; lien social \emph{(Durkheim)} ou aliénation \emph{(Marx)}.
    \item \cle{Dieu} : Être suprême, infini, éternel, créateur. Conceptions : \emph{Théisme} (personnel, providence), \emph{Déisme} (créateur absent, horloger), \emph{Panthéisme} (confondu avec la nature/\emph{Spinoza}), \emph{Athéisme} (négation), \emph{Agnosticisme} (inconnaissable/\emph{Huxley}).
    \item \cle{Foi} : Adhésion libre et totale à une vérité révélée, sans preuve rationnelle démonstrative (Fidéisme : \emph{Tertullien, Pascal, Kierkegaard}).
    \item \cle{Raison} : Faculté de connaître, démonstration, universalité. \cle{Preuve a priori} (ontologique : concept $\to$ existence) vs \cle{Preuve a posteriori} (cosmologique, physico-théologique, morale : fait $\to$ cause).
    \item \cle{Laïcité} : Séparation des institutions étatiques et religieuses, neutralité de l'État, liberté de conscience (Préambule Constitution 1996, Loi n°2004/005 organisation enseignement privé). Ne pas confondre avec athéisme d'État.
\end{itemize}

\textbf{Tableau des preuves classiques et critiques (Bac Tech : savoir les exposer et les critiquer)}
\begin{center}
\begin{tabularx}{\linewidth}{|l|X|X|}\hline
\rowcolor{popblueL}
\textbf{Preuve / Auteur} & \textbf{Principe (Thèse)} & \textbf{Critique majeure (Antithèse)} \\ \hline
\textbf{Ontologique} \newline (Anselme, Descartes) & A priori. Dieu = \emph{ens realissimum} (être parfait). L'existence fait partie de la perfection $\Rightarrow$ Dieu existe nécessairement. & \textbf{Kant} : L'existence n'est pas un prédicat (n'ajoute rien au concept). On ne peut passer du concept à la réalité (usage transcendantal illégitime). \\ \hline
\textbf{Cosmologique} \newline (Thomas d'Aquin, Voie 2) & A posteriori. Chaîne des causes efficientes. Pas de régression à l'infini $\Rightarrow$ Cause première non causée = Dieu. & \textbf{Hume / Kant} : Principe de causalité valable seulement dans l'expérience. Saut illégitime de "cause du monde" à "Dieu personnel". \\ \hline
\textbf{Physico-théologique} \newline (Paley, Newton) & A posteriori. Ordre, finalité, complexité du monde (horloge $\to$ horloger) $\Rightarrow$ Intelligence ordonnatrice. & \textbf{Darwin / Hume} : L'ordre peut venir du hasard + sélection naturelle (évolution). Analogie monde/machine faible. Mal, désordre existants. \\ \hline
\textbf{Morale} \newline (Kant) & A posteriori (pratique). Devoir moral $\to$ Souverain Bien (Vertu + Bonheur). Nécessite postulation : Dieu, Liberté, Immortalité. & Preuve \emph{pratique} (nécessaire pour l'action), non \emph{théorique} (ne prouve pas l'existence objective). "Il faut supposer Dieu", pas "Dieu existe". \\ \hline
\end{tabularx}
\end{center}

\textbf{Méthodes express pour le Bac Tech}
\begin{enumerate}
    \item \textbf{Dissertation (ex: "La religion est-elle un obstacle au progrès ?")} :
    \begin{itemize}
        \item Analyser : Définir \cle{religion} (croyance/rite/lien) et \cle{progrès} (technique, moral, social). Tension : conservatisme vs moteur éthique.
        \item Problématiser : Opposition Marx (opium/frein) vs Durkheim (cohésion) vs Weber (éthique protestante/capitalisme) vs exemples camerounais (Écoles/Hôpitaux confessionnels CBC/PCC/Catholique, fêtes nationales interreligieuses, gestion épidémies, crise anglophone NO/SO).
        \item Plan dialectique : I. Frein (dogmatisme, obscurantisme, conflits). II. Moteur (éthique, éducation, solidarité, cadre moral). III. Dépassement : Laïcité comme condition d'un progrès partagé (distinction sphères).
    \end{itemize}
    \item \textbf{Explication de texte (ex: Kant, \emph{Fondement de la métaphysique des mœurs})} :
    \begin{enumerate}
        \item \textbf{Contexte} : Auteur, œuvre, époque (Aufklärung), thèse générale (postulats de la raison pratique).
        \item \textbf{Structure} : Découper en 2-3 étapes logiques (ex: 1. Définition du Souverain Bien. 2. Impossibilité dans ce monde. 3. Postulation de Dieu/immortalité).
        \item \textbf{Analyse concepts/arguments} : Expliquer "postuler", "souverain bien", "sainteté", distinction théorique/pratique.
        \item \textbf{Portée/Critique} : Valeur (fondation éthique), Limite (pas une preuve théorique, subjectivité), Actualité (laïcité, bioéthique).
    \end{enumerate}
\end{enumerate}

\end{aretenir}

\begin{aretenir}[Checklist Bac -- Leçon 12 : Dieu et la Religion]
\begin{itemize}
    \item $\square$ Je définis \textbf{Religion}, \textbf{Dieu}, \textbf{Foi}, \textbf{Raison}, \textbf{Laïcité} sans confusion.
    \item $\square$ Je distingue les 4 origines du sentiment religieux : \textbf{Peur} (Hume), \textbf{Merveilleux} (Otto), \textbf{Projection sociale} (Feuerbach), \textbf{Illusion/Idéologie} (Freud, Marx).
    \item $\square$ Je classe les 5 conceptions de Dieu : \textbf{Théisme}, \textbf{Déisme}, \textbf{Panthéisme}, \textbf{Athéisme}, \textbf{Agnosticisme} (avec un exemple/auteur chacun).
    \item $\square$ J'expose les \textbf{4 preuves classiques} (Ontologique, Cosmologique, Physico-théologique, Morale) avec leur auteur type et leur structure (a priori / a posteriori).
    \item $\square$ J'énonce la \textbf{critique de Kant} sur la preuve ontologique (existence n'est pas un prédicat) et sur les preuves a posteriori (usage transcendantal illégitime des catégories).
    \item $\square$ Je compare \textbf{Durkheim} (religion = cohésion sociale, sacré/profane) et \textbf{Marx} (religion = opium, soupir de la créature opprimée, idéologie).
    \item $\square$ Je distingue \textbf{Fidéisme} (Foi > Raison), \textbf{Rationalisme théologique} (Raison démontre Dieu), \textbf{Conciliation} (Thomas d'Aquin : grâce parfait nature) et \textbf{Kant} (Raison théorique nie, Raison pratique postule).
    \item $\square$ Je définis la \textbf{Laïcité camerounaise} (Constitution 1996, Loi 2004 : séparation, neutralité, liberté) et je la distingue de l'athéisme d'État.
    \item $\square$ Je mobilise \textbf{2 exemples concrets camerounais} (ex: rôle des Écoles/Hôpitaux confessionnels, fêtes nationales interreligieuses, gestion épidémies, crise anglophone NO/SO).
    \item $\square$ Je sais structurer une \textbf{dissertation} (Analyse, Problématique, Plan dialectique 3 parties, Intro/Conclu types) et une \textbf{explication de texte} (Contexte, Structure, Analyse, Portée).
    \item $\square$ Je connais \textbf{3 citations "passe-partout"} : Pascal (Pari), Kant ("Supprimer le savoir pour faire place à la croyance"), Marx ("L'opium du peuple"), Voltaire ("Si Dieu n'existait pas, il faudrait l'inventer").
\end{itemize}
\end{aretenir}

\findecours

\end{document}
